%% notes-tex/database/database-expansion-bacteria/database-expansion-bacteria.tex
%% [[experiments.database.expansion-bacteria]]
%% https://github.com/Mjvolk3/torchcell/tree/main/notes-tex/database/database-expansion-bacteria/database-expansion-bacteria.tex

\documentclass[10pt,a4paper]{article}
\usepackage[draft]{tcdoc}
\usepackage{float}

\usepackage{pdflscape}

\newcolumntype{L}[1]{>{\raggedright\arraybackslash}p{#1}}

\title{\vspace{-1.5em}Fifty Bacterial Datasets for \emph{E. coli} and \emph{P. putida},
Ranked by Scale and Density}
\author{Michael Volk}
\date{\today}

\begin{document}
\maketitle
\thispagestyle{empty}

\noindent
The database holds 51 schematized and L0--L4-verified \org{Saccharomyces cerevisiae}
datasets and no bacterial dataset at all. Table~\ref{tab:bfinal} ranks 89 candidates for
the two bacterial hosts, \org{Escherichia coli} and \org{Pseudomonas putida}, by
measurements, meaning instances times phenotype dimensionality (Sec.~\ref{sec:rule}), and
then lifts the rows that measure isoprenol, the product the bacterial work targets
(Sec.~\ref{sec:iso}). The first fifty are the recommended builds and the remaining 39 are
a ranked reserve. The
build is staged in two tranches, twenty rows then thirty, with a floor of six rows per
host in the first tranche so that neither host's setup cost waits on the other
(Sec.~\ref{sec:stage}). Of the fifty, 35 are \org{E. coli} and 15 are \org{P. putida}.

One consequence of that ordering is large enough to state before the table rather than
after it (Sec.~\ref{sec:cost}). The rows that connect an engineering intervention to a
production phenotype carry a median of 137 measurements against 226,575 for the
functional-genomics screens, so the measurement order alone leaves thirty of the
thirty-two engineering rows in the reserve. The isoprenol pin returns eleven of them to
the fifty, and Table~\ref{tab:bpins} prints the measurement rank beside the final rank so
the pin can be undone row by row.

Two changes to the schema gate every row, and both were established by running the
validators rather than by reading them (Sec.~\ref{sec:schema}): the gene-identifier
namespace rejects every bacterial gene name, and the genomes tier has no bacterial
assembly set. A third thing expected to block turns out not to, since the reference genome
record already accepts either host unchanged.

For each of the fifty recommended builds, Table~\ref{tab:banalogs} answers whether the
row fits the record types already in use, naming the built yeast dataset it mirrors and
what the loader would still need, and Table~\ref{tab:bsources} gives its citation, link
and data location. Table~\ref{tab:bqueuerows} lists the unverified discovery queue behind
the list, less the fifty; Table~\ref{tab:bexcluded} records what was considered and
dropped; Sec.~\ref{sec:summary} closes with the summary statistics and the confidence
split. Sec.~\ref{sec:acquisition} comes first and is separate from all of it: it holds a
resource whose scale is claimed rather than measured, so it is ranked nowhere and counted
nowhere, and it leads because an unlocated deposit is what a backlog loses silently.

\vspace{0.5em}
\tccontents
\clearpage

%% notes-tex/database/database-expansion-bacteria/sections/acquisition.tex
%% [[experiments.database.expansion-bacteria]]
%% https://github.com/Mjvolk3/torchcell/tree/main/notes-tex/database/database-expansion-bacteria/sections/acquisition.tex
%% SOURCE: acquisition table from
%% experiments/database/scripts/build_bacteria_candidate_datasets_table.py
%% (ACQUISITION list); surrounding prose is hand-authored

\section{Acquisition targets, outside the fifty\secstatus{todo}}\label{sec:acquisition}

Table~\ref{tab:bacquisition} holds one resource that is worth having and is not a row of
the ranked list. A ranked row earns its place with two things: a confirmed accession, and
an instance count read off a released file. A row here has neither. Its scale is asserted
in a project abstract, so it cannot be ordered against rows whose numbers came from data,
and none of its figures enter any total in this document.

The entry is a genome-wide CRISPRi guide library for \org{P. putida} KT2440. Its claim is
roughly 78,932 members at 10 to 15 guides per gene with 798 non-targeting controls,
selected in glucose, in acetate, and in p-coumaric acid at two concentrations. A search on
2026-10-07 found no publication and no deposit carrying those figures. The closest
confirmed resource is the dCpf1 library of row 20, reported at about 16,500 guide targets,
which is a different library at roughly a fifth of the claimed size.

What makes it worth a section of its own is \term{modality}, meaning the physical way a
gene is perturbed. The target campaign of row 6 represses genes with CRISPRi, while the
largest verified supplement available today, the transposon fitness compendium of row 8,
disrupts them with insertions. A model trained on insertion fitness and applied to
repression crosses that boundary, and a repression library of this size would not have to.
Several guides per gene would also supply a graded knockdown rather than a single
all-or-nothing edit.

Acquisition is therefore a prerequisite and not a schedule item: until a file exists with a
\file{sha256}, the library cannot be ranked, costed, or assigned a loader. The reason it
leads the document is that an unlocated deposit is the one kind of entry that leaves a
backlog without anyone deciding to drop it.

%% GENERATED FILE -- do not hand-edit.
%% SOURCE: experiments/database/scripts/build_bacteria_candidate_datasets_table.py
\begingroup
\footnotesize
\begin{longtable}{@{}L{30mm} L{140mm}@{}}
\caption[]{Acquisition targets, which are NOT among the fifty and carry no verified
instance count. A ranked row in Table~\ref{tab:bfinal} has a confirmed accession and a
number read off a file; a row here has neither, so the two cannot be compared and this
table is deliberately kept out of every total. \emph{Claimed} reproduces the source's
own figures and is not a measurement. \emph{Verification} is what was done to check
them, and a search that found nothing is recorded as such.}
\label{tab:bacquisition}\\
\toprule
Field & Statement \\
\midrule
\endfirsthead
\toprule
Field & Statement \\
\midrule
\endhead
\bottomrule
\endfoot
\multicolumn{2}{@{}L{172mm}@{}}{\textbf{Genome-wide KT2440 CRISPRi guide library, pooled growth selection} \quad \org{P. putida}} \\
\addlinespace[3pt]
Modality & CRISPRi repression, multiple guides per gene, pooled selection read out by guide enrichment and depletion \\
Claimed & 78,932 library members at 10 to 15 guides per gene, plus 798 non-targeting controls; selections in glucose, in acetate, and in p-coumaric acid at 10 and 50 mM. About 316,000 guide by selection slots if every member were quantified in all four selections, which is a product of the design and not a count of released values. \\
Claim source & A US Department of Energy project abstract, supplied by the owner on 2026-10-07. Not a peer-reviewed dataset release, and no accession accompanied the figures. \\
Verification & UNVERIFIED. Searched 2026-10-07 and found no publication or deposit carrying these figures. The nearest confirmed genome-wide KT2440 CRISPRi resource is the dCpf1 library of row 20 (Menasalvas 2025), reported at about 16,500 guide targets, so the claimed scale is roughly five times that and is not the same library. Nothing here was measured from a file. \\
Why it matters & It is the only candidate whose perturbation modality matches the Carruthers 2025 target of row 6: repression rather than transposon disruption, several guides per gene, and a pooled selection readout. A transposon compendium transfers across a modality boundary; this would not have to. \\
Blocking action & Locate the deposit before any loader work. Until a file exists with a sha256, this row cannot be ranked, costed or scheduled, and the claimed slot count must not be added to any instance total. \\
\addlinespace[6pt]
\end{longtable}
\endgroup


%% notes-tex/database/database-expansion-bacteria/sections/rule.tex
%% [[experiments.database.expansion-bacteria]]
%% https://github.com/Mjvolk3/torchcell/tree/main/notes-tex/database/database-expansion-bacteria/sections/rule.tex
%% SOURCE: hand-authored prose; every number quoted here comes from
%% experiments/database/scripts/build_bacteria_candidate_datasets_table.py

\section{The selection rule\secstatus{todo}}\label{sec:rule}

\subsection{Rows are ordered by measurements, which is scale times density\secstatus{todo}}

The order is one quantity, applied to every row alike: \term{measurements}, meaning
instances times phenotype dimensionality, descending. \term{Instances} is the number of
genotype $\times$ environment $\times$ replicate records a row can actually yield, and
\term{dimensionality} is how many values each of those records carries. A scalar fitness
screen of $n$ strains contributes $n$ numbers; a proteome panel of $n$ samples at a
depth of 2,000 proteins contributes $2\times10^{3}n$.

Ranking on instances alone gets the order visibly wrong, and it is worth naming the
failure rather than asserting the rule. A quantitative proteome across a couple of dozen
growth conditions holds fewer instances than a small deletion screen, which would rank it
below one, and yet each of its instances is a complete proteome. The product is what
places it correctly, and it is the same normalization the supported-datasets table
already applies through its compressed-signal column and the same one the yeast candidate
list ranks inside its bands with.

Two tie-breaks make the rule total. Instances break a measurement tie, so between two
rows holding the same number of values the one spanning more distinct strain and
condition combinations wins; a wide, shallow row is worth more to a
genotype-to-phenotype model than a narrow, deep one of equal volume. Tier breaks that.

This differs from the yeast list on purpose. There, rows are banded first by what they
supply to a planned single-cell campaign, and scale orders only inside a band. Here scale
and density are the requested rule directly, so no band is applied and no row is promoted
for what it pairs with. What a row combines with is still recorded, in
Table~\ref{tab:banalogs} and in the join column, but it does not move the row.

\subsection{What this ordering costs, stated before the table\secstatus{todo}}\label{sec:cost}

Ranking on measurements has one large consequence and it runs against the reason a
bacterial host is interesting at all. The functional-genomics screens win, and the
metabolic-engineering campaigns lose, by roughly three orders of magnitude.

The thirty-two rows that connect an explicit engineering intervention to a production
phenotype, meaning the production campaigns, the combinatorial designs and the tolerance
series, have a median of 137 measurements and a median rank of 72.5 on the measurement
order. Two of them reach the fifty on that order; thirty sit in the reserve, where they
are thirty of its 39 rows. The seventeen transposon-fitness and chemical-genomics rows
have a median of 226,575 measurements, a factor of about 1,654 more, and every one of
them reaches the fifty.

Neither figure is wrong and the rule is not misapplied. A titer measured on six strains
really is six numbers, and a barcoded library really does yield millions. What the ordering
encodes is that a dataset's worth is its volume, and for a substrate meant to support
inverse strain design that is a defensible default and an incomplete one: the rows that
teach a model what an intervention does to a product are the rows this ordering ranks last.

Two things follow, and both are now done. The ranking is reported as asked, so the
measurement order is recoverable for every row and is printed beside the final order in
Table~\ref{tab:bpins}. And the correction turned out to be one field rather than a
rewrite: a per-row product axis, pinned after ranking, reorders the same records without
recollecting anything. Section~\ref{sec:iso} applies it to the one product the work
targets, where the three-order-of-magnitude penalty had put nine of ten rows in the
reserve. With the pin applied, thirteen engineering rows sit in the fifty and nineteen in
the reserve, and three of the seventeen screens are displaced below the cut. The same
mechanism extends to any other product that becomes a priority, which is the argument for
a pin over a reweighted score.

\subsection{Sequence reconstruction is the gate, not a preference\secstatus{todo}}

A phenotype is only trainable against a genotype that can be written down, so every row
names a \term{sequence basis}: the route by which the total genomic content of one strain
would be reconstructed. The bases are named per host, because the reference differs and
because a b-number and a PP\_ locus tag are different namespaces that must not be
silently merged (Sec.~\ref{sec:schema}).

\begin{itemize}[leftmargin=1.2em,itemsep=1pt,topsep=2pt]
  \item \term{K-12-KO} and \term{KT2440-KO} -- the reference genome minus one
    cataloged gene. The cheapest basis, and the one most of the built yeast set sits on
    in its own reference.
  \item \term{K-12-KO x KO} -- a constructed double mutant, both loci cataloged, which
    is the basis a genetic-interaction row needs.
  \item \term{K-12+transposon} and \term{KT2440+transposon} -- a mapped insertion site
    carrying a barcode. The insertion position is sequenced, so the genotype is realized
    rather than designed, which is what separates a barcoded transposon library from a
    guide library.
  \item \term{K-12+guide}, \term{KT2440+guide} -- the genome is unedited and the
    perturbation is a designed cassette plus a guide target. Regulatory, not
    sequence-level, and it reaches the essential genes a deletion collection cannot hold.
  \item \term{KT2440+promoter} and \term{engineered-chassis+promoter} or
    \term{+RBS} -- a designed regulatory part over a native or heterologous gene.
    Designed rather than verified per strain.
  \item \term{engineered-chassis} -- a named production strain plus heterologous
    cassettes, which is what a titer is measured on.
  \item \term{evolved-WGS} -- a resequenced evolved clone.
  \item \term{reference-only} -- wild type, where the environment carries the
    perturbation.
\end{itemize}

\notesec{The \org{E. coli} bases say K-12 rather than naming one strain, and that is an
admission} The deletion collection is in BW25113 and the knockdown libraries are mostly
MG1655. Those are different genomes: BW25113 carries defined lesions at the lactose,
arabinose and rhamnose loci among others, so a deletion row and a knockdown row do not
sit on the same reference even when both are called K-12. Labeling a collection row with
the MG1655 name would be a sequence-level error in a substrate whose claim is
sequence-level genotype fidelity, and labeling both K-12 records the lineage without
asserting an identity that was not checked. Which background each row actually used is a
per-row provenance item to settle when its loader is written, and it is not settled here.

What fails the gate is a strain whose genome is unknown even in principle. An evolved
population that was never resequenced lands there, and so does a strain described only by
a phenotype in a review table. Those are dropped in Table~\ref{tab:bexcluded} rather than
ranked low, because no amount of phenotypic value substitutes for a missing genotype.

\subsection{Tier, and what it is for here\secstatus{todo}}

Tier no longer orders the list, so it does one job: it records how much of a dataset a
loader would actually get, which is what makes a high measurement count trustworthy or
not.

\begin{itemize}[leftmargin=1.2em,itemsep=1pt,topsep=2pt]
  \item \textbf{Tier 1} clears all three of $\geq$$10^3$ genotypes, $\geq$$10^4$
    instances, and a clean sequence basis.
  \item \textbf{Tier 2} clears two of the three, or is the only dataset covering a
    phenotype the substrate needs for these hosts.
  \item \textbf{Tier 3} is a modality or reference backbone the substrate needs without
    itself carrying a large genotype axis, such as a library definition or a per-gene
    coefficient that applies to every strain alike.
  \item \textbf{Tier 4} is real but small, coarsely labeled, or expensive to extract.
\end{itemize}

The bars apply to what a row contains rather than to its reputation. A curated corpus
with thousands of literature records fails tier 1 whatever its size, because its records
re-serve primary papers and the net-new fraction is unknown until it is split by source.

\subsection{Staging, and why the first tranche carries a floor\secstatus{todo}}\label{sec:stage}

The fifty are built in two tranches, twenty then thirty. The second tranche is the next
thirty by the ordering rule and needs no further explanation. The first twenty carry two
corrections to the rule, the isoprenol pin of Sec.~\ref{sec:iso} and the per-host floor
below. Both are applied after ranking and reported rather than folded into a score.

The correction is a floor of six rows per host, not an even split, and the difference
matters. The two schema changes in Sec.~\ref{sec:schema} are per host, and so is
provisioning a reference release, so neither cost is paid until a row for that host is
actually built; a floor guarantees both hosts start. An even ten-and-ten split would go
much further than that and fight the ordering rule, promoting four rows of a few hundred
measurements over rows of a few hundred thousand.

On these rows the floor does not bind. Ranking by measurements alone already puts
fourteen \org{E. coli} rows and six \org{P. putida} rows in the first tranche, which meets
the floor without promoting anything, and the isoprenol pin then raises the
\org{P. putida} count to nine against eleven because that host carries most of the
isoprenol work. Every
promotion the floor makes is printed, so a future pass in which it does bind will say so
rather than quietly reordering. An even ten-and-ten split would have bound here, which is
the comparison that justifies the floor over a split.

%% notes-tex/database/database-expansion-bacteria/sections/isoprenol.tex
%% [[experiments.database.expansion-bacteria]]
%% https://github.com/Mjvolk3/torchcell/tree/main/notes-tex/database/database-expansion-bacteria/sections/isoprenol.tex
%% SOURCE: hand-authored prose; every number quoted here comes from
%% experiments/database/scripts/build_bacteria_candidate_datasets_table.py and
%% experiments/database/scripts/build_bacteria_discovery_queue.py

\section{Isoprenol, and the pin that the ordering rule needs\secstatus{todo}}\label{sec:iso}

\subsection{The product named as the priority is the one the rule ranks last\secstatus{todo}}

\term{Isoprenol} is 3-methyl-3-buten-1-ol, an aviation fuel precursor, and it is the
product the bacterial strain-design work targets. It appears in the literature under two
names: the earlier papers write \term{isopentenol} and the current ones write isoprenol,
and they are the same molecule, so a row counts as measuring it under either word.

Section~\ref{sec:cost} stated the cost of ordering rows by measurements in general terms.
Isoprenol is where that cost becomes a concrete defect. A titer campaign measures tens to
a few thousand numbers, because a strain yields one titer per condition per replicate and
strains are built by hand. A barcoded fitness screen measures $10^{6}$ to $10^{8}$. No
isoprenol campaign can outrank any screen, whatever its relevance.

The effect is not marginal. Of the ten rows that measure isoprenol, the measurement rule
alone places \emph{nine} outside the fifty recommended builds, spread from rank 55 to rank
89 of 89. One survives, at rank 6, and it survives for a reason unrelated to isoprenol: it
carries a 1,500-protein proteome panel alongside its titers, and the panel depth is what
multiplies its 1,416 instances into 2.1 million measurements. Strip the proteomics and it
would join the other nine. So under the unaided rule the highest-priority product in the
portfolio is represented among the recommended builds by a single row, and by accident.

\subsection{A pin, reported as moves, rather than a tuned weight\secstatus{todo}}

Two repairs were available. A weight on production relevance folded into the score would
reorder the table, and would also make the ordering unauditable: a reader could no longer
recover what scale alone said, and the weight would be chosen by whoever wanted a
particular answer. A pin applied after ranking does the same reordering and stays
reversible, so that is what the generating script does.

The axis divides by the role the molecule plays in the experiment, because that is what
decides whether a row teaches production or teaches survival.

\begin{itemize}[leftmargin=1.2em,itemsep=1pt,topsep=2pt]
  \item \term{direct} -- isoprenol or isopentenol is the measured product, meaning a
    titer. Lifted into the first tranche.
  \item \term{tolerance} -- isoprenol is the stressor and the readout is growth, with no
    titer. Lifted into the first tranche as well, since the molecule of interest is still
    the one being measured against. A row reporting both a titer and a tolerance assay is
    direct, because a row carries one value and the titer is the stronger label.
  \item \term{precursor} -- the measured product is a different isoprenoid or terpenoid,
    so the row carries the mevalonate or MEP pathway that feeds isoprenol without
    carrying the product. Lycopene, pinene, mevalonate and geranic acid are here. Lifted
    above the cut at fifty, not into the first tranche.
  \item \term{analog} -- the stressor is a short-chain alcohol or a biomass inhibitor
    standing in for isoprenol rather than isoprenol itself. Reported on the axis and
    \emph{not} lifted. Substituting one alcohol for another assumes cross-tolerance, and
    a row is not promoted on an assumption.
\end{itemize}

Each lift displaces the weakest row inside the boundary that no pin protects, and every
lift is recorded with its rank before and after. Table~\ref{tab:bpins} prints the rank the
measurement rule gives on its own beside the rank this document prints, so the pin can be
undone row by row. The per-host floor of Sec.~\ref{sec:stage} reads the same protection,
which is what keeps one mechanism from canceling another.

\subsection{What the sweep found, and what it missed\secstatus{todo}}\label{sec:sweep}

A metadata sweep over Europe PMC returned 300 publications for the two hosts, 200
\org{E. coli} and 100 \org{P. putida}, and it is held as a separate artifact with a
separate generating script because none of its record counts have been verified. Every
cell of its instance column reads ``Verify from supplement'', and its own overview states
the limit: ``this is a 300-candidate discovery queue, not 300 download-verified
datasets.'' Table~\ref{tab:bqueue} gives its composition and Table~\ref{tab:bqueueiso}
every isoprenol-relevant row it returned. A queue row is a lead, and promoting one means
fetching the paper, counting its axes and writing a curated row; nothing in the queue may
be cited as though its scale were known.

The sweep flagged six publications as measuring isoprenol directly, four for
\org{P. putida} and two for \org{E. coli}. All four \org{P. putida} papers were already
curated here, which is a useful independent check on coverage. Both \org{E. coli} papers
were absent and are now added, and they matter out of proportion to their size: before
them the table held no isoprenol row for \org{E. coli} at all, so the product existed in
one host only and no cross-host question could be posed.

The sweep also missed three isoprenol papers, which is the more useful finding, and the
third miss is the instructive one. Two were already in this table, a CRISPRi
downregulation campaign and a lignocellulosic-sugar campaign that reports isoprenol
alongside isoprenyl acetate; both are titled for an aviation fuel precursor rather than
for the product. The third sits \emph{in} the queue, at the top of every new row it
returned with a priority score of 70.1, and its isoprenol flag reads none. Its title names
\term{biogasoline}, and reading it settles what that means in its own first sentence:
``Isopentenol (3-methyl-3-buten-1-ol) is an important target compound''. The sweep ranked
the paper first and misclassified its product.

The mechanism is the same in all three cases. Matching titles and abstracts finds a paper
named for its product, not one named for its application or its method, and a molecule
with two accepted names and a fuel-class nickname defeats the match three ways. Three
misses against ten is a 30 percent false-negative rate on the single axis the sweep was
built to find, measured against a list assembled by hand. A scored queue is a source of
leads, and its score is not evidence that it understood the row it ranked.

\subsection{What the same pass added that is not isoprenol\secstatus{todo}}

Fourteen rows were verified and added. Four are the isoprenol and isopentenol rows above.
Three carry the precursor pathway. Of the remaining seven, four change the table's answer
to questions it could not previously answer at all.

\notesec{A measured pairwise axis, which the table did not have} Rachwalski et al.\ 2024
conjugated three essential-gene knockdown constructs into the whole deletion collection,
giving 4,017 deletions crossed against each, and separately crossed the whole 356-guide
collection into one deletion background. That is roughly 12,400 distinct gene pairs, and
the released tables carry a normalized growth value for every pair at every condition
rather than only the hits, so the negatives are present. The pairing is what matters: an
essential gene reachable only by knockdown, crossed with a non-essential gene reachable
only by deletion, is a genotype no deletion collection can produce. It ranks 31st on
measurements and is printed 40th after the pins, which understates it, because a colony-size value is one number and the
axis it spans is two-dimensional.

\notesec{A dose on an essential gene} Hawkins et al.\ 2020 target each essential gene with
100 guides, 10 matched and 9 single-mismatch variants of each, so knockdown is titrated
rather than switched. The knockdown level is a linear-model prediction at a cross-validated
$R^{2}$ of 0.56, not a measurement, so the dose enters as an imputed covariate and carries
its own error. That is a schema requirement, not a caveat to be dropped.

\notesec{A complete matrix, which is rarer than a large one} Royet et al.\ 2025 release a
fold change and a false-discovery-adjusted $q$ for all 5,729 \org{P. putida} genes under
each of four metals, with no missing values, so the 22,916 records include every null. The
design is non-barcoded transposon sequencing, which fixes the addressable genotype at the
\emph{gene}: no per-mutant record can ever be built from it, whatever the effort.

\notesec{Matched transcriptome and translatome} Choe et al.\ 2019 release expression and
ribosome-protected-fragment levels over the same 3,457 genes, which is the densest
per-instance vector in the table. Its genotype side is the weak half: all 31 resequenced
samples are populations, every variant table reports an allele frequency, and the flagship
evolved strain has no variant list of its own, so a per-clone genotype would have to be
manufactured by thresholding frequencies, which the authors never did.

\notesec{One addition the ordering rule still buries, and it is the same defect}
Silvis et al.\ 2021 is the only route to a bacterial cell-shape phenotype under gene
perturbation, and shape is the one modality the built yeast morphology data has no
bacterial counterpart for. It releases four shape summaries for 585 strain-conditions,
distilled from 2.6 million segmented cells. That is 2,340 measurements, so the rule places
it 52nd and the pins move it to 64th, in the reserve either way. Nothing pins it, because the pin built here is for one product
rather than for modality coverage in general, and inventing a second pin unprompted would
be the hand-tuning that Sec.~\ref{sec:iso} argues against. The row is therefore recorded,
reported, and outside the recommended fifty. It is worth saying plainly that this is the
same structural defect isoprenol exposed, showing up on a different axis: a row that is
the unique carrier of a phenotype class is ranked by its volume, and volume is not what
makes it unique. Its images are also undeposited, so the provenance chain for 2.6 million
measured cells ends at a 65-kilobyte workbook.

\clearpage
%% GENERATED FILE -- do not hand-edit.
%% SOURCE: experiments/database/scripts/build_bacteria_candidate_datasets_table.py
\begingroup
\footnotesize
\begin{longtable}{@{}L{40mm} L{13mm} L{17mm} L{32mm} r r L{20mm}@{}}
\caption[]{The isoprenol axis, divided by the role the molecule plays. \emph{Rule} is the rank the measurement ordering gives on its own and \emph{Final} the rank this document prints, so the pin can be undone row by row. \emph{Direct} measures isoprenol or isopentenol as a product, the same molecule under two names; \emph{tolerance} measures growth against it as a stressor; both are lifted into tranche 1. \emph{Precursor} measures another isoprenoid and is lifted above the cut. \emph{Analog} substitutes a different alcohol or inhibitor and is reported but not lifted, because the substitution assumes cross-tolerance.}
\label{tab:bpins}\\
\toprule
\textbf{Dataset} & \textbf{Host} & \textbf{Axis} & \textbf{Product} & \textbf{Rule} & \textbf{Final} & \textbf{Lifted by} \\
\midrule
\endfirsthead
\multicolumn{7}{@{}l}{\footnotesize\emph{Table~\ref{tab:bpins}, continued}}\\
\toprule
\textbf{Dataset} & \textbf{Host} & \textbf{Axis} & \textbf{Product} & \textbf{Rule} & \textbf{Final} & \textbf{Lifted by} \\
\midrule
\endhead
\bottomrule
\endfoot
Carruthers 2025 & \emph{P. putida} & direct & isoprenol & 6 & 6 & -- \\
Foo 2014 isopentenol tolerance & \emph{E. coli} & direct & isopentenol & 92 & 12 & isoprenol \\
Yunus 2026 & \emph{P. putida} & direct & isoprenol & 85 & 13 & isoprenol \\
de Siqueira 2025 & \emph{P. putida} & direct & isoprenol & 84 & 14 & isoprenol \\
Lim 2025 isoprenol TALE & \emph{P. putida} & direct & isoprenol & 77 & 16 & isoprenol \\
Tian 2019 isopentenol CRISPRi & \emph{E. coli} & direct & isopentenol & 74 & 17 & isoprenol \\
Kang 2026 isoprenyl acetate & \emph{P. putida} & direct & isoprenyl acetate, isoprenol & 68 & 18 & isoprenol \\
Wang 2022 P. putida isoprenoids & \emph{P. putida} & direct & isoprenol, epi-isozizaene & 63 & 19 & isoprenol \\
Menasalvas 2025 & \emph{P. putida} & direct & isoprenol & 58 & 20 & isoprenol \\
Wang 2015 isoprenol tolerance & \emph{E. coli} & tolerance & isoprenol & 79 & 15 & isoprenol \\
Niu 2019 pinene evolved & \emph{E. coli} & precursor & pinene & 75 & 48 & precursor \\
Yang 2019 mevalonate from ethanol & \emph{P. putida} & precursor & mevalonate & 65 & 49 & precursor \\
Li 2021 mevalonate flux & \emph{E. coli} & precursor & mevalonate & 61 & 50 & precursor \\
Zou 2022 P. putida furfural & \emph{P. putida} & analog & furfural & 81 & 84 & -- \\
Reyes 2011 n-butanol library & \emph{E. coli} & analog & n-butanol & 90 & 91 & -- \\
\end{longtable}
\endgroup


\clearpage
%% GENERATED FILE -- do not hand-edit.
%% SOURCE: experiments/database/scripts/build_bacteria_discovery_queue.py
\begingroup
\footnotesize
\begin{longtable}{@{}L{34mm} r r r r r@{}}
\caption[]{The discovery queue by class and host. \emph{Curated} counts rows already present in Table~\ref{tab:bfinal}, matched on DOI, so the overlap is measured rather than asserted. \emph{Direct} and \emph{Adjacent} are the source's isoprenol axis. No row carries an instance count, because the source asserts none.}
\label{tab:bqueue}\\
\toprule
\textbf{Class} & \textbf{\emph{E. coli}} & \textbf{\emph{P. putida}} & \textbf{Curated} & \textbf{Direct} & \textbf{Adjacent} \\
\midrule
\endfirsthead
\bottomrule
\endfoot
Joint & 35 & 8 & 9 & 3 & 4 \\
ME breadth & 115 & 67 & 9 & 2 & 13 \\
Scale anchor & 50 & 25 & 11 & 1 & 0 \\
\midrule
\textbf{Total} & \textbf{200} & \textbf{100} & \textbf{29} & \textbf{6} & \textbf{17} \\
\end{longtable}
\endgroup


\clearpage
%% GENERATED FILE -- do not hand-edit.
%% SOURCE: experiments/database/scripts/build_bacteria_discovery_queue.py
\begin{landscape}
\begingroup
\footnotesize
\setlength{\tabcolsep}{3pt}
\begin{longtable}{@{}L{12mm} L{13mm} r L{104mm} L{38mm}@{}}
\caption[]{Every isoprenol-relevant publication the sweep returned, the source's own flag. \emph{Direct} measures isoprenol or isopentenol itself; \emph{adjacent} measures another isoprenoid or terpenoid, so it carries the precursor pathway without the product. \emph{Curated} names the row in Table~\ref{tab:bfinal} when the paper is already there, and is otherwise the remaining work.}
\label{tab:bqueueiso}\\
\toprule
\textbf{Flag} & \textbf{Host} & \textbf{Year} & \textbf{Title} & \textbf{Curated} \\
\midrule
\endfirsthead
\multicolumn{5}{@{}l}{\footnotesize\emph{Table~\ref{tab:bqueueiso}, continued}}\\
\toprule
\textbf{Flag} & \textbf{Host} & \textbf{Year} & \textbf{Title} & \textbf{Curated} \\
\midrule
\endhead
\bottomrule
\endfoot
direct & \emph{E. coli} & 2019 & ``Redirecting Metabolic Flux via Combinatorial Multiplex CRISPRi-Mediated Repression for Isopentenol Production in Escherichia coli.'' & Tian 2019 isopentenol CRISPRi \\
direct & \emph{E. coli} & 2015 & ``Dynamic interplay of multidrug transporters with TolC for isoprenol tolerance in Escherichia coli.'' & Wang 2015 isoprenol tolerance \\
direct & \emph{P. putida} & 2025 & ``Alternate routes to acetate tolerance lead to varied isoprenol production from mixed carbon sources in Pseudomonas putida.'' & de Siqueira 2025 \\
direct & \emph{P. putida} & 2025 & ``Automation and machine learning drive rapid optimization of isoprenol production in Pseudomonas putida.'' & Carruthers 2025 \\
direct & \emph{P. putida} & 2025 & ``Biosensor-driven strain engineering reveals key cellular processes for maximizing isoprenol production in Pseudomonas putida.'' & Menasalvas 2025 \\
direct & \emph{P. putida} & 2025 & ``Evolution-guided tolerance engineering of Pseudomonas putida KT2440 for production of the aviation fuel precursor isoprenol.'' & Lim 2025 isoprenol TALE \\
adjacent & \emph{E. coli} & 2013 & ``Chromosomal evolution of Escherichia coli for the efficient production of lycopene.'' & -- \\
adjacent & \emph{E. coli} & 2021 & ``Fine tuning the glycolytic flux ratio of EP-bifido pathway for mevalonate production by enhancing glucose-6-phosphate dehydrogenase (Zwf) and CRISP...'' & Li 2021 mevalonate flux \\
adjacent & \emph{E. coli} & 2022 & ``Enhancing the Production of Pinene in Escherichia coli by Using a Combination of Shotgun, Product-Tolerance and I-SceI Cleavage Systems.'' & -- \\
adjacent & \emph{E. coli} & 2019 & ``Genomic and transcriptional changes in response to pinene tolerance and overproduction in evolved Escherichia coli.'' & Niu 2019 pinene evolved \\
adjacent & \emph{E. coli} & 2023 & ``Evaluation of Various Escherichia coli Strains for Enhanced Lycopene Production.'' & -- \\
adjacent & \emph{E. coli} & 2024 & ``Rational engineering of Escherichia coli strain for stable and enhanced biosynthesis of pinene.'' & -- \\
adjacent & \emph{E. coli} & 2025 & ``Increased Mevalonate Production Using Engineered Citrate Synthase and Phosphofructokinase Variants of Escherichia coli.'' & -- \\
adjacent & \emph{E. coli} & 2026 & ``Metabolic engineering of endogenous MEP pathway for enhanced lycopene production in Escherichia coli.'' & -- \\
adjacent & \emph{E. coli} & 2025 & ``Enhanced isoprenoid production in Escherichia coli cells harboring the archaeal mevalonate pathway.'' & -- \\
adjacent & \emph{E. coli} & 2025 & ``Synergistic Production of Lycopene and β-Alanine Through Engineered Redox Balancing in Escherichia coli.'' & -- \\
adjacent & \emph{E. coli} & 2026 & ``Heterologous Production of Torularhodin, the Monocyclic Carotenoid with a Terminal Carboxyl Group, in Escherichia coli.'' & -- \\
adjacent & \emph{E. coli} & 2026 & ``Development of a transcription factor-based biosensor strain for reporting α-terpineol production via the alcohol-dependent hemiterpene pathway in ...'' & -- \\
adjacent & \emph{P. putida} & 2019 & ``Mevalonate production from ethanol by direct conversion through acetyl-CoA using recombinant Pseudomonas putida, a novel biocatalyst for terpenoid ...'' & Yang 2019 mevalonate from ethanol \\
adjacent & \emph{P. putida} & 2014 & ``De novo production of the monoterpenoid geranic acid by metabolically engineered Pseudomonas putida.'' & -- \\
adjacent & \emph{P. putida} & 2022 & ``Engineering isoprenoids production in metabolically versatile microbial host Pseudomonas putida.'' & Wang 2022 P. putida isoprenoids \\
adjacent & \emph{P. putida} & 2026 & ``Characterisation of metabolic burden in Pseudomonas putida reveals precursor limitation in heterologous lycopene production.'' & -- \\
adjacent & \emph{P. putida} & 2019 & ``Engineering Pseudomonas putida for isoprenoid production by manipulating endogenous and shunt pathways supplying precursors.'' & -- \\
\end{longtable}
\endgroup
\end{landscape}


%% notes-tex/database/database-expansion-bacteria/sections/schema.tex
%% [[experiments.database.expansion-bacteria]]
%% https://github.com/Mjvolk3/torchcell/tree/main/notes-tex/database/database-expansion-bacteria/sections/schema.tex
%% SOURCE: schema table from
%% experiments/database/scripts/build_bacteria_candidate_datasets_table.py; the
%% validator results quoted in the prose were produced by running the validators,
%% recorded in notes/experiments.database.expansion-bacteria.md

\section{What the schema needs first\secstatus{todo}}\label{sec:schema}

Two changes gate every row in Table~\ref{tab:bfinal}, and a third thing that was
expected to gate them does not. All three were established by running the validators
rather than by reading them, because the question is what the code accepts, not what it
appears to accept.

%% GENERATED FILE -- do not hand-edit.
%% SOURCE: experiments/database/scripts/build_bacteria_candidate_datasets_table.py
\begingroup
\footnotesize
\begin{longtable}{@{}L{52mm} L{46mm} L{72mm}@{}}
\caption[]{What the schema needs before any bacterial row can be written, and how each
was established. The first two are the critical path: every row in
Table~\ref{tab:bfinal} waits on them, and both are additive, so they add classes and
move no served closure. \emph{Evidence} is the result of running the validator, not a
reading of it.}
\label{tab:bschema}\\
\toprule
Change & Where & Evidence \\
\midrule
\endfirsthead
\toprule
Change & Where & Evidence \\
\midrule
\endhead
\bottomrule
\endfoot
\textbf{A bacterial gene-identifier namespace on GenePerturbation, so a b-number or a PP\_ locus tag validates as a systematic gene name} & \file{torchcell/datamodels/schema.py, GenePerturbation.validate_sys_gene_name} & The validator admits only Y[A-P][LR]NNN[WC], QNNNN and YNC[A-Q]NNNN[WC]. Constructing a DeletionPerturbation with b0002, PP\_0001, ECK0002 or thrA raises 'Invalid systematic gene name format'; YAL001C is accepted. Blocks: every row in the table that perturbs a named gene. \\
\addlinespace[5pt]
\textbf{A genomes-tier assembly set per host, so a strain's sequence resolves the way a yeast strain's does} & \file{torchcell/sequence/genome/registry.py and torchcell/sequence/genome/} & The tier defines two assembly sets, both S. cerevisiae (sgd\_S288C\_R64-4-1\_20230830 and peter2018\_1011\_assemblies), and the genome package holds only a scerevisiae subpackage. Blocks: the sequence basis of every row, and so the sequence-reconstruction gate. \\
\addlinespace[5pt]
\textbf{Nothing on the reference side} & \file{torchcell/datamodels/schema.py, ReferenceGenome} & species and strain are free strings, so ReferenceGenome(species='Escherichia coli', strain='MG1655') and the KT2440 equivalent both validate unchanged. Blocks: nothing. \\
\addlinespace[5pt]
\end{longtable}
\endgroup


\notesec{The reference side already admits both hosts} \texttt{ReferenceGenome} carries
\texttt{species} and \texttt{strain} as free strings, so
\texttt{ReferenceGenome(species="Escherichia coli", strain="MG1655")} and the
\org{Pseudomonas putida} KT2440 equivalent both validate with no change at all. The
organism was never hard-coded at the reference; it was hard-coded one level down.

\notesec{The genotype side rejects every bacterial gene name} \texttt{GenePerturbation}
validates \texttt{systematic\_gene\_name} against three patterns, and all three are
\org{S. cerevisiae}: the nuclear open reading frames, the mitochondrial genes, and the
noncoding genes. Constructing a deletion with the \gene{thrA} b-number, a \org{P.
putida} locus tag, an EcoCyc identifier or a bare gene symbol raises \emph{Invalid
systematic gene name format} in each case, while a yeast name is accepted. This is not a
naming inconvenience. \texttt{Genotype} derives its content identity from the sorted
systematic names of its perturbations, so the identifier namespace is load-bearing for
node identity in the graph, and a b-number and a PP\_ locus tag must stay distinguishable
namespaces rather than being merged into one string field.

The fix is additive in the sense the admission gate uses: a bacterial namespace accepted
alongside the yeast patterns adds a validated form and changes no record any served
dataset already holds, so no served closure moves. That is the same shape as the
segregant genotype added for the sixteen-cross yeast panel, where a genotype that carried
no gene edit at all needed a sibling class rather than a widened one.

\notesec{The genomes tier has no bacterial assembly set} The tier that dereferences a
strain to sequence defines two assembly sets, both yeast, and the genome package holds
one organism subpackage. A host needs its reference release pinned there by sha256 before
any row's sequence basis resolves to bytes, so this is ordinary provisioning work rather
than a design question, and it is on the critical path for the same reason the identifier
namespace is.

\subsection{What this buys, and what it does not\secstatus{todo}}\label{sec:schema-scope}

The two changes above are enough to write a deletion, a transposon insertion, a guide
knockdown and a production titer for either host. They are not enough for three things
this table contains, and each carries its own cost in Table~\ref{tab:banalogs}.

A designed combinatorial construct, meaning a ribosome binding site or promoter library
over a heterologous pathway, states what each member was built to be rather than what it
turned out to be. Ingesting one commits to storing the design and a verification status,
which is the same weaker claim the yeast list records for its designed-edit rows.

A resequenced evolved clone is a genotype only where the resequencing exists. An evolved
population that was never sequenced differs from its parent at unmapped positions, so it
fails the sequence-reconstruction gate and is dropped in Table~\ref{tab:bexcluded} rather
than ranked low.

A curated literature corpus is not one dataset. Its records re-serve primary papers, so
the count is not net new until the corpus is split by source identifier, and rows in that
state are marked in Table~\ref{tab:bfinal} rather than counted as reachable.

%% notes-tex/database/database-expansion-bacteria/sections/list.tex
%% [[experiments.database.expansion-bacteria]]
%% https://github.com/Mjvolk3/torchcell/tree/main/notes-tex/database/database-expansion-bacteria/sections/list.tex
%% SOURCE: experiments/database/scripts/build_bacteria_candidate_datasets_table.py

\section{The fifty, in rank order\secstatus{todo}}\label{sec:list}

%% GENERATED FILE -- do not hand-edit.
%% SOURCE: experiments/database/scripts/build_bacteria_candidate_datasets_table.py
\begin{landscape}
\begingroup
\footnotesize
\setlength{\tabcolsep}{3pt}
\renewcommand{\arraystretch}{1.15}
\begin{longtable}{@{}r@{\hspace{3pt}} L{38mm} L{13mm} L{10mm} L{19mm} L{17mm} r@{\hspace{4pt}} r@{\hspace{4pt}} L{19mm} L{74mm}@{}}
\caption[]{The bacterial candidates, ranked by \emph{Meas.} descending, which is
instances times phenotype dimensionality (Sec.~\ref{sec:rule}), with the isoprenol rows
then lifted out of that order (Table~\ref{tab:bpins}). \emph{Genotypes} and
\emph{Env} are the perturbation and condition axes. \emph{Inst.} is
genotype$\times$environment records, $\dagger$ where it is the product of the two axes
rather than a reported count and $\ddagger$ where it is an order-of-magnitude estimate.
\emph{Sequence basis} is the route to each strain's total genomic content; a row with no
route is excluded (Table~\ref{tab:bexcluded}). Superscript \textbf{B} marks a row whose
per-record values are not released, \textbf{A} a corpus that re-serves other papers and
is not net new until split by source. For the fifty recommended builds, rows 1--50, citations and data locations
are in Table~\ref{tab:bsources} and the yeast dataset each row mirrors is in
Table~\ref{tab:banalogs}; the reserve rows below the cut appear in neither.}
\label{tab:bfinal}\\
\toprule
\textbf{\#} & \textbf{Dataset (class; phenotype)} & \textbf{Host} & \textbf{Tier} & \textbf{Genotypes} & \textbf{Env} & \textbf{Inst.} & \textbf{Meas.} & \textbf{Sequence basis} & \textbf{Why} \\
\midrule
\endfirsthead
\multicolumn{10}{@{}l}{\footnotesize\emph{Table~\ref{tab:bfinal}, continued}}\\
\toprule
\textbf{\#} & \textbf{Dataset (class; phenotype)} & \textbf{Host} & \textbf{Tier} & \textbf{Genotypes} & \textbf{Env} & \textbf{Inst.} & \textbf{Meas.} & \textbf{Sequence basis} & \textbf{Why} \\
\midrule
\endhead
\bottomrule
\endfoot
1 & \textbf{Fuhrer 2017}\newline {\scriptsize Metabolome / flux; non-targeted flow-injection TOF-MS metabolite ion intensities} & \org{E. coli} & 1 & 3,807 deletions & 1 condition & $3.4\times 10^{4}$ & $2.6\times 10^{8}$ & K-\allowbreak 12-\allowbreak KO & 3,807 Keio single-gene deletion mutants (from 4,320 after excluding 21 very sick strains and 16 injection failures), two independent clones each in technical duplicate, more than 34,000 mass-spectrometric injections. 7,534 distinct m/z features (3,169 negative plus 4,365 positive mode); 3,130 of them were putatively matched to 1,432 of 2,028 chemical formulas, so annotated depth is far below raw feature depth. Deposited in MassIVE and BioStudies, not MetaboLights. \\
\addlinespace[5pt]
2 & \textbf{Wetmore 2015}\newline {\scriptsize Transposon fitness; barcode competitive fitness (BarSeq)} & \org{E. coli} & 1 & 152,018 insertion mutants & 89 conditions & $1.4\times 10^{7}$$\dagger$ & $1.4\times 10^{7}$ & K-\allowbreak 12+\allowbreak transposon & Table 1 gives 152018 strains with unique bar codes for the E. coli library KEIO\_ML9 and fitness estimates for 3471 genes (84 percent). The E. coli lane included 96 samples with 4 time-zero samples across 47 different carbon or nitrogen sources, of which 89 passed quality metrics. Aggregated to gene level the usable record count is 3471 by 89, about 308919. Across all five bacteria the paper reports 387 successful genome-wide assays. No SRA accession is stated in the text; the LBL supplemental URL was taken from the paper and was not fetched. \\
\addlinespace[5pt]
3 & \textbf{Mutalik 2020}\newline {\scriptsize Transposon fitness; barcode competitive fitness (RB-TnSeq BarSeq)} & \org{E. coli} & 1 & 152,018 insertion mutants & 77 conditions & $1.2\times 10^{7}$$\dagger$ & $1.2\times 10^{7}$ & K-\allowbreak 12+\allowbreak transposon & The K-12 BW25113 RB-TnSeq library carries 152018 barcoded insertions in 3728 genes and was assayed against 14 diverse dsDNA phages in 68 RB-TnSeq assays plus 9 no-phage controls, giving the 77 conditions used here. A parallel BL21 library of about 97000 mutants was assayed against 12 phages in 53 pooled fitness experiments plus 9 controls. The same study also ran CRISPRi and Dub-seq arms on the same phage panel, which makes it a three-modality join on one condition axis. Accessions are quoted from the article's data availability statement; the SRA and figshare pages themselves were not fetched. The related Mutalik 2019 Nature Communications Dub-seq paper (doi 10.1038/s41467-018-08177-8, PMID 30659179) is a barcoded overexpression library of 30558 barcode pairs assayed in 155 fitness experiments over 52 chemicals with SRA PRJNA512427 and figshare 10.6084/m9.figshare.6752753.v1; it is omitted as its own row only because plasmid overexpression has no matching modality value here. \\
\addlinespace[5pt]
4 & \textbf{Tong 2020}\newline {\scriptsize Fitness / chemical genomics; 24 hour solid-agar colony growth curve} & \org{E. coli} & 1 & 3,796 deletions & 30 conditions & $1.1\times 10^{5}$ & $8.2\times 10^{6}$ & K-\allowbreak 12-\allowbreak KO & Kinetic growth measurements for 3796 genes under 30 carbon source conditions in MOPS minimal medium, 1536 colonies per plate on solid agar, giving 113880 individual growth curves and over 8 million data points. Plates were scanned every 20 minutes over 24 hours, so each instance is a time series of roughly 72 points; the 72 figure is derived from the stated interval and duration rather than stated directly. 51 poorly annotated genes showed a low growth phenotype in at least one carbon source. This is the cleanest released Keio by carbon-source matrix and the only row here carrying a real time axis. \\
\addlinespace[5pt]
5 & \textbf{PRECISE-1K}\newline {\scriptsize Transcriptome; RNA-seq expression compendium plus 201 iModulon activities} & \org{E. coli} & 3 & 76 designs & 116 conditions & $1.0\times 10^{3}$ & $4.4\times 10^{6}$ & K-\allowbreak 12-\allowbreak KO & 1,035 samples over 4,257 genes, 5 strain backgrounds, 76 unique gene knockouts, 45 distinct projects, yielding 201 iModulons; a combined public K-12 set reaches 2,710 samples and 194 iModulons. The paper reports no single unique-condition total, so 116 is the sum of its itemized condition variables (9 base media, 18 carbon sources, 38 supplements, 42 heterologous proteins, 4 temperatures, 5 pH values). No GEO SuperSeries was found; Zenodo is the stated repository. \\
\addlinespace[5pt]
6 & \textbf{Carruthers 2025}\newline {\scriptsize Production campaign; isoprenol titer by GC-FID plus paired DIA shotgun proteome} & \org{P. putida} & 3 & 472 guides & 1 condition & $1.4\times 10^{3}$$\dagger$ & $2.1\times 10^{6}$ & KT2440+\allowbreak guide & The full text states 472 unique strains (125 single perturbations and 347 combinations) over six DBTL cycles in triplicate, matching the recalled numbers exactly; machine learning selected roughly 400 constructs out of 800,000 possibilities for a fivefold titer gain. ONE study with eight linked deposits: Dryad 10.5061/dryad.gtht76hzh (proteomic metadata plus a 29.7 MB Top3 peptide quantification table, file list enumerated), a Zenodo and GitHub release of per-cycle titer tables (10.5281/zenodo.17178684, JBEI/Isoprenol\_CRISPRi), and seven per-cycle PRIDE accessions PXD063733 (DBTL0), PXD063737, PXD063738, PXD063740, PXD063743, PXD063744, PXD063746 (DBTL6). An ingestion that expects one PXD per paper silently drops six sevenths of the proteomics. \\
\addlinespace[5pt]
7 & \textbf{Lim 2022 putidaPRECISE321}\,\textsuperscript{\textbf{A}}\newline {\scriptsize Transcriptome; batch-normalized log2 TPM RNA-seq expression, decomposed into 84 iModulons} & \org{P. putida} & 4 & 17 wild type & 118 conditions & 321 & $1.8\times 10^{6}$ & reference-\allowbreak only & All recalled numbers confirmed twice over, from the paper and by recomputation on the released matrix: 321 samples, 118 unique condition groups, 21 projects, a 5,564 gene by 321 sample log\_tpm\_norm matrix, 84 iModulons explaining 75.7 percent of variance with 1,265 member genes. The live iModulonDB API returns the dataset under organism p\_putida, dataset precise321. Replicate structure recomputed as 40 conditions at n=2, 71 at n=3, 7 at n=4, summing to 321 exactly. Gene-level values total 1,786,044. This row is an AGGREGATION of 21 projects: 541 samples were collected and 321 passed QC, of which the paper says 305 were previously published and 16 newly generated. The genotype count of 17 distinct genotype tuples (10 with an annotated perturbation such as relA, finR, crc, fleQ) is an UNDERCOUNT, because engineered muconate producers and ALE clones are encoded in condition strings rather than genotype fields. There is no umbrella GEO accession: the sample table carries 47 distinct SRA or BioProject accessions covering 216 of 321 samples and 9 GEO Series, and 105 samples carry no accession at all, so per-sample raw-data provenance is incomplete by construction. \\
\addlinespace[5pt]
8 & \textbf{Borchert 2024 fModules}\,\textsuperscript{\textbf{A}}\newline {\scriptsize Transposon fitness; gene-level RB-TnSeq relative fitness, log2 barcode ratio against the time-zero baseline} & \org{P. putida} & 1 & 4,732 insertion mutants & 179 conditions & $1.6\times 10^{6}$ & $1.6\times 10^{6}$ & KT2440+\allowbreak transposon & Every recalled number checks out, and the 179 versus 183 discrepancy is real rather than an error: Methods say 183 unique growth conditions were collected, the Abstract says 179 were analyzed, and the released workbook's classifier column has exactly 179 distinct values, so 183 is the collection and 179 is the analyzed matrix. Fitness for 4,732 of 5,564 protein-coding genes. The instance count is VERIFIED, not a product: fModule\_Metadata.xlsx (23.3 MB) was opened and its fitness sheet counted at 4,732 rows by 332 value columns with zero nulls, alongside a t-statistic sheet of identical shape. This row is an AGGREGATION and a superset of the other transposon rows: 254 of 332 samples were pulled from the Fitness Browser and 78 are the Beckham-group sets (Putida\_ML5\_set100 and set101) deposited as SRA PRJNA809672 (15 runs), PRJNA856070 (57 runs) and PRJNA1011287 (39 runs), all counted. Replicate structure: 194 samples from duplicate conditions, 78 from triplicate, 60 singletons. The environment axis is two-slot (condition\_1 plus condition\_2 with units), which is what keeps glucose-plus-stressor tolerance samples distinguishable from sole-carbon catabolism samples; collapsing it to one string destroys that contrast. Caveat: genes lacking a value in any sample were dropped, so condition-specific genes reported by the primary papers are silently absent here. \\
\addlinespace[5pt]
9 & \textbf{Caglar 2017}\newline {\scriptsize Multi-omics campaign; matched RNA-seq and shotgun proteomics plus central carbon fluxes} & \org{E. coli} & 4 & 1 wild type & 34 conditions & 322 & $1.4\times 10^{6}$ & reference-\allowbreak only & 34 growth conditions in E. coli B REL606 with 152 RNA-seq samples, 105 proteomics samples and 65 flux measurements (322 total), spanning exponential and stationary phase plus two time courses from 3 hours to 2 weeks. The paper reports 4,196 distinct mRNAs and proteins as a joint figure, so treat 4,196 as the matched transcript-and-protein gene axis rather than as two independent depths. \\
\addlinespace[5pt]
10 & \textbf{Nichols 2011}\newline {\scriptsize Fitness / chemical genomics; colony size drug-gene fitness score} & \org{E. coli} & 1 & 3,979 deletions & 324 conditions & $1.3\times 10^{6}$$\dagger$ & $1.3\times 10^{6}$ & K-\allowbreak 12-\allowbreak KO & The paper states 3979 mutant strains passing quality control, 324 conditions covering 114 unique stresses, and a total of 13497 significant phenotypes at 5 percent FDR, which is about 1 percent of all condition-gene pairs tested. The URL printed in the paper, http://ecoliwiki.net/tools/chemgen/, now returns HTTP 410 Gone; the same portal is live at the ecoliwiki.org host and serves the full per-strain by per-condition flat file. The Keio background is BW25113, not MG1655 itself. \\
\addlinespace[5pt]
11 & \textbf{Goodall 2018}\newline {\scriptsize Transposon fitness; TraDIS insertion-site read density, essentiality call} & \org{E. coli} & 1 & 901,383 insertion mutants & 1 condition & $9.0\times 10^{5}$$\dagger$ & $9.0\times 10^{5}$ & K-\allowbreak 12+\allowbreak transposon & A transposon mutant library estimated at approximately 3.7 million mutants, sequenced to 8279309 mapped reads over 901383 unique insertion sites, from which 358 genes were identified as essential. Only one condition was used, Luria broth at 37 degrees C to OD600 of 1.0, in two independent DNA extracts TL1 and TL2. 248 genes (59.5 percent) of the essential calls were common to this study, the Keio collection and PEC. The ENA accession is quoted from the paper's data availability statement and the ENA page itself was not fetched. \\
\addlinespace[5pt]
12 & \textbf{Foo 2014 isopentenol tolerance}\newline {\scriptsize Production campaign; isopentenol titer in mg per liter at 48 h, mean and SD over 3 replicates} & \org{E. coli} & 2 & 9 released: 8 tolerance genes plus the control & 1 production condition & 9 & 9 & engineered-\allowbreak chassis & The word biogasoline in the title is isoprenol: the paper opens ``Isopentenol (3-methyl-3-buten-1-ol) is an important target compound'', and isopentenol and isoprenol are the same molecule, so this is an isoprenol row under the older name. It is the only E. coli isoprenol campaign that releases numeric titers. Table 1 gives mg per liter with SD for the 8 tolerance genes that reduced the isopentenol growth lag, metR, ibpA, nrdH, soxS, mdlB, fpr, gidB and yqhD, against an 834 plus or minus 5 mg per liter control, which is what earns tier 2 on one cleared bar as the only released E. coli isoprenol production data. A second layer is genuinely deposited and is larger than the titers: GEO GSE53138 holds 6 arrays, 3 with and 3 without 0.2 percent isopentenol, on a 4,254-probe platform, which is an isoprenol-stress transcriptome rather than a genotype axis because all 6 are one strain. The host is DH1 and the tolerance alleles were amplified from MG1655, so the two backgrounds are mixed within one strain. The 40-gene candidate screen behind the 8 winners is figure-only, and the methods state sampling at 150 and 390 minutes while only the 150-minute arrays were deposited. \\
\addlinespace[5pt]
13 & \textbf{Yunus 2026}\newline {\scriptsize Production campaign; isoprenol titer under predicted CRISPRi knockdowns} & \org{P. putida} & 4 & 20 guides & 1 condition & 60$\ddagger$ & 60 & KT2440+\allowbreak guide & Found while sweeping for CRISPRi isoprenol work and it clearly belongs: FluxRETAP and VAMMPIRE were used to pick downregulation targets and to assemble arrays of up to five sgRNAs, reaching nearly 1.5 g/L isoprenol by knocking down PP\_4118 (alpha-ketoglutarate dehydrogenase). Bibliography verified through NCBI eutils. Strain and instance counts are estimates: the abstract does not state them and the full text is paywalled. No accession is given in the indexed metadata, so data\_location is deliberately empty rather than guessed. Shares authors and the pIY670-era pathway with Carruthers 2025 and Banerjee 2024, so genotype deduplication against those rows is required. \\
\addlinespace[5pt]
14 & \textbf{de Siqueira 2025}\newline {\scriptsize Tolerance / robustness; growth rate on glucose-acetate mixtures plus isoprenol titer by GC-FID} & \org{P. putida} & 4 & 7 resequenced clones & 3 conditions & 63$\dagger$ & 63 & evolved-\allowbreak WGS & Both recalled accessions verified by fetching the repository records. PXD055153 is titled Restoration of growth and isoprenol production in dual carbon source media by an acetate-sensitive Pseudomonas putida strain, submitted 2024-08-23, linked to this DOI. PRJNA1153078 (Pseudomonas putida KT2440 tolerization on acetate and isoprenol) holds exactly 5 SRA experiments and 5 BioSamples, matching the 5 resequenced isolates, four tolerized Sigma-class strains plus the parental producer; 2 of the 7 characterized strains were not resequenced and are blocked at the genotype level. ONE study with linked assets PXD055153 plus PRJNA1153078. Proteome panel size is never stated numerically, so dimensionality is kept at 1 rather than guessed. \\
\addlinespace[5pt]
15 & \textbf{Wang 2015 isoprenol tolerance}\newline {\scriptsize Tolerance / robustness; OD600 at 12 h with and without isoprenol, sample SD over 2 biological replicates} & \org{E. coli} & 2 & 47 released: wild type plus 46 deletions & 2 isoprenol doses & 94$\dagger$ & 94 & K-\allowbreak 12-\allowbreak KO & The only isoprenol tolerance screen in either host over a cataloged deletion collection, which is what earns tier 2 on one cleared bar. Supplementary Table S3 releases OD600 with standard deviation for the wild type and 46 Keio single deletions at 0 and 0.5 percent isoprenol by volume, in 2YT at 30 degrees, and Table S4 adds a 9-gene transporter transcript panel. Keio JW numbers are given per strain, so the deletion alleles resolve to Baba 2006. Two multi-gene strains exist and are the paper's own epistasis test, acrA with acrB and that double with tolC, both reported as non-additive against their singles with a plasmid complementation control; neither appears in any released table, so the combinatorial arm is a provenance gap rather than a record. The 0.75 percent dose is figure-only for the same reason. The paper's mutant count is self-inconsistent, 44 in the results and 45 in the abstract against 46 rows in Tables S2 and S3; the table count is what is ingested. \\
\addlinespace[5pt]
16 & \textbf{Lim 2025 isoprenol TALE}\newline {\scriptsize Tolerance / robustness; maximum specific growth rate under isoprenol challenge plus isoprenol titer with the pIY670 plasmid} & \org{P. putida} & 4 & 46 resequenced clones & 2 conditions & 114$\dagger$ & 114 & evolved-\allowbreak WGS & The recalled design is correct and reported: WT, IPL300 and IPL400 each evolved in four independent lineages (12 isoprenol lineages), plus four extra WT lineages tolerized against formaldehyde as controls, so 16 lineages total, isoprenol ramped 4 to 8.5 g/L over 143 to 353 generations. Only the 46 whole-genome resequenced strains (16 endpoint clones plus about 30 intermediates) have a reconstructable genotype; the evolved populations themselves do not and are blocked at the genotype level. 158 unique mutations across 73 regions; proteomics quantified an average of 2,375 of 5,565 putative proteins. Linked assets: PXD054609 (confirmed, citing this PMID), GEO GSE281392, and ALEdb project Pputida\_isoprenol\_TALE. IMPORTANT: the stated resequencing accession PRJNA1187681 does NOT resolve at NCBI BioProject, NCBI SRA or ENA as of 2026-09-24, so the raw reads are not retrievable; the mutation calls survive only in Supplementary Dataset 1. \\
\addlinespace[5pt]
17 & \textbf{Tian 2019 isopentenol CRISPRi}\,\textsuperscript{\textbf{B}}\newline {\scriptsize Combinatorial design; isopentenol titer by GC plus OD600, replicate count not stated} & \org{E. coli} & 4 & 24: 1 base, 18 single-guide, 3 two-guide, 2 three-guide & 3 induction levels & 84$\dagger$ & 168 & engineered-\allowbreak chassis+\allowbreak guide & The only isoprenol production campaign in either host with a nested combinatorial perturbation axis, and the reason this row is here despite releasing nothing. Three genes, asnA, gldA and prpE, are realized as 3 singles, all 3 pairs and the full triple, each assayed across repression levels of 0, 5 and 10 nM anhydrotetracycline, so the design scores an interaction and a dose rather than a ranking. A second 3-gene array over poxB, ackA and pta is matched against a triple knockout of the same genes, which is a knockdown-versus-deletion comparison on one gene set. The host is DH1, not BW25113 or MG1655, carrying the KG1R10 mevalonate pathway as two plasmids: atoB, HMGS, HMGR, MK and PMK under Ptrc, then nudB and PMD under PlacUV5. Every construct has a JBEI registry number. The paper's own target count is self-inconsistent, 21 genes in the introduction against 18 guides in the results and 15 in the abstract, and one printed symbol, arcC, is not a K-12 gene name. \\
\addlinespace[5pt]
18 & \textbf{Kang 2026 isoprenyl acetate}\newline {\scriptsize Production campaign; isoprenyl acetate and isoprenol titer from mixed glucose and xylose} & \org{P. putida} & 4 & 25 pathway designs & 3 conditions & 225$\ddagger$ & 225 & engineered-\allowbreak chassis & Found in a ProteomeXchange sweep rather than the recalled list, and it belongs only if the downstream ester is admitted: the product is isoprenyl acetate, an ATF1-esterified derivative of isoprenol, at 1.9 g/L in fed-batch with a yield of 0.067 g per g total sugar. PXD067010 confirmed by fetching its PRIDE record, which names this DOI and PMID and lists Petzold as lab head. Engineering combines a heterologous alcohol acetyltransferase, esterase deletions to stop product degradation, glucose-xylose co-utilization, and reinforced acetyl-CoA flux, across tube, flask and 2 L fed-batch in biological triplicate. Strain and instance counts are estimates; the deposit does not state a run count. Shares Carruthers and Lee with the flagship isoprenol rows, so deduplicate genotypes. \\
\addlinespace[5pt]
19 & \textbf{Wang 2022 P. putida isoprenoids}\,\textsuperscript{\textbf{B}}\newline {\scriptsize Production campaign; isoprenol, epi-isozizaene, limonene and cineole titers in mg per liter, plus growth and glucose} & \org{P. putida} & 4 & 18 registry-resolved, 16 multi-gene & 20 medium and supplement conditions & 380$\dagger$ & 380 & engineered-\allowbreak chassis & The first isoprenol production in this host, and the row with the most multi-gene genotypes of any isoprenol campaign: 16 of its 18 registry-resolved strains carry either the 3-gene phaABC deletion at PP\_5003 to PP\_5005, a 5-to-8-gene heterologous operon, or both. Best titer 104 mg per liter from 2 percent glucose at 48 h, and 25 mg per liter from p-coumarate, which makes it the only isoprenol row with an aromatic carbon source. Two further deletions are informative for cross-referencing: crc at PP\_5292, and PP\_2675, the isoprenol catabolism gene, which sits immediately beside the PP\_2674 ethanol dehydrogenase that Yang 2019 deletes in the same host. Every per-strain titer is figure-only, which is why this is blocked rather than buildable: the data statement points at JBEI's Experiment Data Depot with no study identifier, and that host returned HTTP 502 when checked. Four further assayed genotypes, two single and one double deletion of ppc and pycAB plus a crc overexpression strain, carry no strain identifier and no locus tags at all. The authors also record that ``the polyploid property nature of P. putida may increase the instability of using a high-copy plasmid'', with colony-to-colony variation observed. \\
\addlinespace[5pt]
20 & \textbf{Menasalvas 2025}\newline {\scriptsize Production campaign; isoprenol titer by GC-FID after biosensor-coupled pyrF growth selection} & \org{P. putida} & 3 & 165 deletions & 2 conditions & $1.3\times 10^{3}$$\dagger$ & $1.3\times 10^{3}$ & engineered-\allowbreak chassis & ATTRIBUTION CORRECTION: no author named Teng is on this paper. The Dryad DOI belongs to Menasalvas et al. with Thomas Eng as co-corresponding author, so the recalled Teng is almost certainly Eng. Over 165 deletion and overexpression strains across 70 gene loci, quadruplicate, titers at 24 and 48 h, reaching about 900 mg/L (36-fold). ONE study, four linked deposits, all confirmed: Dryad 10.5061/dryad.sbcc2frjq (2 files, 78.98 MB), PRIDE PXD061547 (An Isoprenol Biosensor for Combinatorial High Throughput Strain Engineering in Pseudomonas putida, 97 raw files with 8, 24 and 48 h timepoints), BioProject PRJNA1226229 (3 WGS runs), Zenodo 10.5281/zenodo.17155686 (AlphaFold output only). IMPORTANT on the CRISPRi library question: the paper does contain a pooled dCpf1 library of about 16,500 guides (3 per gene over 5,591 coding sequences), but there is NO released per-guide fitness or abundance matrix. Dryad holds only baseline per-guide read counts, a lost-guide list, per-round enriched-guide hit lists, and representative rather than per-replicate ONT reads, so a guide-by-sample matrix cannot be reconstructed. Proteomics covers about 2,500 to 3,000 proteins on a five-context subset, not every strain, so dimensionality is kept at 1. \\
\addlinespace[5pt]
\midrule \multicolumn{10}{@{}l}{\textbf{End of tranche 1. Rows 1--20 are the first builds, 6 per host.}}\\ \midrule
21 & \textbf{Price 2018}\newline {\scriptsize Transposon fitness; barcode competitive fitness (BarSeq log ratio)} & \org{E. coli} & 1 & 3,789 insertion mutants & 162 conditions & $6.1\times 10^{5}$$\dagger$ & $6.1\times 10^{5}$ & K-\allowbreak 12+\allowbreak transposon & The E. coli subset is separable: the published per-organism page for orgId Keio (Escherichia coli BW25113) states 207 condition samples with 162 successful, and its fit\_logratios\_good.tab was downloaded and counted at 162 experiment columns by 3789 gene rows. Whole-study totals are 32 bacteria, 4870 genome-wide fitness experiments meeting consistency criteria, and 11779 poorly annotated protein-coding genes with phenotypes. A 2021 data correction removes the sucrose and D-mannitol experiments for this organism. The live Fitness Browser sits behind a Cloudflare challenge and could not be read, so its current counts may exceed these. \\
\addlinespace[5pt]
22 & \textbf{Choe 2025}\newline {\scriptsize CRISPR library screen; sgRNA abundance by next-generation sequencing} & \org{E. coli} & 2 & 39,591 guides & 12 conditions & $4.8\times 10^{5}$$\dagger$ & $4.8\times 10^{5}$ & K-\allowbreak 12+\allowbreak guide & The initial library holds 39591 sgRNA sequences at 99.96 percent coverage, targeting 4198 coding sequences. Twelve chemical conditions were evaluated: CCCP, polymyxin B, pyocyanin, rifampicin, sulfamethizole, verapamil, puromycin, erythromycin, phleomycin, mitomycin C, MMS and novobiocin. This is the genuine antibiotic-panel CRISPRi screen in E. coli and is the correct source for that claim rather than Rousset 2018. The ENA project was fetched and confirmed, released December 2024, submitted by KAIST. \\
\addlinespace[5pt]
23 & \textbf{Rapp 2026 CRISPRi metabolome}\newline {\scriptsize CRISPR library screen; flow-injection MS metabolome (ion intensities), plus carotenoid titer in the engineering follow-up} & \org{E. coli} & 3 & 1,515 guides & 1 condition & $1.5\times 10^{3}$ & $4.5\times 10^{5}$ & K-\allowbreak 12+\allowbreak guide & Arrayed CRISPRi library covering all 1515 metabolic genes of iML1515, each strain profiled by fast flow-injection MS, so per-design metabolome vectors are recoverable; 36\% of iML1515-predicted metabolites accumulated specifically in some knockdown, and LC-MS/MS spectra were generated for 102 previously uncharacterized metabolites. The exact number of annotated ions per strain was not verified from the full text, so dimensionality 300 is an estimate. The MassIVE record was fetched and confirmed. \\
\addlinespace[5pt]
24 & \textbf{Schmidt 2022 nitrogen}\newline {\scriptsize Transposon fitness; gene-level RB-TnSeq fitness across nitrogen-containing sole nitrogen sources} & \org{P. putida} & 1 & 4,778 insertion mutants & 71 conditions & $3.4\times 10^{5}$$\dagger$ & $3.4\times 10^{5}$ & KT2440+\allowbreak transposon & Not on the requested list but it clearly belongs, and it is the largest single-study condition axis for this organism: assimilation of 52 different nitrogen-containing compounds across 71 tested conditions, with significant fitness phenotypes in 672 genes including 100 transcriptional regulators and 112 transport proteins. Same JBEI-1 library and same Deutschbauer and Keasling pipeline as the lysine, valerolactam and fatty-acid rows, and it is one of the source studies cited by the Borchert 2024 compendium, so it must be deduplicated against that row. The compendium's nitrogen-source experiment group holds 104 of the 332 samples, which is where most of these land. Fitness Browser only, Cloudflare-blocked, so unconfirmed. \\
\addlinespace[5pt]
25 & \textbf{Wang 2018}\newline {\scriptsize CRISPR library screen; sgRNA abundance by Illumina sequencing, gene and sgRNA fitness} & \org{E. coli} & 2 & 56,071 guides & 5 conditions & $2.8\times 10^{5}$$\dagger$ & $2.8\times 10^{5}$ & K-\allowbreak 12+\allowbreak guide & Two corrections to the commonly cited form of this row. The venue is Nature Communications 9:2475, not Science; no Science 2018 genome-wide E. coli CRISPRi paper could be found. The library is 55671 targeting sgRNAs plus 400 negative controls, not about 92000; the 92000-guide figure belongs to the Bikard-lab libraries of Cui 2018 and Rousset 2018. Coverage is at least one sgRNA for 98.6 percent of 4140 protein-coding genes and 79.8 percent of 178 RNA-coding genes. The five screened phenotypes are essentiality in LB, auxotrophy in MOPS versus LB, L-tryptophan biosynthesis, furfural tolerance and isobutanol tolerance. Processed per-gene and per-sgRNA fitness are released as supplementary data; no raw-sequencing accession could be confirmed. \\
\addlinespace[5pt]
26 & \textbf{Rousset 2018}\newline {\scriptsize CRISPR library screen; guide abundance log2 fold change (relative sgRNA fitness)} & \org{E. coli} & 2 & 59,000 guides & 4 conditions & $2.4\times 10^{5}$$\dagger$ & $2.4\times 10^{5}$ & K-\allowbreak 12+\allowbreak guide & The screen starts from a pool of about 92000 sgRNAs targeting random chromosomal positions and is filtered to about 59000 guides. Conditions are growth in rich medium over 17 generations for essentiality plus infection by phages lambda, T4 and 186 at MOI 1, all in triplicate; the paper screens no antibiotics, so any title citing antibiotic exposure for this paper is a conflation. 379 genes were selected for follow-up, including 235 annotated as essential. The ENA accession was verified directly at EBI. \\
\addlinespace[5pt]
27 & \textbf{Shiver 2016}\newline {\scriptsize Fitness / chemical genomics; colony opacity fitness score (Iris)} & \org{E. coli} & 1 & 3,975 deletions & 57 conditions & $2.3\times 10^{5}$$\dagger$ & $2.3\times 10^{5}$ & K-\allowbreak 12-\allowbreak KO & 3975 Keio deletion mutants against 26 new stresses for 57 conditions in total, scored from colony opacity with the Iris image analysis software, and designed as an extension of the Nichols 2011 condition set. The release is raw plate images plus Iris output and ranked score files rather than a clean gene by condition matrix, so building the matrix requires rerunning the published scoring scripts. The Zenodo mirror of the Dryad deposit was fetched and its file list confirmed. \\
\addlinespace[5pt]
28 & \textbf{Thompson 2020 fatty acid and alcohol}\newline {\scriptsize Transposon fitness; gene-level RB-TnSeq fitness on fatty acid and alcohol carbon sources} & \org{P. putida} & 1 & 4,778 insertion mutants & 23 conditions & $2.2\times 10^{5}$$\dagger$ & $2.2\times 10^{5}$ & KT2440+\allowbreak transposon & The recalled 13 fatty acids and 10 alcohols is confirmed verbatim in the abstract and the compounds were enumerated: straight-chain C3 to C10 plus C12 and C14, the esters Tween 20 and butyl stearate, oleic acid; and ethanol, butanol, pentanol, 1,2-propanediol, 1,3-butanediol, 1,4-butanediol, 1,5-pentanediol, isopentanol, isoprenol, 2-methyl-1-butanol. Biological duplicate throughout, corroborated by the compendium metadata showing exactly 2 samples per compound, hence 23 x 2 x 4,778; the condition-collapsed count is about 109,900. Directly relevant to the isoprenol rows: this paper postulates the catabolic routes for isoprenol and isopentanol through leucine metabolism after oxidation and CoA activation, which is the product-loss pathway the isoprenol producers delete. A published correction exists (AEM 87(8):e00177-21, DOI 10.1128/aem.00177-21). Fitness Browser only, Cloudflare-blocked, so unconfirmed; not every published condition survived the compendium's completeness filter, so the full matrix is not recoverable from the workbook alone. \\
\addlinespace[5pt]
29 & \textbf{Borchert 2023 lignin tolerance}\newline {\scriptsize Transposon fitness; gene-level RB-TnSeq fitness under lignin-stream inhibitors on a fixed glucose background} & \org{P. putida} & 1 & 4,732 insertion mutants & 16 conditions & $2.1\times 10^{5}$$\dagger$ & $2.1\times 10^{5}$ & KT2440+\allowbreak transposon & This is the lignin-stream tolerance row with reconstructed engineering strategies, and it is a tolerance rather than a catabolism screen: every stressor is overlaid on M9 plus 20 mM glucose. The 15 stressors were enumerated from the SRA sample titles (ferulate, 4-coumarate, levulinate, beta-ketoadipate, protocatechuate, cis,cis-muconate, glycolate, lactate, acetate, Na2SO4, NaCl, vanillate, vanillin, 4-hydroxybenzoate, 4-hydroxybenzaldehyde) plus a glucose-only control, in triplicate: the 57 SRA runs decompose exactly as 15 x 3 selective, plus 3 plus 3 glucose-control enrichments, plus 3 plus 3 time-zero samples. PRJNA856070 confirmed by fetching the BioProject page (57 SRA experiments, 57 BioSamples). Reconstructed strategies from the fitness hits: delta-gacAS, delta-fleQ, delta-lapAB, delta-ttgR with Ptac-ttgABC, Ptac-PP\_1150-PP\_1152, delta-relA, delta-PP\_1430, several of which also improved growth in a complex lignin-stream mimic. These samples are NOT in the Fitness Browser, which is precisely why the Borchert 2024 compendium reached into SRA for them; the 39 samples of Putida\_ML5\_set100 are recoverable at full precision from that workbook. \\
\addlinespace[5pt]
30 & \textbf{Schastnaya 2021}\newline {\scriptsize Metabolome / flux; metabolome of phosphosite point mutants and matched deletion mutants} & \org{E. coli} & 4 & 89 designs & 5 conditions & 445$\dagger$ & $2.0\times 10^{5}$ & engineered-\allowbreak chassis & 52 known phosphosites on 23 central metabolic enzymes were mutated in E. coli MG1655 delta-mutS, giving 89 phosphomutants profiled on 5 carbon sources against matched Keio deletion strains, with 460 to 500 annotated metabolites by flow-injection TOF-MS. This is the best genotype-by-condition phospho-regulation panel, but it generated no phosphoproteome of its own; its phosphosite list is drawn from earlier studies such as Potel 2018. \\
\addlinespace[5pt]
31 & \textbf{Cui 2018}\newline {\scriptsize CRISPR library screen; guide abundance log2 fold change over 17 generations} & \org{E. coli} & 2 & 92,000 guides & 2 conditions & $1.8\times 10^{5}$$\dagger$ & $1.8\times 10^{5}$ & K-\allowbreak 12+\allowbreak guide & The library is about 92000 unique guide RNAs targeting random positions along the MG1655 genome with an NGG PAM requirement, averaging 19 targets per gene; an exact targeted-gene count is not stated. The two genome-wide screens differ only in dCas9 expression level (strains LC-E18 and LC-E75), both in rich medium with anhydrotetracycline in triplicate, with no stress panel. Data availability names Supplementary Data 4 for the screen results and otherwise defers to the corresponding author on request, so there is no SRA, ENA or GEO accession for this paper. \\
\addlinespace[5pt]
32 & \textbf{Schmidt 2016}\newline {\scriptsize Proteome; absolute protein copies per cell and per cell volume} & \org{E. coli} & 4 & 3 wild type & 22 conditions & 66$\dagger$ & $1.6\times 10^{5}$ & reference-\allowbreak only & 2,359 proteins quantified across all 22 conditions, about 55 percent of predicted ORFs, calibrated to copies per cell with flow cytometry and synthetic peptide standards. Primary strain BW25113, with MG1655 and NCM3722 at selected conditions; biological triplicates for the main dataset, so 66 is a product not a reported sample count. \\
\addlinespace[5pt]
33 & \textbf{Butland 2008}\,\textsuperscript{\textbf{B}}\newline {\scriptsize Genetic interaction; colony size interaction (S) score} & \org{E. coli} & 1 & 155,415 double mutants & 1 condition & $1.6\times 10^{5}$$\ddagger$ & $1.6\times 10^{5}$ & K-\allowbreak 12-\allowbreak KO x KO & Only the citation fields and the abstract-level method description could be confirmed: a quantitative screening procedure based on conjugation of E. coli deletion or hypomorphic strains to create double mutants on a genome-wide scale. The article is closed access with no PMC deposit, and nature.com, Springer static content and mirror sites all refused fetching, so the number of query genes, array strains and screened pairs are all unconfirmed. A search snippet suggesting 39 genome-wide screens, which would give roughly 39 by 3985 or 155415 pairs, could not be verified on any fetched page, so the counts above are an unverified estimate and the row is marked blocked pending a paywalled retrieval. \\
\addlinespace[5pt]
34 & \textbf{Mori 2021}\newline {\scriptsize Proteome; absolute protein mass fractions by DIA/SWATH with xTop inference} & \org{E. coli} & 4 & 3 wild type & 60 conditions & 66 & $1.5\times 10^{5}$ & reference-\allowbreak only & 2,335 proteins with absolute mass fractions over about 60 growth conditions in 66 samples, covering carbon, nitrogen, phosphate and oxygen limitation, non-metabolic stresses and biofilm states, in K-12 NCM3722 and derivatives plus MG1655 sub-strain EQ353 and E. coli Nissle 1917. No ProteomeXchange or PRIDE accession was stated in the text examined; the deposit found is the SWATHAtlas spectral library. An author correction was published in October 2024 and should be read with the paper. \\
\addlinespace[5pt]
35 & \textbf{Ishii 2007}\newline {\scriptsize Multi-omics campaign; paired transcriptome, proteome, metabolome and 13C flux in glucose-limited chemostats} & \org{E. coli} & 4 & 25 designs & 5 conditions & 29 & $1.2\times 10^{5}$ & K-\allowbreak 12-\allowbreak KO & Design confirmed from a secondary open-access source: 24 single-gene disruptants at a fixed dilution rate of 0.2 per hour plus wild type at 5 dilution rates, measured by DNA microarray, 2D-DIGE, CE-TOFMS and metabolic flux analysis. Per-layer depth (transcripts, proteins, metabolites, fluxes) is NOT confirmed: the paper is paywalled and absent from PMC, and no accessible secondary source quotes those counts, so 4,300 is a genome-wide-array estimate for the dominant layer and must be replaced before use in any ranking. \\
\addlinespace[5pt]
36 & \textbf{Campos 2018}\newline {\scriptsize Fitness / chemical genomics; single-cell microscopy morphology and cell cycle features} & \org{E. coli} & 2 & 4,227 deletions & 1 condition & $4.2\times 10^{3}$$\dagger$ & $1.1\times 10^{5}$ & K-\allowbreak 12-\allowbreak KO & 4227 strains of the Keio collection, covering 98 percent of the non-essential genome and 87 percent of the complete genome, imaged in one condition (M9 with 0.1 percent casamino acids and 0.2 percent glucose at 30 degrees C). 26 quantitative features per strain: 19 morphological, 5 cell cycle, and 2 growth-related. On average about 360 cells were imaged per strain, roughly 1.3 million cells after filtering. This is the highest-dimensionality E. coli deletion phenotype release found, and the natural counterpart to a yeast morphology or transcriptome profile rather than to a scalar fitness score. \\
\addlinespace[5pt]
37 & \textbf{Typas 2008}\,\textsuperscript{\textbf{B}}\newline {\scriptsize Genetic interaction; colony pixel count interaction score} & \org{E. coli} & 1 & 91,655 double mutants & 2 conditions & $9.2\times 10^{4}$$\dagger$ & $9.2\times 10^{4}$ & K-\allowbreak 12-\allowbreak KO x KO & This is a method paper (GIANT-coli). Its released per-pair data cover only two demonstration screens: a 12 by 12 cross yielding 66 distinct pairwise double mutants plus 12 self-matings, and the pal and yraP screens in Supplementary Tables 2A and 2B. The genome-wide work is reported only as a linkage curve from 9 by 3985 crosses in LB and 14 by 3985 in M9, which is the 91655 figure above and is a product, not a released matrix. Status is blocked because no genome-wide double-mutant score matrix was released. The toolkit couples the Keio kanamycin deletion library with the ASKA chloramphenicol library, about 4000 single-gene deletions each. \\
\addlinespace[5pt]
38 & \textbf{Girgis 2009}\newline {\scriptsize Transposon fitness; microarray genetic footprinting z-score per locus} & \org{E. coli} & 1 & 500,000 insertion mutants & 17 conditions & $7.3\times 10^{4}$$\ddagger$ & $7.3\times 10^{4}$ & K-\allowbreak 12+\allowbreak transposon & The library is about 5 by 10\textasciicircum{}5 mutants each with a single transposon insertion, selected in 17 antibiotics, read out by hybridizing transposon-junction DNA against genomic DNA on spotted microarrays. The released data are per-locus z-scores, with Dataset S5 the combined z-scores across all loci, so usable records are locus by drug rather than strain by drug; the number of array loci is not stated, so the instance count here is an estimate from an assumed roughly 4300 loci and is labeled as such. No GEO or ArrayExpress accession exists for this dataset; a GEO search returned only the unrelated expression series GSE10855. The companion motility screen is Girgis HS, Liu Y, Ryu WS, Tavazoie S 2007 PLoS Genetics 3(9):e154, doi 10.1371/journal.pgen.0030154, PMID 17941710, same library size, four selection regimes released as 13 supplementary microarray datasets with no accession. \\
\addlinespace[5pt]
39 & \textbf{Banerjee 2025}\newline {\scriptsize Production campaign; indigoidine titer from p-coumarate plus growth rate, with a global DIA proteome} & \org{P. putida} & 4 & 12 promoter variants & 3 conditions & 36$\ddagger$ & $7.2\times 10^{4}$ & KT2440+\allowbreak promoter & This is the recalled genome-scale design-tradeoff campaign, and it is the cleanest promoter-variant genotype axis I found for KT2440: a four-gene growth-coupling cutset for p-coumarate to glutamine to indigoidine, whose fourth gene PP\_0897 (a fumarate hydratase isomer) proved multifunctional and rate-limiting, so its native promoter was swapped for the Anderson-collection pJ23109 or the weak PP\_0415 promoter. PXD050285 verified two ways (ProteomeCentral and the PRIDE v3 API) and its files enumerated: 26 raw MS runs, Orbitrap Exploris 480, DIA searched library-free with DIA-NN at 1 percent global FDR, comprising an 8-run promoter-variant arm (pJ and 0415 driving PP\_0897, 4 replicates each) and an 18-run cross-feeding arm (strains 2370 and 2487 across M9 alanine, M9 alanine malate, M9 p-CA alanine malate, n=3). The about 2,000 quantified proteins is REPORTED verbatim in the paper, not estimated. A separate 190-substrate BIOLOG grid in the same paper is a much larger environment axis if it is ingested. Note the PRIDE submission title differs from the published title; both are correct for their object. \\
\addlinespace[5pt]
40 & \textbf{Mohiuddin 2022 promoter library}\newline {\scriptsize Transcriptome; GFP fluorescence time course over 9 reads, and fold change against untreated} & \org{E. coli} & 1 & 1,930 promoter-GFP reporter strains, 1,809 distinct promoters & untreated plus ampicillin, ofloxacin and gentamicin & $7.7\times 10^{3}$$\dagger$ & $6.9\times 10^{4}$ & K-\allowbreak 12+\allowbreak promoter-\allowbreak reporter & Supplemental Data Set 1 releases the full grid twice over: 69,480 raw GFP readings, 1,930 reporter strains at nine times under four conditions, and 52,110 fold changes against the untreated arm. The library is the E. coli K-12 MG1655 collection of more than 1,900 native promoters fused to a fast-folding GFP on a low-copy plasmid, treated in early stationary phase with ampicillin at 200, ofloxacin at 5 and gentamicin at 50 micrograms per milliliter for five hours. Nothing in the genome is perturbed, which is why it sits beside the screens rather than among them, and it is the reason to take it: the store serves drug-response FITNESS for E. coli from four rows and no drug-response EXPRESSION at all, so this is the half that is missing. Two of its three drugs are also in the Choe 2025 CRISPRi panel, so a gene can be read as both a knockdown fitness cost and an induced promoter under the same compound. The time axis is the second reason: nine reads across the exposure, where the RNA-seq compendia are endpoint. \\
\addlinespace[5pt]
41 & \textbf{Teteneva 2024 lake water}\newline {\scriptsize Transposon fitness; gene-level RB-TnSeq fitness (Wetmore normalization)} & \org{E. coli} & 1 & 3,691 genes carrying insertions & 3 lake water samples x filtered or not x 3 days & $6.6\times 10^{4}$ & $6.6\times 10^{4}$ & K-\allowbreak 12+\allowbreak transposon & Counted off the file rather than the abstract: Supplementary Table S4 holds 11,027 gene by water-sample rows and 66,162 non-empty fitness values across its six columns, filtered and non-filtered water at days 2, 4 and 8, for 3,691 distinct genes. The library is 430,849 unique insertions in 3,833 genes of E. coli K-12 W3110 RpoS+, built by the Wetmore protocol, and fitness is averaged over three technical replicates per sample. Two things make it worth the slot. The host is W3110 rather than the BW25113 that Price 2018, Shiver 2016 and Tong 2020 all sit on, so joining it to them tests whether a gene-level fitness effect survives a change of K-12 background, which no pair of rows already on the list can ask. And the environment is an oligotrophic natural water rather than a defined medium or a drug, which is the only such condition in either host here. The rpoS direction is the opposite sign to the usual stress result, since rpoS mutants gain fitness in filtered water and lose viability, so the row carries a documented growth-versus-survival split that a single fitness number compresses. \\
\addlinespace[5pt]
42 & \textbf{Rachwalski 2024 mobile CRISPRi}\newline {\scriptsize Genetic interaction; normalized colony growth, mean of 2 replicates, no dispersion released} & \org{E. coli} & 1 & 12,404 gene pairs plus 357 knockdowns & 2 media at 6 inducer levels & $6.4\times 10^{4}$ & $6.4\times 10^{4}$ & K-\allowbreak 12-\allowbreak KO+\allowbreak guide & The only measured pairwise genetic-interaction axis at scale in either host, and it is released in full rather than as hits. Three knockdown constructs, of lolA, pssA and mreD, were conjugated into the whole Keio collection, giving 4,017 distinct deletions crossed against each, and the entire 356-guide collection was separately crossed into a lpp deletion. Supplementary Tables S2A, S3 and S4A carry a normalized growth value for every pair at every condition with no gaps, so the negative results are present and not only the 68 suppressors and 9 enhancers of lolA knockdown. The pairing is what makes it valuable here: an essential gene reachable only by knockdown, crossed with a non-essential gene reachable only by deletion, is a combination no deletion collection can produce. Three cautions for the loader. The methods list five inducer levels and omit 100 nanograms per milliliter, which the table headers do carry, so trusting the prose drops a sixth of the dose response. Table S4A holds 4,542 rows for 4,017 unique deletions, with 485 labels repeated up to five times and no plate or well key to separate them. And replicate TYPE is contradictory, called technical in the methods and biological in three figure captions. \\
\addlinespace[5pt]
43 & \textbf{Fang 2025 FFA CRISPRi-FACS}\,\textsuperscript{\textbf{B}}\newline {\scriptsize CRISPR library screen; sgRNA enrichment (FACS sort plus NGS), then titer (g/L) on validated strains} & \org{E. coli} & 2 & 55,671 guides & 1 condition & $5.6\times 10^{4}$ & $5.6\times 10^{4}$ & K-\allowbreak 12+\allowbreak guide & The screen reuses a previously published plasmid library of 55671 sgRNAs, sorted with a fluorescent free-fatty-acid biosensor and read out by NGS; the engineered pcnBi-acrDi-fadR+ strain reached 35.1 g/L FFAs in fed-batch. Pooled enrichment gives a guide-level score rather than a per-design titer, so only the handful of individually reconstructed strains carry phenotype values; no sequencing accession was confirmed. \\
\addlinespace[5pt]
44 & \textbf{Babu 2014}\newline {\scriptsize Genetic interaction; colony size epistasis (S) score} & \org{E. coli} & 1 & 671,071 double mutants & 1 condition & $4.3\times 10^{4}$ & $4.3\times 10^{4}$ & K-\allowbreak 12-\allowbreak KO x KO & The largest E. coli digenic screen that could be verified. It reports over 600000 digenic mutant combinations from 163 query donor genes crossed into an array of 3968 non-essential deletions plus 149 hypomorphic strains, which is 671071 nominal pairs, with donors transferred by conjugation from an Hfr-Cavalli chloramphenicol-marked background. Only the high-confidence subset is released, 25239 aggravating plus 17466 alleviating for 42705 pairs in Table S2, so the full matrix is not recoverable and the instance count reflects the released subset. The dedicated portal http://ecoli.med.utoronto.ca/esga no longer resolves to the Emili lab (the host presents a certificate for lymnaea.org), and BioGRID has no entry for PMID 24586182. The earlier Babu 2011 PLoS Genetics map (doi 10.1371/journal.pgen.1002377, PMID 22125496) is smaller at over 235000 combinations and likewise released only significant pairs. Kumar 2016 Cell Reports (doi 10.1016/j.celrep.2015.12.060, PMID 26774489) covers 102255 viable digenic mutants and is the one E. coli interaction set BioGRID serves, at 31644 curated interactions. \\
\addlinespace[5pt]
45 & \textbf{Gupta 2024 turnover}\newline {\scriptsize Translation / turnover; per-protein degradation and turnover rates} & \org{E. coli} & 3 & 1 wild type & 13 conditions & 13$\ddagger$ & $4.2\times 10^{4}$ & reference-\allowbreak only & About 3,200 proteins, 77 percent of all genes, with turnover rates across 13 growth conditions in strain NCM3722, by heavy 15N dynamic labeling with TMTproC complement-reporter mass spectrometry. Replicates exist but the per-condition replicate count is unconfirmed, so instances is a floor. Cite the 2024 Nature Communications version, not the 2022 bioRxiv preprint. \\
\addlinespace[5pt]
46 & \textbf{MCF2Chem 2023}\,\textsuperscript{\textbf{A}}\newline {\scriptsize Aggregation / support; titer, yield, productivity, content per literature record} & \org{E. coli} & 3 & 8,888 designs & 1 condition & $8.9\times 10^{3}$ & $3.6\times 10^{4}$ & reference-\allowbreak only & 8888 production records covering 1231 compounds in 590 microbial cell factories, with titer, yield, productivity and content plus strain culture information; bacteria are 60\% of species, and E. coli, S. cerevisiae, Y. lipolytica and C. glutamicum together account for 78\% of the compounds, with E. coli alone producing about a quarter of the compounds among the top 20 species. No exact E. coli record count is published, so the E. coli subset size could not be determined. These are counts of literature records, not measured strain-condition observations, and must be resolved to primary source DOIs before they count as independent datasets. \\
\addlinespace[5pt]
47 & \textbf{Choe 2019 genome-reduced ALE}\newline {\scriptsize Multi-omics campaign; matched transcriptome and ribosome-profiling RPKM over 3,457 genes} & \org{E. coli} & 3 & 3 strains, 31 resequenced populations & 1 evolving condition & 10 & $3.5\times 10^{4}$ & evolved-\allowbreak WGS & The only matched transcriptome and translatome in either host, and the densest per-instance vector in the table: Supplementary Data 6 releases RPKM for all 3,457 genes across 6 RNA-seq samples and Supplementary Data 7 the ribosome-protected-fragment RPKM plus translational efficiency over the same gene set, so expression and translation are measured on one strain pair. Supplementary Data 5 adds 839 sigma-70 binding peaks. The genotype side is the problem and it is not a small one: all 31 resequenced samples are POPULATIONS, every variant table gives allele frequency against a timepoint, and the flagship evolved strain has no released variant list of its own, so a per-clone genotype would have to be produced by thresholding day-62 frequencies, a call the authors never made. The reduced parent MS56 has no sequence accession, only a laboratory URL behind a bot wall, so total genomic content is reachable for the MG1655 lineage and not for the reduced one. The evolution environment is non-stationary by design, lysogeny broth supplement falling from 0.1 percent to zero over 62 days, and the released trajectory is cell density rather than growth rate, so early and late passages are not the same condition. Gene identifiers in both expression tables are bare gene names with no locus tag. \\
\addlinespace[5pt]
48 & \textbf{Niu 2019 pinene evolved}\newline {\scriptsize Combinatorial design; ratio of growth and of pinene titer against the unactivated parent, 3 replicates} & \org{E. coli} & 4 & 77 single-target strains plus 2 six-target combinations & 1 pinene dose & 77 & 154 & evolved-\allowbreak WGS & Two properties worth the row. First, a complete released perturbation result rather than a hit list: Supplementary Tables 3 and 4 carry both a growth ratio and a pinene ratio for all 57 activation and all 20 interference targets, with every guide sequence in Table 1, and two six-target combination strains cross activation of flgFGH, sufBCDS, dusB, rpoA, yehA and hslU in one strain and interference of ydiJ, yjbQ, prpR, marR, fabR and cedA in another. Second, a 374-variant table that is the only evolved-strain genotype in either host carrying b-numbers and absolute coordinates with reference and alternate bases, all at frequency 1.00. That table comes with a caveat that must travel with it: reads were aligned to MG1655 while the parent is BW25113, with no stated parental subtraction, so an unknown and probably large share of the 374 calls are background differences between the two K-12 strains rather than evolution-acquired mutations. It is a genotype description against MG1655, not a validated mutation list. The reads themselves were never deposited. No absolute pinene titer appears anywhere, only ratios, and the 182-gene quantitative PCR panel is figure-only. A 2020 erratum attaches to the citation and corrects only a competing-interest statement. \\
\addlinespace[5pt]
49 & \textbf{Yang 2019 mevalonate from ethanol}\,\textsuperscript{\textbf{B}}\newline {\scriptsize Production campaign; mevalonate, ethanol and acetate titers with growth and medium pH} & \org{P. putida} & 4 & 11 strains, 9 multi-gene, up to 9 perturbed loci & 6 conditions including 3 pH levels & 63$\dagger$ & 315 & engineered-\allowbreak chassis & The deepest deletion stacking of any row in either host and the cleanest sequence definition, which is the opposite of what its data statement suggests. Five chromosomal deletions accumulate in order, endA, endX, qedH-I, qedH-II and phaG, each identified by a GenBank gene identifier and each with its upstream and downstream homology-arm primers released, so the deletion BOUNDARIES are pinned rather than implied; the method is markerless with sucrose counter-selection, so no marker is left unaccounted for. All six heterologous genes are released as literal codon-optimized nucleotide sequences totaling 9,283 bases including the ribosome-binding regions. The deepest strain carries 5 deletions plus 4 heterologous genes in one genotype. Two cautions. The data statement reads ``Not applicable'' in full, so every titer is figure-only and only a handful of values with standard deviations survive in prose, which is why this is blocked. And the locus tags are absent from the paper: PP\_3375, PP\_2451, PP\_2674, PP\_2679 and PP\_1408 were DERIVED from the stated gene identifiers rather than quoted, and must be recorded as derived. PP\_2674 sits beside the PP\_2675 isoprenol catabolism gene that Wang 2022 deletes in the same host. \\
\addlinespace[5pt]
50 & \textbf{Li 2021 mevalonate flux}\newline {\scriptsize Metabolome / flux; fitted net flux for 198 reactions with 90 percent confidence bounds} & \org{E. coli} & 4 & 3 strains with a released flux map & 1 labeling condition & 3 & 594 & K-\allowbreak 12-\allowbreak KO & Worth the row for its flux layer alone, which is released in full where the titers are not: Additional file 2 carries 198 reaction rows per strain with a best fit plus lower and upper 90 percent bounds, in both absolute and normalized units, with the fit diagnostics printed, so a carbon-flux phenotype vector can be ingested directly. The perturbation design pairs a 5-level synthetic promoter series over zwf, whose relative strengths are released, with CRISPRi of pfkA in a pfkA-positive host, because deleting pfkA in this background caused a severe growth defect. Serious defects to record. The paper states NO replicate count, has no statistics section and reports no error statistic anywhere, so nothing from it can carry an uncertainty. Its CRISPRi axis is three guide positions, not a strength series, and the repression depth was explicitly never measured, so there is no dose coordinate. Mevalonate titers are figure-only for about 22 of its 27 strains. And the flux supplement covers a different strain set than the text claims: the control's map is absent and a fourth strain's map is present but never discussed, so the released fluxes cannot be matched to the published figure. \\
\addlinespace[5pt]
\midrule \multicolumn{10}{@{}l}{\textbf{End of tranche 2. Rows 21--50 complete the fifty; rows below are the ranked reserve.}}\\ \midrule
51 & \textbf{Hawkins 2020 mismatch-CRISPRi}\newline {\scriptsize CRISPR library screen; relative fitness in doublings against wild type, mean and SD over 4 replicates} & \org{E. coli} & 1 & about 27,000 guides over 270 essential genes & 1 medium & $2.7\times 10^{4}$$\dagger$ & $2.7\times 10^{4}$ & K-\allowbreak 12+\allowbreak guide & A dose axis on an essential gene, which no deletion collection can carry. Each essential gene is targeted by 100 guides, 10 fully matched plus 9 singly mismatched variants of each, so knockdown is titrated rather than switched, and the paper reports a per-gene expression-fitness curve over 17 sliding activity bins for all 270 E. coli essential genes. Table S3 releases per-guide fitness with SD and the 20-nucleotide spacer, which is the element that distinguishes the genotypes, and 1,000 non-targeting controls give a measured noise floor of 0.0825. The guide count is arithmetic over two sourced figures, 270 analyzed genes times 100 guides: the paper states only that the libraries exceed 30,000 elements and never splits that between the two species, so the exact E. coli library size is unstated. Three cautions. The knockdown LEVEL is a linear-model prediction, not a measurement, at a cross-validated R-squared of 0.56, and a preliminary version of that model was used to design the libraries, so the dose must be typed as an imputed covariate. Guides whose fully matched member proved non-functional had their whole series excluded before the released curves, so Table S3 and the curve tables are different populations. And the article is not open access, so the supplementary workbooks need a manual retrieval before they can be mirrored. \\
\addlinespace[5pt]
52 & \textbf{Wang 2024 rifampicin Tn-seq}\newline {\scriptsize Transposon fitness; TRANSIT resampling log2 fold change, with p and adjusted p} & \org{E. coli} & 1 & 4,419 genes & 3 rifampicin multiples of MIC x 2 exposure times & $2.7\times 10^{4}$ & $2.7\times 10^{4}$ & K-\allowbreak 12+\allowbreak transposon & Table S2 is the whole matrix and not the hits: six sheets, one per condition, each carrying all 4,419 genes with insertion-site count, mean control and experimental read depth, log2 fold change, raw p and adjusted p, for 26,514 gene by condition rows. The separate hit table S3 holds 911 genes, so the released negatives outnumber the reported positives by about twenty-nine to one, which is the ratio that decides whether a row can train anything. Host is K-12 MG1655 with a Tn5 library at 71.4 insertions per kilobase, two biological replicates, and the design is a dose by time grid, 0.25x, 4x and 20x MIC at 1 and 3 hours, which is a shape no other row here has: every other drug screen on the list is one dose per drug. It joins Choe 2025 and Shiver 2016 on the drug axis and Girgis 2009 on the modality, so the same gene under the same antibiotic can be read as knockdown, deletion and insertion fitness. \\
\addlinespace[5pt]
53 & \textbf{Royet 2025 KT2440 metal Tn-seq}\newline {\scriptsize Transposon fitness; gene-level log2 fold change of insertion reads, metal against LB, with a permutation q} & \org{P. putida} & 1 & 5,729 genes, 600 essential in LB & 4 metals at one dose each & $2.3\times 10^{4}$$\dagger$ & $2.3\times 10^{4}$ & KT2440+\allowbreak transposon & A complete released matrix, which is rarer than a large one: Supplementary Table S5 carries a log2 fold change and a Benjamini-Hochberg q for all 5,729 genes under each of cobalt, copper, zinc and cadmium, with no missing values, so the 22,916 records include every null and not only the 25 retained hits. Table S4 adds a four-level essentiality call per gene from a hidden Markov model, 600 essential and 4,458 non-essential, on 105,349 within-gene TA sites against 129,002 genome-wide. Sequence provenance is unusually good: the laboratory isolate was itself resequenced for this work and deposited as CP036494 at 100 percent average nucleotide identity to the canonical AE015451.2, and the 12 raw runs map one-to-one onto the six pools by an SRA isolate attribute whose read counts match Table S3 exactly. Two limits are structural rather than fixable. The design is non-barcoded mariner Tn-seq, so the only addressable genotype unit is the GENE and no per-mutant record can ever be built from it, which is the opposite of a barcoded library. And there is no multi-gene axis: every screened genotype carries one insertion, with only a roxS roxR double and a pvdMNOE operon deletion built by hand for validation. All supplementary numbers are stored as text strings under merged headers. \\
\addlinespace[5pt]
54 & \textbf{CeCaFDB flux compendium}\,\textsuperscript{\textbf{A}}\newline {\scriptsize Aggregation / support; curated published 13C flux distributions on a normalized central-carbon network} & \org{E. coli} & 3 & mixed & not stated & 297 & $2.3\times 10^{4}$ & reference-\allowbreak only & 581 flux distributions across 36 organisms from 118 references (1995 to 2013); E. coli alone contributes 297 flux distributions from 32 references, the largest single-organism set (S. cerevisiae is next at 76). The normalized comparison network has 66 metabolites and 76 reactions, which is the per-map dimensionality bound, not necessarily each source study's own network size. The site returned an expired-TLS error on 2026-09-24, so current operation is unconfirmed. \\
\addlinespace[5pt]
55 & \textbf{Thompson 2019 lysine}\newline {\scriptsize Transposon fitness; gene-level RB-TnSeq fitness on glucose, 5-aminovalerate, D-lysine and L-lysine} & \org{P. putida} & 1 & 4,778 insertion mutants & 4 conditions & $1.9\times 10^{4}$$\dagger$ & $1.9\times 10^{4}$ & KT2440+\allowbreak transposon & This is where the name JBEI-1 is defined, verbatim: the JBEI-1 library was created by diluting a 1 mL aliquot of the previously described P. putida RB-TnSeq library (Rand 2017) into 500 mL of LB plus kanamycin. So JBEI-1 is a regrown working stock of Putida\_ML5, not a separate library, which is why every transposon row here shares one join key. Four sole-carbon conditions at 10 mM in MOPS minimal medium, 48 h, with three time-zero aliquots; Methods state no per-condition replicate count. The paper reports 39 genes with fitness below -2 at absolute t above 4. Fitness data are said to be publicly available at fit.genomics.lbl.gov with no per-experiment accession, and that site returns a Cloudflare challenge, so unconfirmed. Strains and plasmids at public-registry.jbei.org folder 391. \\
\addlinespace[5pt]
56 & \textbf{Rand 2017 Putida\_ML5 library}\newline {\scriptsize Modality / backbone; defines the barcoded insertion pool every fitness row is scored on; itself assayed on four carbon sources} & \org{P. putida} & 3 & 185,401 insertion mutants & 4 conditions & $1.9\times 10^{4}$$\ddagger$ & $1.9\times 10^{4}$ & KT2440+\allowbreak transposon & CORRECTION to the recalled library definition on two points. First, the library is named Putida\_ML5 and was built here in the Deutschbauer and Arkin lab from the pKMW3 mariner vector library of Wetmore 2015; JBEI-1 is not a separate library but a regrown, re-aliquoted working stock of Putida\_ML5 named later in Thompson 2019 mBio. The compendium metadata's mutantLibrary column carries exactly two values, Putida\_ML5 (25 samples) and Putida\_ML5\_JBEI (307 samples). Second, the recalled sizes of about 100,000 insertion mutants and about 4,800 nonessential genes are BOTH wrong. The only quantitative description in the literature is in Borchert 2022 (ACS Synth Biol 11:2015-2021, DOI 10.1021/acssynbio.2c00119, PMID 35657709) citing this paper: 185,401 uniquely barcoded transposon insertions, of which 32,591 are intergenic and 152,810 map to 5,213 of 5,661 genes. The about 4,800 figure is a FITNESS-coverage number, not library coverage: the Fitness Browser states fitness data for 4,778 of 5,661 genes, and Borchert 2024 reports 4,732 of 5,564 after filtering. Rand 2017 itself says only thousands of kanamycin-resistant colonies. Its own screen covers 4 carbon sources in duplicate over two days. Join on barcode to mapped insertion to PP\_ locus tag against AE015451 (6,181,873 nt, taxid 160488); no study releases per-barcode fitness, all publish gene-level aggregates. The Fitness Browser organism page is orgId=Putida and reported 314 experiments across 54 carbon sources, 52 nitrogen sources, 2 stress compounds and 27 other conditions, but it is behind a Cloudflare managed JavaScript challenge today, so those figures were read from Internet Archive snapshots (2025-03-10 and 2025-04-24) and must NOT be recorded as live 2026 values. Hence accession\_confirmed false. \\
\addlinespace[5pt]
57 & \textbf{Balakrishnan 2022}\newline {\scriptsize Multi-omics campaign; matched mRNA and protein mass fractions plus promoter activity and mRNA degradation rates} & \org{E. coli} & 4 & 2 designs & 4 conditions & 8$\dagger$ & $1.5\times 10^{4}$ & reference-\allowbreak only & The cleanest transcript-and-protein join on identical samples: more than 1,900 proteins with matched mRNA, plus about 2,700 genome-wide mRNA degradation rates, in E. coli K-12 NCM3722 (and a delta-rsd mutant) across a reference glucose minimal medium plus carbon limitation, anabolic limitation and translational inhibition, growth rates 0.3 to 0.9 per hour. Replicate counts vary by assay and are not uniformly reported, so instances is a product. \\
\addlinespace[5pt]
58 & \textbf{Brunk 2016}\newline {\scriptsize Multi-omics campaign; 86 metabolites plus 55 protein complexes, with product titers (isopentenol, limonene, bisabolene)} & \org{E. coli} & 4 & 9 pathway designs & 9 conditions & 81$\dagger$ & $1.1\times 10^{4}$ & engineered-\allowbreak chassis & Nine strains (eight engineered plus wild-type E. coli DH1) carrying three versions of a heterologous mevalonate pathway, sampled across a 0 to 72 hour batch fermentation: the joint panel is 86 metabolites and 55 protein complexes at 9 time points, while the aggregate metabolomics set is described as 9 strains x 13 time points x 86 metabolites, so instances 81 is the product for the joint panel and the metabolome alone is larger. Variances come from triplicate measurements where available. No omics repository accession was confirmed. \\
\addlinespace[5pt]
59 & \textbf{Thompson 2019 valerolactam}\newline {\scriptsize Transposon fitness; gene-level RB-TnSeq fitness on lactam carbon sources, then valerolactam titer in deletion strains} & \org{P. putida} & 2 & 4,778 insertion mutants & 2 conditions & $9.6\times 10^{3}$$\dagger$ & $9.6\times 10^{3}$ & KT2440+\allowbreak transposon & Two selective RB-TnSeq conditions, 10 mM valerolactam and 10 mM 5-aminovalerate in MOPS minimal medium, with three time-zero aliquots; the two 500 uL aliquots per condition are pooling, not replication, so the paper reports no replicate count. RB-TnSeq is the small half of this paper; the payload worth ingesting is the engineering ladder it produced, wild type 0 mg/L to delta-oplBA 9.27 to delta-oplBA delta-davT 85.19 to delta-oplBA delta-davT delta-alr 91.97 mg/L valerolactam at 48 h, with strains at public-registry.jbei.org folder 456. Fitness lives only in the Fitness Browser with no per-experiment accession, and that site is behind a Cloudflare managed challenge to both WebFetch and curl, so accession\_confirmed is false. The two valerolactam carbon-source samples are however confirmed to exist, appearing by name (set6IT064, set7IT045) in the Borchert 2024 compendium workbook. \\
\addlinespace[5pt]
60 & \textbf{D2Cell 2026}\,\textsuperscript{\textbf{A}}\newline {\scriptsize Aggregation / support; literature-reported production metric per gene-modification record} & \org{E. coli} & 3 & 8,134 designs & 1 condition & $8.1\times 10^{3}$ & $8.1\times 10^{3}$ & reference-\allowbreak only & Confirmed verbatim in the preprint: 8134 experimentally reported single- and double-gene modification records for 73 products in E. coli, merged for training with 11643 single-gene-modification rows for the same 73 products simulated by the genome-scale model (FSEOF-style), so the simulated rows originate from a separate generation step and are separable in principle, but the preprint states no explicit flag distinguishing them in the released file; the whole database is 29006 entries, 1210 products, 751 organisms. These are counts of literature records, not measured strain-condition observations, and must be resolved to primary source DOIs before they count as independent datasets. \\
\addlinespace[5pt]
61 & \textbf{Li 2014}\newline {\scriptsize Translation / turnover; absolute protein synthesis rates from ribosome profiling} & \org{E. coli} & 3 & 1 wild type & 2 conditions & 2$\dagger$ & $6.1\times 10^{3}$ & reference-\allowbreak only & 3,041 genes accounting for more than 96 percent of total protein synthesized in rich defined medium, with a similar number in glucose minimal medium, at 90 million fragments per sample. Method is ribosome profiling calibrated to total protein per doubling; the quantitative mass spectrometry comparison is to published external datasets, not generated here. Strain designation is unconfirmed (main-text methods do not name it). \\
\addlinespace[5pt]
62 & \textbf{Potel 2018 phosphoproteome}\newline {\scriptsize Proteome; phosphosite identifications including phosphohistidine} & \org{E. coli} & 4 & 1 wild type & 2 conditions & 2$\ddagger$ & $4.3\times 10^{3}$ & reference-\allowbreak only & The deepest E. coli phosphoproteome found, but with almost no breadth: one wild-type MG1655 strain on glycerol, exponential versus stationary phase. The 2,129 total phosphosites and 246 phosphohistidine sites on 173 proteins come from a secondary search snippet, not from a fetched primary text; the abstract confirms only that about 10 percent of sites are phosphohistidine, so dimensionality needs primary confirmation before ranking. \\
\addlinespace[5pt]
63 & \textbf{Baba 2006}\newline {\scriptsize Modality / backbone; deletion viability (in-frame knockout obtained or not)} & \org{E. coli} & 3 & 3,985 deletions & 1 condition & $4.0\times 10^{3}$ & $4.0\times 10^{3}$ & K-\allowbreak 12-\allowbreak KO & Of 4288 genes targeted, mutants were obtained for 3985, held in duplicate for 7970 strains; 303 genes including 37 of unknown function could not be disrupted and are the paper's essential-gene candidates. Strain background is E. coli K-12 BW25113. The distribution URL given in the paper (GenoBase, http://ecoli.aist-nara.ac.jp/) is stale; NBRP E. coli currently lists 3909 Keio entries ready to distribute. This is the strain resource and the join axis for Nichols 2011, Tong 2020, Campos 2018, Shiver 2016 and the conjugation-based interaction screens, not a phenotype screen itself. \\
\addlinespace[5pt]
64 & \textbf{Gerdes 2003}\newline {\scriptsize Transposon fitness; genetic footprinting essentiality assertion} & \org{E. coli} & 2 & 3,746 insertion mutants & 1 condition & $3.7\times 10^{3}$ & $3.7\times 10^{3}$ & K-\allowbreak 12+\allowbreak transposon & Unambiguous essentiality assessments were made for 3746 (87 percent) of E. coli protein-encoding genes or ORFs, of which 620 (14 percent) were asserted essential and 3126 (73 percent) non-essential; no assertion was possible for 327 genes on technical grounds and evidence was insufficient for 218. Method is Tn5-based genetic footprinting in rich aerobic media. The per-gene Table S1 is still served live at the Wisconsin mirror, which makes this the cleanest hash-pinnable per-gene essentiality table for MG1655. Note that the Database of Essential Genes lists 609 essential genes for this study and 296 for Baba 2006, both of which disagree with the source papers (620 and 303); use the paper tables. \\
\addlinespace[5pt]
65 & \textbf{Kochanowski 2017}\newline {\scriptsize Multi-omics campaign; promoter activities and absolute central metabolite concentrations on shared conditions} & \org{E. coli} & 4 & 2 wild type & 26 conditions & 26 & $3.7\times 10^{3}$ & reference-\allowbreak only & 95 fluorescent promoter reporters (an existing library expanded by 28, with 31 inactive promoters discarded) measured at steady state across 26 conditions spanning carbon sources, amino-acid supplementation and sub-lethal chloramphenicol, growth rates 0.1 to 1.5 per hour, in BW25113 plus a Crp deletion mutant. 47 central metabolites were quantified absolutely by ion-pairing UPLC-MS/MS in 23 of the 26 conditions, so dimensionality 142 is the union of 95 promoter activities and 47 metabolites, not a per-sample vector measured in every condition. No fluxes were measured. \\
\addlinespace[5pt]
66 & \textbf{Oyetunde 2019}\,\textsuperscript{\textbf{A}}\newline {\scriptsize Aggregation / support; titer (g/L), rate (g/L/h), yield (g/g)} & \org{E. coli} & 3 & 1,200 designs & 1 condition & $1.2\times 10^{3}$ & $3.6\times 10^{3}$ & reference-\allowbreak only & About 1200 experimentally realized E. coli cell factories hand-curated from about 100 papers covering more than 20 products, with features in six categories (carbon source, bioprocess conditions, genetic modifications, product properties, production metrics, undocumented factors) and TRY targets; data availability statement says all data are in the supplementary files. These are literature records, not measured strain-condition observations, so they must be resolved to primary source DOIs before counting as independent datasets. \\
\addlinespace[5pt]
67 & \textbf{Silvis 2021 CRISPRi morphology}\newline {\scriptsize Morphology / imaging; median cell length and width with their coefficients of variation} & \org{E. coli} & 2 & 548 targeting guides over 522 genes & 2 induction levels & 585 & $2.3\times 10^{3}$ & K-\allowbreak 12+\allowbreak guide & The only route to a bacterial cell-shape phenotype under gene perturbation, which is the one modality the built yeast morphology data has no bacterial counterpart for, and that is what earns tier 2 on one cleared bar. Table S3 releases median length, median width and a robust coefficient of variation for each, over 346 uninduced and 239 induced strain-conditions, summarizing 2,646,096 segmented cells of which 2,257,926 passed a stated width filter. Table S2 adds pooled relative fitness with 4 replicate columns and individually measured growth rate and lag. Correcting a common assumption about this paper, there is NO RNA-seq in it: the transcriptional half of the title is a 4-gene quantitative PCR panel plus promoter reporters, both figure-only, so no expression vector can be taken from it. Three further limits. None of the images, contours or per-cell measurements were deposited anywhere, so the provenance chain for 2.6 million measured cells terminates at a 65-kilobyte workbook. The imaging replicate count is never stated. And induction, timepoint and plate format are confounded, uninduced read at 3.5 hours in 96-well plates with a plate-median correction and induced at 5.5 hours in 384-well plates, so the two are not a dose axis. 163 of 548 strains are flagged contaminated and the flag must travel with the record. \\
\addlinespace[5pt]
68 & \textbf{Taniguchi 2010}\newline {\scriptsize Multi-omics campaign; single-cell protein copy number distribution and mRNA copy number per tagged gene} & \org{E. coli} & 3 & 1,018 designs & 1 condition & $1.0\times 10^{3}$ & $2.0\times 10^{3}$ & engineered-\allowbreak chassis & 1,018 confirmed strains of 1,400 attempted, each carrying a chromosomal C-terminal YFP fusion to one native gene at its native locus, imaged at about 4,000 cells per strain (96 strains per microfluidic device, 25 seconds each), with mRNA by RNA-seq and single-molecule FISH (137 genes had matched FISH). Per-strain dimensionality is only 2 summary values, so its value is the 1,018-gene tagged panel and the single-cell distributions, not depth per sample. The CGSC availability claim is unverified. \\
\addlinespace[5pt]
69 & \textbf{Haverkorn van Rijsewijk 2011}\newline {\scriptsize Metabolome / flux; 13C metabolic flux ratios in central carbon metabolism} & \org{E. coli} & 4 & 91 deletions & 2 conditions & 182$\dagger$ & $1.5\times 10^{3}$ & K-\allowbreak 12-\allowbreak KO & 91 regulator deletion mutants (81 transcription factors plus 10 sigma and anti-sigma factors) on glucose and galactose, giving 8 metabolic flux ratios per strain by GC-MS of proteinogenic amino acids with FiatFlux. Dimensionality is genuinely small; the value of this row is the 91-strain deletion axis, not depth. Do not confuse with the Nielsen commentary (DOI 10.1038/msb.2011.10, PMID 21451588). \\
\addlinespace[5pt]
70 & \textbf{Opgenorth 2019}\newline {\scriptsize Multi-omics campaign; 1-dodecanol titer (g/L) plus targeted proteomics of the engineered pathway proteins} & \org{E. coli} & 4 & 60 RBS variants & 2 conditions & 180$\ddagger$ & 900 & engineered-\allowbreak chassis+\allowbreak RBS & Two DBTL cycles over 60 engineered E. coli MG1655 strains; cycle 1 tested 36 strains varying ribosome-binding sites and reductase enzymes in one pathway operon, and cycle 2 designs from the machine-learning models raised titer 21\% to 0.83 g/L. The abstract states dodecanol plus all engineered-pathway protein concentrations were measured; the exact proteomics panel size and replicate count were not verified because the full text is paywalled, so dimensionality 5 and instances 180 are estimates. Conditions counted as 2 for the two DBTL cycles. \\
\addlinespace[5pt]
71 & \textbf{Valencia 2022 free fatty acids}\newline {\scriptsize Production campaign; free fatty acid and fatty acid methyl ester chain-length distribution at 48 h} & \org{P. putida} & 4 & 48 pathway designs & 3 conditions & 144$\ddagger$ & 864 & engineered-\allowbreak chassis & Included because its readout is genuinely multivariate: every sample yields a chain-length vector (C8, C10, C12:1, C14, C16, C16:1) rather than a scalar titer, so encoding it at dimensionality 1 would discard the study. Eight host backgrounds from deletions of three fatty acyl-CoA ligases, crossed with about six plasmid constructs carrying protein-engineered TesA thioesterase variants or MmFAMT, across three media including fourfold-diluted sorghum hydrolysate, n=3; maxima 670.9 mg/L total free fatty acid, 253.6 mg/L C8, 302.4 mg/L total methyl ester, 561.1 mg/L on hydrolysate. The genotype count is an estimate because the host by plasmid design is not fully crossed. Source data for Figures 2 to 6 sit in Supplementary Data 1, but data availability cites only the bare public-registry.jbei.org root with no folder number, so nothing is retrievable by accession. \\
\addlinespace[5pt]
72 & \textbf{Mohamed 2020 hydroxycinnamic TALE}\newline {\scriptsize Tolerance / robustness; growth rate and lag phase on p-coumarate and ferulate at increasing concentrations} & \org{P. putida} & 4 & 105 resequenced clones & 4 conditions & 420$\ddagger$ & 420 & evolved-\allowbreak WGS & Not on the recalled list but it is the best deposited P. putida tolerance evolution campaign. Six independent parallel replicates each for p-coumarate, ferulate and an equal-mass mixture, plus four replicates for a glucose control ALE, ramped to the solubility limit (40 g/L ferulate) or a 67-day cap; outcomes include a 37 h decrease in lag phase at 20 g/L p-coumarate and a 2.4-fold growth-rate increase at 30 g/L ferulate. PRJNA624660 confirmed by fetching the BioProject page: 102 SRA experiments and 105 BioSamples, with the ALEdb project named P. putida Hydroxycinnamic TALE. CAVEAT on genotype reconstruction: the deposit mixes endpoint populations with clonal isolates (1 to 2 intermediate isolates, the endpoint population, and a single endpoint isolate per condition), and only the isolates have a reconstructable single genotype; the population samples are blocked at the genotype level even though their reads are public. \\
\addlinespace[5pt]
73 & \textbf{Fang 2021 FFA CRISPRi}\newline {\scriptsize CRISPR library screen; free fatty acid titer (g/L) by GC} & \org{E. coli} & 4 & 108 guides & 1 condition & 324$\ddagger$ & 324 & K-\allowbreak 12+\allowbreak guide & Arrayed rather than pooled: 108 synthetic sgRNAs repressing 108 chromosomal genes across five metabolic modules, each strain measured for FFA titer by GC in three biological replicates, giving recoverable per-design phenotypes; 30 beneficial genes from the screen plus 26 more from omics, 56 total, and 30.0 g/L in fed-batch. Accessions are named in the paper but were not fetched. \\
\addlinespace[5pt]
74 & \textbf{Jones 2015 ePathOptimize}\newline {\scriptsize Combinatorial design; violacein titer (mg/L) by HPLC} & \org{E. coli} & 4 & 307 promoter variants & 1 condition & 307 & 307 & engineered-\allowbreak chassis+\allowbreak promoter & Five IPTG-inducible mutant T7 promoters of graded strength were placed on each of the five violacein pathway genes vioABCDE in E. coli BL21star(DE3), a theoretical 3125-member library; 107 colonies were screened first (3.4\% of the library) and a second enriched library of 200 mutants was screened, giving 307 genotype-linked titers. Best screen hit 238 mg/L, 63-fold over control, rising to 1829 plus or minus 46 mg/L after fermentation optimization. Deoxyviolacein and pathway intermediates were also quantified for selected strains, so the per-instance panel may exceed 1. \\
\addlinespace[5pt]
75 & \textbf{Jervis 2019}\newline {\scriptsize Combinatorial design; limonene titer (mg per L organic phase)} & \org{E. coli} & 4 & 88 RBS variants & 1 condition & 264$\ddagger$ & 264 & engineered-\allowbreak chassis+\allowbreak RBS & Two RBS libraries in E. coli DH10-beta: pGLlib pairs 12 designed RBS variants each for trAg-gpps and trMs-limS, 144 possible combinations, from which 360 colonies were screened in deep-well plates and 64 clones sequenced, covering 56 of the 144 combinations; a second 5184-member pMVA2 library had 156 colonies screened and a designed 32-combination reduced library built and tested. Genotypes 88 counts only sequence-resolved designs (56 plus 32), since unsequenced colonies have titers without a recoverable genotype; instances assumes triplicate. \\
\addlinespace[5pt]
76 & \textbf{Schmidt 2025 3-hydroxyacid PKS}\newline {\scriptsize Production campaign; LC-MS titer of three 3-hydroxyacids in supernatant at 48 h} & \org{P. putida} & 4 & 12 pathway designs & 2 conditions & 72$\ddagger$ & 216 & engineered-\allowbreak chassis & This is the modular type I PKS row. JBEI public registry folder 887 (Schmidt et al 2025) was confirmed through the REST endpoint: 58 STRAIN entries with public read access, which is an upper bound on strains built and the reason the genotype count is a figure-level estimate rather than a stated total. Three products quantified (3H24DMPA, 3H4MPA at 18.2 mg/L, 3H4MHA at 6.4 mg/L), n=3 for titers, up to seven media variants that are not fully crossed, single 48 h harvest. The genotype classes here exceed the usual four and need extra encoding: alongside deletions and plasmid PKS expression there are mini-Tn7 and BxB1 chromosomal integrations, C-terminal ssrA degron tags on FabD (LAA active, DAS intermediate, LDD inert), and acyltransferase-domain swaps in the LipPKS1 loading module. Accompanying RB-TnSeq is said to be at fit.genomics.lbl.gov, unconfirmed. A sibling deposit from the same group, PXD069956 (Engineering an Extremely Hybrid PKS for Adipic Acid Production, Klass et al., ACS Synth Biol, DOI 10.1021/acssynbio.5c00972, PMID 42339614), covers the same PKS platform in P. putida and E. coli if the modality is worth widening. \\
\addlinespace[5pt]
77 & \textbf{Alper 2005 lycopene}\newline {\scriptsize Combinatorial design; lycopene content per cell dry weight} & \org{E. coli} & 4 & 64 designs & 1 condition & 192$\ddagger$ & 192 & K-\allowbreak 12-\allowbreak KO & Two independent target-discovery routes, a genome-scale-model systematic search and a transposon-insertion combinatorial search, produced two disjoint gene sets whose exhaustive combination gave a defined set of 64 knockout strains, the best 8.5-fold over recombinant K-12 wild type and 2-fold over the engineered parent. The deletion set is in a K-12 background that also carries the heterologous lycopene pathway; transposon insertions were used for discovery, not as the final genotypes. Replicate counts were not verified. \\
\addlinespace[5pt]
78 & \textbf{ActiveOpt valine (Kumar 2021)}\newline {\scriptsize Combinatorial design; valine elemental carbon yield (\% of maximum theoretical yield from glucose plus acetate)} & \org{E. coli} & 4 & 93 RBS variants & 1 condition & 186$\ddagger$ & 186 & engineered-\allowbreak chassis+\allowbreak RBS & New experimental dataset reported in this paper: 39 plasmids expressing ilvBN*DE and ilvIH*C/C*-ygaZH with varied RBS strengths, tested as 89 pairwise plasmid combinations in strain PYR003a (BW25113 delta-aceE delta-gdhA delta-poxB delta-ldhA delta-recA), plus 4 new combinations chosen by ActiveOpt, so 93 genotype-linked yields; best strain 45\% elemental carbon yield, 54.7\% of maximum theoretical. A minimum of two biological replicate colonies per experiment, 48 h shake flask endpoint, no time course. \\
\addlinespace[5pt]
79 & \textbf{Banerjee 2024}\newline {\scriptsize Production campaign; isoprenol titer plus growth and glucose consumption, shake flask and fed-batch} & \org{P. putida} & 4 & 19 deletions & 3 conditions & 171$\ddagger$ & 171 & KT2440-\allowbreak KO & The recalled 3.5 g/L fed-batch figure is confirmed in the abstract, against 1.1 g/L in batch and over tenfold above the starting strain. Eight deletion targets came from bilevel optimization and constrained minimal cut sets on iJN1462, and the best strain carries five deletions. The strain count is an estimate read from the preprint because the published version is paywalled. Proteomics here is targeted MRM of pathway enzymes only, not a discovery panel, so dimensionality stays 1. Linked assets: PXD039868 on PanoramaPublic plus six Experimental Data Depot studies at edd.jbei.org, which were not opened. \\
\addlinespace[5pt]
80 & \textbf{Wang 2018 phosphatase CRISPRi}\newline {\scriptsize CRISPR library screen; lycopene content (and beta-carotene in follow-up strains)} & \org{E. coli} & 4 & 56 guides & 1 condition & 168$\ddagger$ & 168 & engineered-\allowbreak chassis & A CRISPRi knockdown library over 56 phosphatase-encoding genes built in a lycopene overproducer; 28 of the 56 knockdowns impaired lycopene synthesis, and combinatorial knockdowns plus knockouts were then tested in lycopene and beta-carotene overproducers. The gene set is small and fully named, so per-design values should be recoverable from the figures, but replicate counts were not verified so instances is an estimate. \\
\addlinespace[5pt]
81 & \textbf{Jain 2015 1,2-propanediol}\newline {\scriptsize Production campaign; 1,2-propanediol titer (g/L) and yield (g/g glucose)} & \org{E. coli} & 4 & 20 designs & 1 condition & 60$\ddagger$ & 120 & engineered-\allowbreak chassis & I could not confirm the existence of a distinct growth-coupled 1,2-propanediol design-validation campaign: repeated searches surfaced only computational growth-coupling frameworks (OptKnock, minimal cut sets, gcFront) plus separate 1,2-PDO engineering papers, so the 1,2-propanediol row is filled by the best-documented reconstructable E. coli campaign, which reached 5.13 g/L and 0.48 g/g glucose via an optimal minimal enzyme set, flux channeling, raised NADH availability and improved anaerobic growth. The citation is sourced; the genotype and instance counts are unverified estimates, hence confidence recall. \\
\addlinespace[5pt]
82 & \textbf{Farasat 2014 neurosporene (ActiveOpt case 2)}\newline {\scriptsize Combinatorial design; specific neurosporene productivity (microgram per gCDW per hour)} & \org{E. coli} & 4 & 101 RBS variants & 1 condition & 101 & 101 & engineered-\allowbreak chassis+\allowbreak RBS & Kept as a row separate from the valine campaign because the design space is independent (crtEBI RBS library, different host and readout). The neurosporene measurements are Farasat et al. 2014's, not new work: 73 designed crtEBI expression constructs measured as exploration experiments plus 28 kinetic-model-designed extrapolation constructs, productivity rising from 196.3 to 286 microgram per gCDW per hour. ActiveOpt (Kumar 2021, doi 10.1016/j.ymben.2021.06.009) is a retrospective reanalysis of that dataset and generated no new neurosporene strains, so the primary citation is Farasat 2014. \\
\addlinespace[5pt]
83 & \textbf{Lim 2020 ionic liquid ALE}\,\textsuperscript{\textbf{B}}\newline {\scriptsize Tolerance / robustness; growth performance at high protic ionic liquid concentrations} & \org{P. putida} & 4 & 10 resequenced clones & 3 conditions & 90$\ddagger$ & 90 & evolved-\allowbreak WGS & Bibliography confirmed through CrossRef (Green Chemistry 22(17):5677-5690); the journal is not indexed in PubMed, so the PMID is deliberately left empty rather than guessed. Evolved strains grow at up to 4 percent triethanolammonium acetate and 8 percent triethylammonium hydrogen sulfate where the wild type cannot, with mutations in relA, gacS, oprB, fleQ, tktA and uvrY for minimal media and PP\_5350, emrE, oprD and PP\_5324 for the ionic-liquid-specific arms; PP\_5350 (an RpiR-family regulator upregulating the glyoxylate cycle) and the emrE efflux pump were validated by reverse engineering. Marked BLOCKED because no sequencing deposit could be located: the publisher page returns HTTP 403, OSTI dropped the connection, and no BioProject or SRA accession surfaced anywhere, so the evolved genotypes are not retrievable even though the mutations are named in the text. Strain and instance counts are estimates. \\
\addlinespace[5pt]
84 & \textbf{Zou 2022 P. putida furfural}\newline {\scriptsize Tolerance / robustness; growth plus furfural, HMF and acetate concentrations, 2 replicates} & \org{P. putida} & 4 & 20 strains, a 5-gene deletion parent and an evolved derivative & 6 inhibitor conditions & 18 & 72 & evolved-\allowbreak WGS & A cross-host counterpart to the built yeast furfural screen, which is the reason to want it: furfural tolerance measured in a bacterial host on a defined genotype series, against the same inhibitor a yeast CRISPR screen already covers. Supplementary Tables S3 to S6 release the mock-hydrolysate time courses as numbers, with optical density, furfural, HMF and acetate per timepoint for the parent and the evolved strain. The evolved reads are deposited as SRR19970103, which makes this the only evolved row here whose variant calls can be reproduced. The parent is itself a 5-gene deletion, of gcd plus the four gtsABCD genes, and the evolved derivative carries 37 further mutations after four months of escalating furfural. Three limits. Only 11 of the 37 mutations are individually listed; the other 26 are counted by class in Table S7 and never named, so the evolved genotype is incomplete as released and recoverable only by re-calling the deposited reads. The parent was never sequenced and variants were called against the public reference, so the 37 calls conflate furfural selection with the deletion construction and with lab drift, which the paper does not acknowledge. Every recombinant strain phenotype is figure-only, and every uncertainty in the paper is a two-point standard deviation. \\
\addlinespace[5pt]
85 & \textbf{Ling 2022 muconate}\newline {\scriptsize Production campaign; cis,cis-muconate titer, molar yield and volumetric rate on glucose and xylose} & \org{P. putida} & 4 & 11 deletions & 4 conditions & 66$\ddagger$ & 66 & engineered-\allowbreak chassis & The best documented muconate campaign with per-strain titers: 11 strains carry titer or growth measurements and Table 1 lists 24 strains, across four carbon conditions in flask, plate reader and 0.5 L fed-batch, reaching 33.7 g/L at 0.18 g/L/h and 46 percent molar yield. PRJNA783062 confirmed by fetching the BioProject page, titled Sequencing of 5 ALE isolates on xylose (NREL, 2021-11-23), so only the genome resequencing is deposited and there is no proteomics accession. Genotype encoding is unusually mixed and does not fit a single modality: rational deletions and plasmid pathways sit alongside ALE-derived promoter mutations in the PP\_2569 promoter, xylE point mutations (A62V, A455V) and a roughly 227.8 kb genome duplication from PP\_5050 to PP\_5242, none of which are designed edits. \\
\addlinespace[5pt]
86 & \textbf{Werner 2023 beta-ketoadipate}\newline {\scriptsize Production campaign; beta-ketoadipate titer, productivity and molar yield from lignin-derived aromatics} & \org{P. putida} & 4 & 8 deletions & 4 conditions & 64$\ddagger$ & 64 & KT2440-\allowbreak KO & Eight strains with titers across four aromatic feeds (p-coumarate, ferulate, a 3:1 molar mixture, and corn stover alkaline pretreated liquor extractives), reaching 44.5 g/L at 1.15 g/L/h from model compounds and 25 g/L at 0.66 g/L/h from real corn stover aromatics; genotypes are deletions of pcaIJ, crc, pobAR, lvaE, gacA and gacS plus chromosomal overexpression at the fpvA locus. It has the best per-strain numerical table of this cluster, with every main-text data point tabulated in data S1. But there is no repository accession at all and strains are available only from NREL under a material transfer agreement, so accession\_confirmed is false because no accession exists. One arm is informative as a negative: the gacA variants underperformed. \\
\addlinespace[5pt]
87 & \textbf{Banerjee 2020 indigoidine}\newline {\scriptsize Production campaign; indigoidine titer, rate and yield plus growth, scaled from deep-well plate to 2 L bioreactor} & \org{P. putida} & 4 & 2 guides & 8 conditions & 48$\ddagger$ & 48 & KT2440+\allowbreak guide & This is the growth-coupled indigoidine CRISPRi campaign, and it needs a CORRECTION to the recalled framing. The counts are right but the object is not: 63 constrained-minimal-cut-set solutions were analyzed, one feasible cut set hit 14 metabolic reactions, which mapped through GPRs to 16 single-copy genes, of which 2 (mqo-I and cynT) were dropped as essential by genome-wide RB-TnSeq, leaving 14 genes knocked down. Those 14 guides sit in ONE multiplexed dCpf1 gRNA array in a SINGLE strain. There are not 14 CRISPRi strains, so the genotype axis here is 2 (engineered versus empty-vector control), not 14, and anything that ingests this as a 14-strain panel is wrong. Conditions are 2 carbon sources crossed with 4 culture formats. RNA-seq is genome-wide (about 5,500 genes, estimated from the annotation, not stated); the proteomics is targeted SRM, so it is a selected reaction list rather than a discovery panel. BioProjects PRJNA580539 through PRJNA580574 (both endpoints resolve; the paper typos the upper bound); proteomics at PanoramaWeb, not opened. \\
\addlinespace[5pt]
88 & \textbf{Wang 2023 BATCH biosensor CRISPRi}\,\textsuperscript{\textbf{B}}\newline {\scriptsize CRISPR library screen; biosensor fluorescence sort, then titer (mg/L) on selected sgRNA variants} & \org{E. coli} & 4 & 20 guides & 2 conditions & 40$\ddagger$ & 40 & engineered-\allowbreak chassis & Doubly mismatched sgRNA pools give graded knockdown of 20 central-carbon-metabolism genes, screened with a PadR-based p-coumaric acid biosensor and an HpdR-based butyrate biosensor: pfkA plus ptsI variants raised p-coumarate 40.6\% to 1308.6 mg/L from glycerol, and sucA or ldhA variants raised butyrate 19.0\% and 25.2\%. The total number of mismatched sgRNA variants in the pools was not confirmed, and the pooled sort yields no per-variant phenotype except for the few reconstructed winners, hence blocked. Conditions 2 counts the two product-biosensor campaigns. \\
\addlinespace[5pt]
89 & \textbf{Czajka 2022}\newline {\scriptsize Production campaign; indigoidine titer by A612 in DMSO plus growth rate} & \org{P. putida} & 4 & 6 guides & 1 condition & 36$\dagger$ & 36 & KT2440+\allowbreak guide & The direct follow-up to Banerjee 2020 on the same 14-target multiplexed CRISPRi strain, adding an optimized Cpf1 ribosome binding site for a 1.6-fold gain, PHA-operon deletion strains, RNA-seq, targeted proteomics and 13C metabolic flux analysis. Six strains carry titers in one condition (M9, 10 g/L glucose, arabinose plus IPTG), n at least 6 for production profiling, sampled at 24, 48, 72 and 120 h. Kept as its own row rather than folded into Banerjee 2020 because the genotypes and the readout differ. Data availability names only public-registry.jbei.org with no folder number and states explicitly that everything else is in the manuscript and supplementary files, so there is no GEO, SRA, PRIDE or MassIVE deposit and nothing is machine-retrievable. \\
\addlinespace[5pt]
90 & \textbf{Fong 2005 growth-coupled lactate}\newline {\scriptsize Production campaign; lactate titer (g/L), growth rate, lactate secretion rate} & \org{E. coli} & 4 & 3 deletions & 1 condition & 11 & 33 & K-\allowbreak 12-\allowbreak KO & Three OptKnock-designed knockout strains were built and adaptively evolved to yield 11 evolved lactate-producing strains, growing on 2 g/L glucose at 37 C with titers 0.87 to 1.75 g/L and secretion rates coupled to growth rate; computational post-evolution growth-rate predictions were within 10\% in 38 of 50 cases. Only 3 distinct engineered genotypes exist, so this is a small validation campaign rather than a library. \\
\addlinespace[5pt]
91 & \textbf{Reyes 2011 n-butanol library}\newline {\scriptsize Tolerance / robustness; two growth-rate statistics per gene with a t-test p, 4 biological replicates} & \org{E. coli} & 4 & 14 released of 140 assayed & 2 butanol levels & 14 & 28 & K-\allowbreak 12+\allowbreak ORF-\allowbreak plasmid & The only row that scores overexpression and deletion of the same trait side by side from cataloged collections: 55 ASKA overexpression clones and 84 Keio deletions were assayed individually after a pooled enrichment of a roughly 14,000-clone genomic library at sevenfold coverage. Both halves resolve to public collections, so the sequence basis is clean on both signs of perturbation. What limits it badly is release: only the 14 winners carry values, 11 overexpression and 3 deletion, and the 126 other assayed genotypes are not figure-locked but unrecoverable, never plotted per gene. The array-CGH enrichment layer is deposited as GEO GSE26223, 16 two-color arrays on a 4,062-row platform, but its replicates are technical rather than biological and the mapping from the 193 enriched genes to the 55 that were picked is not released. The paper disagrees with itself on the candidate count, 193 in one place and 194 in another, and on the first enrichment rung, 0.5 percent in the results against 0 percent in the methods; both readings must be recorded. One of the 3 deletion hits, rph, is self-flagged as an artifact because the host background already carries an inactivating frameshift there, so that record needs a reliability flag. The overexpression induction level is never stated. \\
\addlinespace[5pt]
92 & \textbf{Yu 2016 PHBA}\newline {\scriptsize Production campaign; para-hydroxybenzoic acid titer from glucose via chorismate, plus carbon yield} & \org{P. putida} & 4 & 7 deletions & 1 condition & 21$\dagger$ & 21 & KT2440-\allowbreak KO & The recalled six strains plus wild type is exact: S0 wild-type control and S1 through S6, built from markerless deletions of pobA, pheA, trpE and hexR combined with plasmid-borne E. coli ubiC and feedback-resistant aroG-D146N; S6 reached 1.73 g/L at 18.1 percent C-mol per C-mol. One defined medium, one 44 h endpoint, biological triplicates, which makes it the cheapest row here to ingest. A single S6 fed-batch run adds a short time series. Up to eight co-measured scalars exist per instance (PHBA, biomass, glucose, acetate, alpha-ketogluconate, pyruvate, lactate, succinate) if each counts as a phenotype. No proteomics, no transcriptomics, no repository deposit of any kind, so accession\_confirmed is false because no accession exists to confirm. \\
\addlinespace[5pt]
\end{longtable}
\endgroup
\end{landscape}


\clearpage
%% GENERATED FILE -- do not hand-edit.
%% SOURCE: experiments/database/scripts/build_bacteria_candidate_datasets_table.py
\begin{landscape}
\begingroup
\footnotesize
\setlength{\tabcolsep}{4pt}
\renewcommand{\arraystretch}{1.15}
\begin{longtable}{@{}r@{\hspace{4pt}} L{86mm} L{74mm} L{78mm}@{}}
\caption[]{Sources for the fifty recommended builds, rows 1--50 of
Table~\ref{tab:bfinal}, in the same order; the ranked reserve below the cut is not listed.
\emph{Data} is where the per-record values live; an entry marked unconfirmed was not
fetched live and must be checked before a loader is written. Every link is clickable.}
\label{tab:bsources}\\
\toprule
\textbf{\#} & \textbf{Dataset, citation and link} & \textbf{Modality and readout} & \textbf{Data} \\
\midrule
\endfirsthead
\multicolumn{4}{@{}l}{\footnotesize\emph{Table~\ref{tab:bsources}, continued}}\\
\toprule
\textbf{\#} & \textbf{Dataset, citation and link} & \textbf{Modality and readout} & \textbf{Data} \\
\midrule
\endhead
\bottomrule
\endfoot
1 & \textbf{Fuhrer 2017}\newline Fuhrer T, Zampieri M, Sevin DC, Sauer U, Zamboni N. Genomewide landscape of gene-metabolome associations in Escherichia coli. Molecular Systems Biology 2017\newline \href{https://doi.org/10.15252/msb.20167150}{doi:10.\allowbreak 15252/\allowbreak msb.\allowbreak 20167150} & gene deletion\newline {\scriptsize endpoint or steady state} & MassIVE MSV000078963; EBI BioStudies S-\allowbreak BSST5; Datasets EV1-\allowbreak EV4 \\
\addlinespace[5pt]
2 & \textbf{Wetmore 2015}\newline Wetmore KM, Price MN, Waters RJ, Lamson JS, He J, Hoover CA, Blow MJ, Bristow J, Butland G, Arkin AP, Deutschbauer A. Rapid quantification of mutant fitness in diverse bacteria by sequencing randomly bar-coded transposons. mBio 2015\newline \href{https://doi.org/10.1128/mBio.00306-15}{doi:10.\allowbreak 1128/\allowbreak mBio.\allowbreak 00306-\allowbreak 15} & transposon insertion\newline {\scriptsize endpoint or steady state} & http:/\allowbreak /\allowbreak genomics.\allowbreak lbl.\allowbreak gov/\allowbreak supplemental/\allowbreak rbarseq/\allowbreak  ; code at https:/\allowbreak /\allowbreak bitbucket.\allowbreak org/\allowbreak berkeleylab/\allowbreak feba\newline {\scriptsize (unconfirmed)} \\
\addlinespace[5pt]
3 & \textbf{Mutalik 2020}\newline Mutalik VK, Adler BA, Rishi HS, Piya D, Zhong C, Koskella B, Kutter EM, Calendar R, Novichkov PS, Price MN, Deutschbauer AM, Arkin AP. High-throughput mapping of the phage resistance landscape in E. coli. PLoS Biology 2020\newline \href{https://doi.org/10.1371/journal.pbio.3000877}{doi:10.\allowbreak 1371/\allowbreak journal.\allowbreak pbio.\allowbreak 3000877} & transposon insertion\newline {\scriptsize endpoint or steady state} & SRA BioProject PRJNA645443; RB-\allowbreak TnSeq data at figshare 10.\allowbreak 6084/\allowbreak m9.\allowbreak figshare.\allowbreak 11413128; CRISPRi data at 10.\allowbreak 6084/\allowbreak m9.\allowbreak figshare.\allowbreak 11859216.\allowbreak v1; Dub-\allowbreak seq data at 10.\allowbreak 6084/\allowbreak m9.\allowbreak figshare.\allowbreak 11838879.\allowbreak v2\newline {\scriptsize (unconfirmed)} \\
\addlinespace[5pt]
4 & \textbf{Tong 2020}\newline Tong M, French S, El Zahed SS, Ong WK, Karp PD, Brown ED. Gene dispensability in Escherichia coli grown in thirty different carbon environments. mBio 2020\newline \href{https://doi.org/10.1128/mBio.02259-20}{doi:10.\allowbreak 1128/\allowbreak mBio.\allowbreak 02259-\allowbreak 20} & gene deletion\newline {\scriptsize sampled series} & Supplementary Tables S1 to S6 via PMC7527729; CarPE web application at https:/\allowbreak /\allowbreak edbrownlab.\allowbreak shinyapps.\allowbreak io/\allowbreak CarPE/\allowbreak  \\
\addlinespace[5pt]
5 & \textbf{PRECISE-1K}\newline Lamoureux CR, Decker KT, Sastry AV, Rychel K, Gao Y, McConn JL, Zielinski DC, Palsson BO. A multi-scale expression and regulation knowledge base for Escherichia coli. Nucleic Acids Research 2023\newline \href{https://doi.org/10.1093/nar/gkad750}{doi:10.\allowbreak 1093/\allowbreak nar/\allowbreak gkad750} & mixed\newline {\scriptsize endpoint or steady state} & https:/\allowbreak /\allowbreak imodulondb.\allowbreak org (dataset E.\allowbreak  coli PRECISE-\allowbreak 1K); https:/\allowbreak /\allowbreak doi.\allowbreak org/\allowbreak 10.\allowbreak 5281/\allowbreak zenodo.\allowbreak 8284223; https:/\allowbreak /\allowbreak github.\allowbreak com/\allowbreak SBRG/\allowbreak precise1k \\
\addlinespace[5pt]
6 & \textbf{Carruthers 2025}\newline Carruthers DN, Kinnunen PC, Li Y, Chen Y, Gin JW, Yunus IS, et al. Automation and machine learning drive rapid optimization of isoprenol production in Pseudomonas putida. Nature Communications 2025\newline \href{https://doi.org/10.1038/s41467-025-66304-8}{doi:10.\allowbreak 1038/\allowbreak s41467-\allowbreak 025-\allowbreak 66304-\allowbreak 8} & CRISPRi\newline {\scriptsize endpoint or steady state} & https:/\allowbreak /\allowbreak doi.\allowbreak org/\allowbreak 10.\allowbreak 5061/\allowbreak dryad.\allowbreak gtht76hzh \\
\addlinespace[5pt]
7 & \textbf{Lim 2022 putidaPRECISE321}\newline Lim HG, Rychel K, Sastry AV, Bentley GJ, Mueller J, Schindel HS, et al. Machine-learning from Pseudomonas putida KT2440 transcriptomes reveals its transcriptional regulatory network. Metabolic Engineering 2022\newline \href{https://doi.org/10.1016/j.ymben.2022.04.004}{doi:10.\allowbreak 1016/\allowbreak j.\allowbreak ymben.\allowbreak 2022.\allowbreak 04.\allowbreak 004} & environment-only\newline {\scriptsize endpoint or steady state} & https:/\allowbreak /\allowbreak github.\allowbreak com/\allowbreak SBRG/\allowbreak modulome\_\allowbreak ppu \\
\addlinespace[5pt]
8 & \textbf{Borchert 2024 fModules}\newline Borchert AJ, Bleem AC, Lim HG, Rychel K, Dooley KD, Kellermyer ZA, Hodges TL, Palsson BO, Beckham GT. Machine learning analysis of RB-TnSeq fitness data predicts functional gene modules in Pseudomonas putida KT2440. mSystems 2024\newline \href{https://doi.org/10.1128/msystems.00942-23}{doi:10.\allowbreak 1128/\allowbreak msystems.\allowbreak 00942-\allowbreak 23} & transposon insertion\newline {\scriptsize endpoint or steady state} & https:/\allowbreak /\allowbreak github.\allowbreak com/\allowbreak beckham-\allowbreak lab/\allowbreak fModule \\
\addlinespace[5pt]
9 & \textbf{Caglar 2017}\newline Caglar MU, Houser JR, Barnhart CS, Boutz DR, Carroll SM, Dasgupta A, Lenoir WF, Smith BL, Sridhara V, Sydykova DK, Vander Wood D, Marx CJ, Marcotte EM, Barrick JE, Wilke CO. The E. coli molecular phenotype under different growth conditions. Scientific Reports 2017\newline \href{https://doi.org/10.1038/srep45303}{doi:10.\allowbreak 1038/\allowbreak srep45303} & environment-only\newline {\scriptsize sampled series} & GSE67402 and GSE94117 (GEO); PXD002140 and PXD005721 (PRIDE); flux data doi:10.\allowbreak 18738/\allowbreak T8/\allowbreak UG3TUR \\
\addlinespace[5pt]
10 & \textbf{Nichols 2011}\newline Nichols RJ, Sen S, Choo YJ, Beltrao P, Zietek M, Chaba R, Lee S, Kazmierczak KM, Lee KJ, Wong A, Shales M, Lovett S, Winkler ME, Krogan NJ, Typas A, Gross CA. Phenotypic landscape of a bacterial cell. Cell 2011\newline \href{https://doi.org/10.1016/j.cell.2010.11.052}{doi:10.\allowbreak 1016/\allowbreak j.\allowbreak cell.\allowbreak 2010.\allowbreak 11.\allowbreak 052} & gene deletion\newline {\scriptsize endpoint or steady state} & https:/\allowbreak /\allowbreak ecoliwiki.\allowbreak org/\allowbreak tools/\allowbreak chemgen/\allowbreak  (coli\_\allowbreak FinalData2.\allowbreak txt, 12.\allowbreak 3 Mb; nichols.\allowbreak tgz, 61.\allowbreak 5 Mb; conditions.\allowbreak txt); Cell Table S2 via PMC3060659 \\
\addlinespace[5pt]
11 & \textbf{Goodall 2018}\newline Goodall ECA, Robinson A, Johnston IG, Jabbari S, Turner KA, Cunningham AF, Lund PA, Cole JA, Henderson IR. The essential genome of Escherichia coli K-12. mBio 2018\newline \href{https://doi.org/10.1128/mBio.02096-17}{doi:10.\allowbreak 1128/\allowbreak mBio.\allowbreak 02096-\allowbreak 17} & transposon insertion\newline {\scriptsize endpoint or steady state} & ENA PRJEB24436\newline {\scriptsize (unconfirmed)} \\
\addlinespace[5pt]
12 & \textbf{Foo 2014 isopentenol tolerance}\newline Foo JL, Jensen HM, Dahl RH, George K, Keasling JD, Lee TS, Leong S, Mukhopadhyay A. Improving microbial biogasoline production in Escherichia coli using tolerance engineering. mBio 2014;5:e01932-14\newline \href{https://doi.org/10.1128/mBio.01932-14}{doi:10.\allowbreak 1128/\allowbreak mBio.\allowbreak 01932-\allowbreak 14} & single-gene overexpression on a production chassis\newline {\scriptsize titers sampled at 24, 48 and 72 h; Table 1 releases 48 h} & GEO GSE53138 (GSM1282891-\allowbreak GSM1282896, platform GPL14649); mBio Table 1 and Tables S2-\allowbreak S4 \\
\addlinespace[5pt]
13 & \textbf{Yunus 2026}\newline Yunus IS, Carruthers DN, Chen Y, Gin JW, Baidoo EEK, Petzold CJ, Garcia Martin H, Adams PD, Mukhopadhyay A, Lee TS. Predictive CRISPR-mediated gene downregulation for enhanced production of sustainable aviation fuel precursor in Pseudomonas putida. Metabolic Engineering 2026\newline \href{https://doi.org/10.1016/j.ymben.2025.11.007}{doi:10.\allowbreak 1016/\allowbreak j.\allowbreak ymben.\allowbreak 2025.\allowbreak 11.\allowbreak 007} & CRISPRi\newline {\scriptsize endpoint or steady state} & \newline {\scriptsize (unconfirmed)} \\
\addlinespace[5pt]
14 & \textbf{de Siqueira 2025}\newline de Siqueira GMV, Srinivasan A, Chen Y, Gin JW, Petzold CJ, Lee TS, Guazzaroni ME, Eng T, Mukhopadhyay A. Alternate routes to acetate tolerance lead to varied isoprenol production from mixed carbon sources in Pseudomonas putida. Applied and Environmental Microbiology 2025\newline \href{https://doi.org/10.1128/aem.02123-24}{doi:10.\allowbreak 1128/\allowbreak aem.\allowbreak 02123-\allowbreak 24} & laboratory evolution\newline {\scriptsize sampled series} & https:/\allowbreak /\allowbreak www.\allowbreak ebi.\allowbreak ac.\allowbreak uk/\allowbreak pride/\allowbreak archive/\allowbreak projects/\allowbreak PXD055153 \\
\addlinespace[5pt]
15 & \textbf{Wang 2015 isoprenol tolerance}\newline Wang C, Yang L, Shah AA, Choi ES, Kim SW. Dynamic interplay of multidrug transporters with TolC for isoprenol tolerance in Escherichia coli. Scientific Reports 2015;5:16505\newline \href{https://doi.org/10.1038/srep16505}{doi:10.\allowbreak 1038/\allowbreak srep16505} & gene deletion\newline {\scriptsize endpoint or steady state} & Springer ESM 41598\_\allowbreak 2015\_\allowbreak BFsrep16505\_\allowbreak MOESM1\_\allowbreak ESM.\allowbreak pdf (Tables S2, S3, S4); strains are Keio JW numbers, NBRP National Institute of Genetics \\
\addlinespace[5pt]
16 & \textbf{Lim 2025 isoprenol TALE}\newline Lim HG, Srinivasan A, Menchavez R, Yunus IS, Noh MH, White M, et al. Evolution-guided tolerance engineering of Pseudomonas putida KT2440 for production of the aviation fuel precursor isoprenol. Metabolic Engineering 2025\newline \href{https://doi.org/10.1016/j.ymben.2025.05.007}{doi:10.\allowbreak 1016/\allowbreak j.\allowbreak ymben.\allowbreak 2025.\allowbreak 05.\allowbreak 007} & laboratory evolution\newline {\scriptsize sampled series} & https:/\allowbreak /\allowbreak www.\allowbreak ebi.\allowbreak ac.\allowbreak uk/\allowbreak pride/\allowbreak archive/\allowbreak projects/\allowbreak PXD054609 \\
\addlinespace[5pt]
17 & \textbf{Tian 2019 isopentenol CRISPRi}\newline Tian T, Kang JW, Kang A, Lee TS. Redirecting metabolic flux via combinatorial multiplex CRISPRi-mediated repression for isopentenol production in Escherichia coli. ACS Synthetic Biology 2019;8:391-402\newline \href{https://doi.org/10.1021/acssynbio.8b00429}{doi:10.\allowbreak 1021/\allowbreak acssynbio.\allowbreak 8b00429} & multiplex CRISPRi knockdown\newline {\scriptsize sampled at 24 and 48 h, with the methods also stating 76 h against 72 h in the results} & JBEI public registry (https:/\allowbreak /\allowbreak public-\allowbreak registry.\allowbreak jbei.\allowbreak org), JPUB numbers per plasmid in Table 1; no sequence-\allowbreak data or value accession\newline {\scriptsize (unconfirmed)} \\
\addlinespace[5pt]
18 & \textbf{Kang 2026 isoprenyl acetate}\newline Kang CW, Carruthers DN, McCauley J, Chen Y, Gin JW, Petzold CJ, Simmons BA, Lee TS. Multi-layered metabolic remodeling of Pseudomonas putida for efficient conversion of lignocellulosic sugars to the precursors of advanced aviation fuel. Metabolic Engineering Communications 2026\newline \href{https://doi.org/10.1016/j.mec.2026.e00274}{doi:10.\allowbreak 1016/\allowbreak j.\allowbreak mec.\allowbreak 2026.\allowbreak e00274} & pathway plasmid\newline {\scriptsize sampled series} & https:/\allowbreak /\allowbreak www.\allowbreak ebi.\allowbreak ac.\allowbreak uk/\allowbreak pride/\allowbreak archive/\allowbreak projects/\allowbreak PXD067010 \\
\addlinespace[5pt]
19 & \textbf{Wang 2022 P. putida isoprenoids}\newline Wang X, Baidoo EEK, Kakumanu R, Xie S, Mukhopadhyay A, Lee TS. Engineering isoprenoids production in metabolically versatile microbial host Pseudomonas putida. Biotechnology for Biofuels and Bioproducts 2022;15:137\newline \href{https://doi.org/10.1186/s13068-022-02235-6}{doi:10.\allowbreak 1186/\allowbreak s13068-\allowbreak 022-\allowbreak 02235-\allowbreak 6} & heterologous pathway plus chromosomal deletion\newline {\scriptsize sampled at 0, 24, 48 h, and to 72 h for the sesquiterpene} & JBEI registry JPUB\_\allowbreak 019914 to JPUB\_\allowbreak 019988; JBEI Experiment Data Depot with no study id (HTTP 502 when checked)\newline {\scriptsize (unconfirmed)} \\
\addlinespace[5pt]
20 & \textbf{Menasalvas 2025}\newline Menasalvas J, Kulakowski S, Chen Y, Gin JW, Akyuz Turumtay E, Baral NR, et al. Biosensor-driven strain engineering reveals key cellular processes for maximizing isoprenol production in Pseudomonas putida. Science Advances 2025\newline \href{https://doi.org/10.1126/sciadv.ady2677}{doi:10.\allowbreak 1126/\allowbreak sciadv.\allowbreak ady2677} & gene deletion\newline {\scriptsize sampled series} & https:/\allowbreak /\allowbreak doi.\allowbreak org/\allowbreak 10.\allowbreak 5061/\allowbreak dryad.\allowbreak sbcc2frjq \\
\addlinespace[5pt]
21 & \textbf{Price 2018}\newline Price MN, Wetmore KM, Waters RJ, Callaghan M, Ray J, Liu H, Kuehl JV, Melnyk RA, Lamson JS, Suh Y, Carlson HK, Esquivel Z, Sadeeshkumar H, Chakraborty R, Zane GM, Rubin BE, Wall JD, Visel A, Bristow J, Blow MJ, Arkin AP, Deutschbauer AM. Mutant phenotypes for thousands of bacterial genes of unknown function. Nature 2018\newline \href{https://doi.org/10.1038/s41586-018-0124-0}{doi:10.\allowbreak 1038/\allowbreak s41586-\allowbreak 018-\allowbreak 0124-\allowbreak 0} & transposon insertion\newline {\scriptsize endpoint or steady state} & https:/\allowbreak /\allowbreak genomics.\allowbreak lbl.\allowbreak gov/\allowbreak supplemental/\allowbreak bigfit/\allowbreak html/\allowbreak Keio/\allowbreak  (fit\_\allowbreak logratios\_\allowbreak good.\allowbreak tab, fit\_\allowbreak t.\allowbreak tab, strain\_\allowbreak fit.\allowbreak tab, all.\allowbreak poolcount); figshare 10.\allowbreak 6084/\allowbreak m9.\allowbreak figshare.\allowbreak 5134840 and 10.\allowbreak 6084/\allowbreak m9.\allowbreak figshare.\allowbreak 5134837; Fitness Browser https:/\allowbreak /\allowbreak fit.\allowbreak genomics.\allowbreak lbl.\allowbreak gov \\
\addlinespace[5pt]
22 & \textbf{Choe 2025}\newline Choe D, Lee E, Kim K, Hwang S, Jeong KJ, Palsson BO, Cho BK, Cho S. Rapid identification of key antibiotic resistance genes in E. coli using high-resolution genome-scale CRISPRi screening. iScience 2025\newline \href{https://doi.org/10.1016/j.isci.2025.112435}{doi:10.\allowbreak 1016/\allowbreak j.\allowbreak isci.\allowbreak 2025.\allowbreak 112435} & CRISPRi\newline {\scriptsize endpoint or steady state} & ENA PRJEB33267 (E.\allowbreak  coli genome-\allowbreak wide CRISPRi under multiple chemical treatments) \\
\addlinespace[5pt]
23 & \textbf{Rapp 2026 CRISPRi metabolome}\newline Rapp J, Verhulsdonk A, Garcke A, Stadelmann A, Farke N, Trossmann F, Kronenberger T, Alvarado A, Petras D, Link H. The metabolome of an E. coli CRISPRi library identifies benefits of minimal metabolite levels and targets for engineering. Cell Systems 2026\newline \href{https://doi.org/10.1016/j.cels.2025.101518}{doi:10.\allowbreak 1016/\allowbreak j.\allowbreak cels.\allowbreak 2025.\allowbreak 101518} & CRISPRi\newline {\scriptsize endpoint or steady state} & MassIVE MSV000095534 ('A metabolism-\allowbreak wide CRISPRi library expands the measurable E.\allowbreak  coli metabolome -\allowbreak  FI-\allowbreak MS dataset', submitted by H.\allowbreak  Link, 2024-\allowbreak 08-\allowbreak 07); Zenodo https:/\allowbreak /\allowbreak zenodo.\allowbreak org/\allowbreak records/\allowbreak 15640293 \\
\addlinespace[5pt]
24 & \textbf{Schmidt 2022 nitrogen}\newline Schmidt M, Pearson AN, Incha MR, Thompson MG, Baidoo EEK, Kakumanu R, et al. Nitrogen Metabolism in Pseudomonas putida: Functional Analysis Using Random Barcode Transposon Sequencing. Applied and Environmental Microbiology 2022\newline \href{https://doi.org/10.1128/aem.02430-21}{doi:10.\allowbreak 1128/\allowbreak aem.\allowbreak 02430-\allowbreak 21} & transposon insertion\newline {\scriptsize endpoint or steady state} & https:/\allowbreak /\allowbreak fit.\allowbreak genomics.\allowbreak lbl.\allowbreak gov/\allowbreak cgi-\allowbreak bin/\allowbreak org.\allowbreak cgi?orgId=Putida\newline {\scriptsize (unconfirmed)} \\
\addlinespace[5pt]
25 & \textbf{Wang 2018}\newline Wang T, Guan C, Guo J, Liu B, Wu Y, Xie Z, Zhang C, Xing XH. Pooled CRISPR interference screening enables genome-scale functional genomics study in bacteria with superior performance. Nature Communications 2018\newline \href{https://doi.org/10.1038/s41467-018-04899-x}{doi:10.\allowbreak 1038/\allowbreak s41467-\allowbreak 018-\allowbreak 04899-\allowbreak x} & CRISPRi\newline {\scriptsize endpoint or steady state} & Supplementary Data 05 to 10 of https:/\allowbreak /\allowbreak www.\allowbreak nature.\allowbreak com/\allowbreak articles/\allowbreak s41467-\allowbreak 018-\allowbreak 04899-\allowbreak x ; code at https:/\allowbreak /\allowbreak github.\allowbreak com/\allowbreak zhangchonglab/\allowbreak CRISPRi-\allowbreak functional-\allowbreak genomics-\allowbreak in-\allowbreak prokaryotes ; library at Addgene (Chong Zhang E.\allowbreak  coli Genome-\allowbreak wide Inhibition Library)\newline {\scriptsize (unconfirmed)} \\
\addlinespace[5pt]
26 & \textbf{Rousset 2018}\newline Rousset F, Cui L, Siouve E, Becavin C, Depardieu F, Bikard D. Genome-wide CRISPR-dCas9 screens in E. coli identify essential genes and phage host factors. PLoS Genetics 2018\newline \href{https://doi.org/10.1371/journal.pgen.1007749}{doi:10.\allowbreak 1371/\allowbreak journal.\allowbreak pgen.\allowbreak 1007749} & CRISPRi\newline {\scriptsize endpoint or steady state} & ENA PRJEB28256 (secondary study ERP110440); supplementary tables S1 to S10 \\
\addlinespace[5pt]
27 & \textbf{Shiver 2016}\newline Shiver AL, Osadnik H, Kritikos G, Li B, Krogan N, Typas A, Gross CA. A chemical-genomic screen of neglected antibiotics reveals illicit transport of kasugamycin and blasticidin S. PLoS Genetics 2016\newline \href{https://doi.org/10.1371/journal.pgen.1006124}{doi:10.\allowbreak 1371/\allowbreak journal.\allowbreak pgen.\allowbreak 1006124} & gene deletion\newline {\scriptsize endpoint or steady state} & Dryad 10.\allowbreak 5061/\allowbreak dryad.\allowbreak f3kc0, mirrored at https:/\allowbreak /\allowbreak zenodo.\allowbreak org/\allowbreak records/\allowbreak 5085832 (imagefile.\allowbreak tar 10.\allowbreak 5 GB, irisfile.\allowbreak tar 65.\allowbreak 1 MB, ImageKeyFinal.\allowbreak xls, InputFiles.\allowbreak tar); scripts at https:/\allowbreak /\allowbreak github.\allowbreak com/\allowbreak AnthonyShiverMicrobes/\allowbreak fitness\_\allowbreak score \\
\addlinespace[5pt]
28 & \textbf{Thompson 2020 fatty acid and alcohol}\newline Thompson MG, Incha MR, Pearson AN, Schmidt M, Sharpless WA, Eiben CB, et al. Fatty Acid and Alcohol Metabolism in Pseudomonas putida: Functional Analysis Using Random Barcode Transposon Sequencing. Applied and Environmental Microbiology 2020\newline \href{https://doi.org/10.1128/aem.01665-20}{doi:10.\allowbreak 1128/\allowbreak aem.\allowbreak 01665-\allowbreak 20} & transposon insertion\newline {\scriptsize endpoint or steady state} & https:/\allowbreak /\allowbreak fit.\allowbreak genomics.\allowbreak lbl.\allowbreak gov/\allowbreak cgi-\allowbreak bin/\allowbreak org.\allowbreak cgi?orgId=Putida\newline {\scriptsize (unconfirmed)} \\
\addlinespace[5pt]
29 & \textbf{Borchert 2023 lignin tolerance}\newline Borchert AJ, Bleem A, Beckham GT. RB-TnSeq identifies genetic targets for improved tolerance of Pseudomonas putida towards compounds relevant to lignin conversion. Metabolic Engineering 2023\newline \href{https://doi.org/10.1016/j.ymben.2023.04.007}{doi:10.\allowbreak 1016/\allowbreak j.\allowbreak ymben.\allowbreak 2023.\allowbreak 04.\allowbreak 007} & transposon insertion\newline {\scriptsize endpoint or steady state} & https:/\allowbreak /\allowbreak www.\allowbreak ncbi.\allowbreak nlm.\allowbreak nih.\allowbreak gov/\allowbreak bioproject/\allowbreak PRJNA856070 \\
\addlinespace[5pt]
30 & \textbf{Schastnaya 2021}\newline Schastnaya E, Raguz Nakic Z, Gruber CH, Doubleday PF, Krishnan A, Johns NI, Park J, Wang HH, Sauer U. Extensive regulation of enzyme activity by phosphorylation in Escherichia coli. Nature Communications 2021\newline \href{https://doi.org/10.1038/s41467-021-25988-4}{doi:10.\allowbreak 1038/\allowbreak s41467-\allowbreak 021-\allowbreak 25988-\allowbreak 4} & mixed\newline {\scriptsize endpoint or steady state} & PXD027243 (PRIDE); MSV000087795 (MassIVE metabolomics) \\
\addlinespace[5pt]
31 & \textbf{Cui 2018}\newline Cui L, Vigouroux A, Rousset F, Varet H, Khanna V, Bikard D. A CRISPRi screen in E. coli reveals sequence-specific toxicity of dCas9. Nature Communications 2018\newline \href{https://doi.org/10.1038/s41467-018-04209-5}{doi:10.\allowbreak 1038/\allowbreak s41467-\allowbreak 018-\allowbreak 04209-\allowbreak 5} & CRISPRi\newline {\scriptsize endpoint or steady state} & Supplementary Data 4 of https:/\allowbreak /\allowbreak www.\allowbreak nature.\allowbreak com/\allowbreak articles/\allowbreak s41467-\allowbreak 018-\allowbreak 04209-\allowbreak 5 (no repository accession)\newline {\scriptsize (unconfirmed)} \\
\addlinespace[5pt]
32 & \textbf{Schmidt 2016}\newline Schmidt A, Kochanowski K, Vedelaar S, Ahrne E, Volkmer B, Callipo L, Knoops K, Bauer M, Aebersold R, Heinemann M. The quantitative and condition-dependent Escherichia coli proteome. Nature Biotechnology 2016\newline \href{https://doi.org/10.1038/nbt.3418}{doi:10.\allowbreak 1038/\allowbreak nbt.\allowbreak 3418} & environment-only\newline {\scriptsize endpoint or steady state} & PXD000498 (ProteomeXchange/\allowbreak PRIDE, DOI 10.\allowbreak 6019/\allowbreak PXD000498) \\
\addlinespace[5pt]
33 & \textbf{Butland 2008}\newline Butland G, Babu M, Diaz-Mejia JJ, Bohdana F, Phanse S, Gold B, Yang W, Li J, Gagarinova AG, Pogoutse O, Mori H, Wanner BL, Lo H, Wasniewski J, Christopolous C, Ali M, Venn P, Safavi-Naini A, Sourour N, Caron S, Choi JY, Laigle L, Nazarians-Armavil A, Deshpande A, Joe S, Datsenko KA, Yamamoto N, Andrews BJ, Boone C, Ding H, Sheikh B, Moreno-Hagelsieb G, Greenblatt JF, Emili A. eSGA: E. coli synthetic genetic array analysis. Nature Methods 2008\newline \href{https://doi.org/10.1038/nmeth.1239}{doi:10.\allowbreak 1038/\allowbreak nmeth.\allowbreak 1239} & double deletion\newline {\scriptsize endpoint or steady state} & \newline {\scriptsize (unconfirmed)} \\
\addlinespace[5pt]
34 & \textbf{Mori 2021}\newline Mori M, Zhang Z, Banaei-Esfahani A, Lalanne JB, Okano H, Collins BC, Schmidt A, Schubert OT, Lee DS, Li GW, Aebersold R, Hwa T, Ludwig C. From coarse to fine: the absolute Escherichia coli proteome under diverse growth conditions. Molecular Systems Biology 2021\newline \href{https://doi.org/10.15252/msb.20209536}{doi:10.\allowbreak 15252/\allowbreak msb.\allowbreak 20209536} & environment-only\newline {\scriptsize endpoint or steady state} & PASS01421 (SWATHAtlas spectral library); Datasets EV1-\allowbreak EV12 \\
\addlinespace[5pt]
35 & \textbf{Ishii 2007}\newline Ishii N, Nakahigashi K, Baba T, Robert M, Soga T, Kanai A, Hirasawa T, Naba M, Hirai K, Hoque A, Ho PY, Kakazu Y, Sugawara K, Igarashi S, Harada S, Masuda T, Sugiyama N, Togashi T, Hasegawa M, Takai Y, Yugi K, Arakawa K, Iwata N, Toya Y, Nakayama Y, Nishioka T, Shimizu K, Mori H, Tomita M. Multiple high-throughput analyses monitor the response of E. coli to perturbations. Science 2007\newline \href{https://doi.org/10.1126/science.1132067}{doi:10.\allowbreak 1126/\allowbreak science.\allowbreak 1132067} & mixed\newline {\scriptsize endpoint or steady state} & Science supporting online material; no repository accession found\newline {\scriptsize (unconfirmed)} \\
\addlinespace[5pt]
36 & \textbf{Campos 2018}\newline Campos M, Govers SK, Irnov I, Dobihal GS, Cornet F, Jacobs-Wagner C. Genomewide phenotypic analysis of growth, cell morphogenesis, and cell cycle events in Escherichia coli. Molecular Systems Biology 2018\newline \href{https://doi.org/10.15252/msb.20177573}{doi:10.\allowbreak 15252/\allowbreak msb.\allowbreak 20177573} & gene deletion\newline {\scriptsize endpoint or steady state} & Dataset EV2 (corrected and normalized scores) with the article at https:/\allowbreak /\allowbreak www.\allowbreak embopress.\allowbreak org/\allowbreak doi/\allowbreak full/\allowbreak 10.\allowbreak 15252/\allowbreak msb.\allowbreak 20177573 \\
\addlinespace[5pt]
37 & \textbf{Typas 2008}\newline Typas A, Nichols RJ, Siegele DA, Shales M, Collins SR, Lim B, Braberg H, Yamamoto N, Takeuchi R, Wanner BL, Mori H, Weissman JS, Krogan NJ, Gross CA. High-throughput, quantitative analyses of genetic interactions in E. coli. Nature Methods 2008\newline \href{https://doi.org/10.1038/nmeth.1240}{doi:10.\allowbreak 1038/\allowbreak nmeth.\allowbreak 1240} & double deletion\newline {\scriptsize endpoint or steady state} & Supplementary Tables 2A and 2B via PMC2700713; author manuscript at https:/\allowbreak /\allowbreak escholarship.\allowbreak org/\allowbreak content/\allowbreak qt0b4526ns/\allowbreak qt0b4526ns.\allowbreak pdf\newline {\scriptsize (unconfirmed)} \\
\addlinespace[5pt]
38 & \textbf{Girgis 2009}\newline Girgis HS, Hottes AK, Tavazoie S. Genetic architecture of intrinsic antibiotic susceptibility. PLoS ONE 2009\newline \href{https://doi.org/10.1371/journal.pone.0005629}{doi:10.\allowbreak 1371/\allowbreak journal.\allowbreak pone.\allowbreak 0005629} & transposon insertion\newline {\scriptsize endpoint or steady state} & PLoS ONE Supporting Information Datasets S1 to S5 at https:/\allowbreak /\allowbreak journals.\allowbreak plos.\allowbreak org/\allowbreak plosone/\allowbreak article?id=10.\allowbreak 1371/\allowbreak journal.\allowbreak pone.\allowbreak 0005629 (no repository accession)\newline {\scriptsize (unconfirmed)} \\
\addlinespace[5pt]
39 & \textbf{Banerjee 2025}\newline Banerjee D, Menasalvas J, Chen Y, Gin JW, Baidoo EEK, Petzold CJ, Eng T, Mukhopadhyay A. Addressing genome scale design tradeoffs in Pseudomonas putida for bioconversion of an aromatic carbon source. npj Systems Biology and Applications 2025\newline \href{https://doi.org/10.1038/s41540-024-00480-z}{doi:10.\allowbreak 1038/\allowbreak s41540-\allowbreak 024-\allowbreak 00480-\allowbreak z} & promoter variant\newline {\scriptsize sampled series} & https:/\allowbreak /\allowbreak proteomecentral.\allowbreak proteomexchange.\allowbreak org/\allowbreak cgi/\allowbreak GetDataset?ID=PXD050285 \\
\addlinespace[5pt]
40 & \textbf{Mohiuddin 2022 promoter library}\newline Mohiuddin SG, Massahi A, Orman MA. High-Throughput Screening of a Promoter Library Reveals New Persister Mechanisms in Escherichia coli. Microbiology Spectrum 2022;10(1):e0225321\newline \href{https://doi.org/10.1128/spectrum.02253-21}{doi:10.\allowbreak 1128/\allowbreak spectrum.\allowbreak 02253-\allowbreak 21} & promoter-GFP reporter\newline {\scriptsize 9 reads across a 5 hour exposure} & Supplemental File 2 spectrum02253-\allowbreak 21\_\allowbreak supp\_\allowbreak 2\_\allowbreak seq5.\allowbreak xlsx, sheets Raw Data and Fold Change \\
\addlinespace[5pt]
41 & \textbf{Teteneva 2024 lake water}\newline Teteneva N, Sanches-Medeiros A, Sourjik V. Genome-wide screen of genetic determinants that govern Escherichia coli growth and persistence in lake water. The ISME Journal 2024;18(1):wrae096\newline \href{https://doi.org/10.1093/ismejo/wrae096}{doi:10.\allowbreak 1093/\allowbreak ismejo/\allowbreak wrae096} & transposon insertion\newline {\scriptsize endpoint or steady state} & Supplementary Table S4 (xlsx, gene fitness); NCBI BioProject PRJNA1043681 (lake water) and PRJNA1073534 (Tn5 library annotation); code at github.\allowbreak com/\allowbreak NataliyaTeteneva/\allowbreak Tn5\_\allowbreak library\_\allowbreak analysis \\
\addlinespace[5pt]
42 & \textbf{Rachwalski 2024 mobile CRISPRi}\newline Rachwalski K, Tu MM, Madden SJ, French S, Hansen DM, Brown ED. A mobile CRISPRi collection enables genetic interaction studies for the essential genes of Escherichia coli. Cell Reports Methods 2024;4:100693\newline \href{https://doi.org/10.1016/j.crmeth.2023.100693}{doi:10.\allowbreak 1016/\allowbreak j.\allowbreak crmeth.\allowbreak 2023.\allowbreak 100693} & essential-gene knockdown crossed with a cataloged deletion\newline {\scriptsize endpoint or steady state} & Zenodo 10.\allowbreak 5281/\allowbreak zenodo.\allowbreak 10214517 (1.\allowbreak 4 GB and 364 MB image archives plus analysis code); Tables S1-\allowbreak S4 as xlsx; vector pFD152 is Addgene 125546 \\
\addlinespace[5pt]
43 & \textbf{Fang 2025 FFA CRISPRi-FACS}\newline Fang L, Hao X, Fan J, Liu X, Chen Y, Wang L, Huang X, Song H, Cao Y. Genome-scale CRISPRi screen identifies pcnB repression conferring improved physiology for overproduction of free fatty acids in Escherichia coli. Nature Communications 2025\newline \href{https://doi.org/10.1038/s41467-025-58368-3}{doi:10.\allowbreak 1038/\allowbreak s41467-\allowbreak 025-\allowbreak 58368-\allowbreak 3} & CRISPRi\newline {\scriptsize endpoint or steady state} & https:/\allowbreak /\allowbreak www.\allowbreak nature.\allowbreak com/\allowbreak articles/\allowbreak s41467-\allowbreak 025-\allowbreak 58368-\allowbreak 3\newline {\scriptsize (unconfirmed)} \\
\addlinespace[5pt]
44 & \textbf{Babu 2014}\newline Babu M, Arnold R, Bundalovic-Torma C, Gagarinova A, Wong KS, Kumar A, Stewart G, Samanfar B, Aoki H, Wagih O, Vlasblom J, Phanse S, Lad K, Yeou Hsiung Yu A, Graham C, Jin K, Brown E, Golshani A, Kim P, Moreno-Hagelsieb G, Greenblatt J, Houry WA, Parkinson J, Emili A. Quantitative genome-wide genetic interaction screens reveal global epistatic relationships of protein complexes in Escherichia coli. PLoS Genetics 2014\newline \href{https://doi.org/10.1371/journal.pgen.1004120}{doi:10.\allowbreak 1371/\allowbreak journal.\allowbreak pgen.\allowbreak 1004120} & double deletion\newline {\scriptsize endpoint or steady state} & Table S2 (XLS, high-\allowbreak confidence epistatic gene pairs) at https:/\allowbreak /\allowbreak journals.\allowbreak plos.\allowbreak org/\allowbreak plosgenetics/\allowbreak article?id=10.\allowbreak 1371/\allowbreak journal.\allowbreak pgen.\allowbreak 1004120 \\
\addlinespace[5pt]
45 & \textbf{Gupta 2024 turnover}\newline Gupta M, Johnson ANT, Cruz ER, Costa EJ, Guest RL, Li SHJ, Hart EM, Nguyen T, Stadlmeier M, Bratton BP, Silhavy TJ, Wingreen NS, Gitai Z, Wuhr M. Global protein turnover quantification in Escherichia coli reveals cytoplasmic recycling under nitrogen limitation. Nature Communications 2024\newline \href{https://doi.org/10.1038/s41467-024-49920-8}{doi:10.\allowbreak 1038/\allowbreak s41467-\allowbreak 024-\allowbreak 49920-\allowbreak 8} & environment-only\newline {\scriptsize sampled series} & PXD042444 (ProteomeXchange/\allowbreak PRIDE); code at https:/\allowbreak /\allowbreak github.\allowbreak com/\allowbreak wuhrlab/\allowbreak ProteinTurnoverEcoli \\
\addlinespace[5pt]
46 & \textbf{MCF2Chem 2023}\newline Cai P, Liu S, Zhang D, Hu QN. MCF2Chem: A manually curated knowledge base of biosynthetic compound production. Biotechnology for Biofuels and Bioproducts 2023\newline \href{https://doi.org/10.1186/s13068-023-02419-8}{doi:10.\allowbreak 1186/\allowbreak s13068-\allowbreak 023-\allowbreak 02419-\allowbreak 8} & mixed\newline {\scriptsize endpoint or steady state} & https:/\allowbreak /\allowbreak mcf.\allowbreak lifesynther.\allowbreak com\newline {\scriptsize (unconfirmed)} \\
\addlinespace[5pt]
47 & \textbf{Choe 2019 genome-reduced ALE}\newline Choe D, Lee JH, Yoo M, Hwang S, Sung BH, Cho S, Palsson B, Kim SC, Cho BK. Adaptive laboratory evolution of a genome-reduced Escherichia coli. Nature Communications 2019;10:935\newline \href{https://doi.org/10.1038/s41467-019-08888-6}{doi:10.\allowbreak 1038/\allowbreak s41467-\allowbreak 019-\allowbreak 08888-\allowbreak 6} & adaptive laboratory evolution on a reduced genome\newline {\scriptsize populations resequenced at 20 timepoints over 62 days, 807 generations} & ENA PRJEB21199 (51 runs: 31 WGS, 6 RNA-\allowbreak seq, 8 Ribo-\allowbreak seq, 6 ChIP-\allowbreak seq); Supplementary Data 1-\allowbreak 7 as xlsx; MG1655 is NC\_\allowbreak 000913.\allowbreak 3, MS56 has no accession \\
\addlinespace[5pt]
48 & \textbf{Niu 2019 pinene evolved}\newline Niu FX, Huang YB, Ji LN, Liu JZ. Genomic and transcriptional changes in response to pinene tolerance and overproduction in evolved Escherichia coli. Synthetic and Systems Biotechnology 2019;4:113-119\newline \href{https://doi.org/10.1016/j.synbio.2019.05.001}{doi:10.\allowbreak 1016/\allowbreak j.\allowbreak synbio.\allowbreak 2019.\allowbreak 05.\allowbreak 001} & CRISPR activation and interference on an evolved host\newline {\scriptsize endpoint or steady state} & Supplementary mmc1.\allowbreak docx holds the 374-\allowbreak variant table with b-\allowbreak numbers, all guide sequences, and both ratio tables; no sequence-\allowbreak data accession \\
\addlinespace[5pt]
49 & \textbf{Yang 2019 mevalonate from ethanol}\newline Yang J, Son JH, Kim H, Cho S, Na JG, Yeon YJ, Lee J. Mevalonate production from ethanol by direct conversion through acetyl-CoA using recombinant Pseudomonas putida, a novel biocatalyst for terpenoid production. Microbial Cell Factories 2019;18:168\newline \href{https://doi.org/10.1186/s12934-019-1213-y}{doi:10.\allowbreak 1186/\allowbreak s12934-\allowbreak 019-\allowbreak 1213-\allowbreak y} & stacked chromosomal deletion plus heterologous pathway\newline {\scriptsize metabolic time courses to 27 h in flask, 3 days in the fermenter} & GenBank gene ids for all five deleted loci and accessions for all six heterologous genes; Additional file 2 holds the codon-\allowbreak optimized sequences; no data deposit \\
\addlinespace[5pt]
50 & \textbf{Li 2021 mevalonate flux}\newline Li Y, Xian H, Xu Y, Zhu Y, Sun Z, Wang Q, Qi Q. Fine tuning the glycolytic flux ratio of EP-bifido pathway for mevalonate production by enhancing glucose-6-phosphate dehydrogenase and CRISPRi suppressing 6-phosphofructose kinase in Escherichia coli. Microbial Cell Factories 2021;20:32\newline \href{https://doi.org/10.1186/s12934-021-01526-1}{doi:10.\allowbreak 1186/\allowbreak s12934-\allowbreak 021-\allowbreak 01526-\allowbreak 1} & promoter replacement plus CRISPRi knockdown\newline {\scriptsize endpoint or steady state} & Additional file 2 (xlsx) holds the 198-\allowbreak reaction flux maps; promoters are iGEM Anderson registry parts; no data accession \\
\addlinespace[5pt]
\end{longtable}
\endgroup
\end{landscape}


%% notes-tex/database/database-expansion-bacteria/sections/queue.tex
%% [[experiments.database.expansion-bacteria]]
%% https://github.com/Mjvolk3/torchcell/tree/main/notes-tex/database/database-expansion-bacteria/sections/queue.tex
%% SOURCE: queue-rows table from
%% experiments/database/scripts/build_bacteria_discovery_queue.py;
%% surrounding prose is hand-authored, and every number in it is printed by that
%% script's "queue-rows:" lines

\section{The discovery queue, less the fifty\secstatus{todo}}\label{sec:queue}

The 300-publication discovery queue of Sec.~\ref{sec:sweep} is a list of leads rather
than datasets: no row carries a verified instance count, because the source asserts
none. Table~\ref{tab:bqueuerows} prints every queue row except the 19 that are already
among the fifty recommended builds, matched on DOI, which leaves 281 publications, 191
for \org{E. coli} and 90 for \org{P. putida}. Ten of those already sit in the ranked
reserve of Table~\ref{tab:bfinal} and are marked; the rest have no curated row.

The source sorts each row into one of three portfolio classes, defined in its own field
guide. \term{Joint} is ``Direct production engineering combined with scalable CRISPR,
screen, evolution, or omics measurements.'' \term{ME breadth} is ``Independent
product-oriented metabolic-engineering campaign.'' \term{Scale anchor} is ``Large
perturbation, phenotype, tolerance, or omics resource useful for model pretraining or
host-feature learning.''

\subsection{A second pass over the queue, and what the ordering rule kept\secstatus{todo}}
\label{sec:queue2}

The queue was read a second time with the selection order stated as scale times density
first, then overlap with a modality the store already serves for \org{E. coli}, then the
isoprenol axis for \org{P. putida}. Running it produced three rows, and the shape of the
result is more informative than the rows.

Most of the 271 uncurated rows are single-strain production campaigns, which this rule
ranks last by construction, so the pass reduced to the handful that claim a genome-scale
measurement. Six of those were checked against their deposits. Three release the matrix
and three release only the outcome, so the gate that decided this pass was not scale but
what the supplement contains. The three that passed are
\emph{Mohiuddin 2022} at final rank 40, a library of 1,930 promoter-GFP reporters read
nine times under four antibiotic conditions, 69,480 values; \emph{Teteneva 2024} at rank
41, barcoded transposon fitness for 3,691 genes in lake water, 66,162 values; and
\emph{Wang 2024} at rank 52, a transposon library under rifampicin on a dose-by-time
grid, 26,514 gene-by-condition rows with the negatives present. Each count was read off
the deposited file.

The three that failed are in Table~\ref{tab:bexcluded} and they fail the same way. A
screen of $10^{5}$ genomic fragments under furfural deposited 268 enriched genes and no
array data. A copper screen of the whole Keio collection at three doses scored growth
qualitatively and released 68 named mutants. A ciprofloxacin screen of the same
collection states that all data are in the supplement, and the supplement is a PDF
holding 48 mutants. In each case the abstract supports a measurement count in the tens
of thousands and the deposit supports a few hundred, which is the one failure mode that
a title-and-abstract sweep cannot detect and that ranking by the sweep's own score would
have rewarded.

Four more claim genome-scale measurement and were not resolved: a Keio screen against
benzoic acid and a screen on oleate, neither open access, an aggregated resequencing
corpus whose deposit resolves to analysis code rather than to the mutation table, and a
tunable-guide library whose only supplement is a PDF. They stay in the queue as leads and
none of them is counted anywhere, because a row is ranked on a number read off a file.

Two further things this pass settled. The \org{P. putida} half of the queue is now
exhausted against this rule: of its 82 uncurated rows not one is a genome-scale
measurement, and the single large \org{P. putida} dataset in the whole sweep was already
curated, carrying the sweep's \emph{lowest} priority score of the hundred. And adding two
rows inside the cut pushed two out, \emph{Hawkins 2020} and \emph{Royet 2025}, the second
of which is \org{P. putida}. Nothing pins them, because the pin of Sec.~\ref{sec:iso} is
for one product and neither row measures it. That displacement is the ordering working as
specified, and it is recorded here rather than absorbed silently.

%% GENERATED FILE -- do not hand-edit.
%% SOURCE: experiments/database/scripts/build_bacteria_discovery_queue.py
\begin{landscape}
\begingroup
\footnotesize
\setlength{\tabcolsep}{3pt}
\renewcommand{\arraystretch}{1.1}
\begin{longtable}{@{}r@{\hspace{4pt}} L{14mm} L{96mm} r L{18mm} L{45mm} L{14mm} L{38mm}@{}}
\caption[]{The discovery queue less the fifty recommended builds: every row of the 300-publication sweep except the 19 whose DOI matches one of rows 1--50 of Table~\ref{tab:bfinal} (9 \org{E. coli}, 10 \org{P. putida}), leaving 281: 191 \org{E. coli} and 90 \org{P. putida}. The match is on the DOI, lowercased and without its resolver prefix. No row carries a verified instance count, because the source asserts none. \emph{Rank} is the source's rank within its host, and a superscript \textbf{R} marks the 10 rows already curated in the ranked reserve below the cut. \emph{Class} is the source's portfolio class (Sec.~\ref{sec:queue}) and \emph{Isoprenol} its isoprenol flag. \emph{Product or readout} is the source's field as given: 101 rows carry one of its two generic labels and 180 a phrase taken from the title, 44 of those cut mid-word. Where the source gives no DOI (1 row), the DOI column links the source record instead.}
\label{tab:bqueuerows}\\
\toprule
\textbf{Rank} & \textbf{Host} & \textbf{Dataset or campaign} & \textbf{Year} & \textbf{Class} & \textbf{Product or readout} & \textbf{Isoprenol} & \textbf{DOI} \\
\midrule
\endfirsthead
\multicolumn{8}{@{}l}{\footnotesize\emph{Table~\ref{tab:bqueuerows}, continued}}\\
\toprule
\textbf{Rank} & \textbf{Host} & \textbf{Dataset or campaign} & \textbf{Year} & \textbf{Class} & \textbf{Product or readout} & \textbf{Isoprenol} & \textbf{DOI} \\
\midrule
\endhead
\bottomrule
\endfoot
3 & \org{E. coli} & \sourcetext{Hybrid metabolic engineering enables xylose-driven co-production of polyhydroxybutyrate and violacein in Escherichia coli.} & 2026 & ME breadth & \sourcetext{polyhydroxybutyrate and violacein} & -- & \href{https://doi.org/10.1186/s13036-026-00653-w}{10.\allowbreak 1186/\allowbreak s13036-\allowbreak 026-\allowbreak 00653-\allowbreak w} \\
\addlinespace[2pt]
4 & \org{E. coli} & \sourcetext{Demonstration and technoeconomic analysis of dodecanol production from acetate using metabolically engineered Escherichia coli.} & 2026 & ME breadth & \sourcetext{Demonstration and technoeconomic analysis of dodecanol} & -- & \href{https://doi.org/10.1016/j.ymben.2025.10.007}{10.\allowbreak 1016/\allowbreak j.\allowbreak ymben.\allowbreak 2025.\allowbreak 10.\allowbreak 007} \\
\addlinespace[2pt]
5 & \org{E. coli} & \sourcetext{Biosynthesis of Limonene Using Acetate as the Carbon Source in Metabolically Engineered Escherichia coli.} & 2026 & ME breadth & \sourcetext{Limonene Using Acetate as the Carbon Source in Metabolically Engineered Escherichia coli} & -- & \href{https://doi.org/10.1021/acs.jafc.5c14638}{10.\allowbreak 1021/\allowbreak acs.\allowbreak jafc.\allowbreak 5c14638} \\
\addlinespace[2pt]
6 & \org{E. coli} & \sourcetext{Genome-wide screening of the Keio collection identifies genetic determinants of Escherichia coli sensitivity and tolerance to benzoic acid.} & 2026 & Scale anchor & \sourcetext{Genome-scale/omics phenotype} & -- & \href{https://doi.org/10.1007/s11274-026-05211-6}{10.\allowbreak 1007/\allowbreak s11274-\allowbreak 026-\allowbreak 05211-\allowbreak 6} \\
\addlinespace[2pt]
8 & \org{E. coli} & \sourcetext{Long-term experimental evolution decouples size and production costs in Escherichia coli.} & 2022 & Joint & \sourcetext{Long-term experimental evolution decouples size and} & -- & \href{https://doi.org/10.1073/pnas.2200713119}{10.\allowbreak 1073/\allowbreak pnas.\allowbreak 2200713119} \\
\addlinespace[2pt]
9 & \org{E. coli} & \sourcetext{Combining directed evolution of pathway enzymes and dynamic pathway regulation using a quorum-sensing circuit to improve the production of 4-hydroxyphenylacetic acid in Escherichia coli.} & 2019 & Joint & \sourcetext{4-hydroxyphenylacetic acid} & -- & \href{https://doi.org/10.1186/s13068-019-1438-3}{10.\allowbreak 1186/\allowbreak s13068-\allowbreak 019-\allowbreak 1438-\allowbreak 3} \\
\addlinespace[2pt]
10 & \org{E. coli} & \sourcetext{ATP and NADPH engineering of Escherichia coli to improve the production of 4-hydroxyphenylacetic acid using CRISPRi.} & 2021 & Joint & \sourcetext{4-hydroxyphenylacetic acid using CRISPRi} & -- & \href{https://doi.org/10.1186/s13068-021-01954-6}{10.\allowbreak 1186/\allowbreak s13068-\allowbreak 021-\allowbreak 01954-\allowbreak 6} \\
\addlinespace[2pt]
11 & \org{E. coli} & \sourcetext{Effect of adaptive laboratory evolution of engineered Escherichia coli in acetate on the biosynthesis of succinic acid from glucose in two-stage cultivation.} & 2024 & Joint & \sourcetext{succinic acid from glucose in two-stage cultivation} & -- & \href{https://doi.org/10.1186/s40643-024-00749-5}{10.\allowbreak 1186/\allowbreak s40643-\allowbreak 024-\allowbreak 00749-\allowbreak 5} \\
\addlinespace[2pt]
12 & \org{E. coli} & \sourcetext{Fine-Tuning of hemB Using CRISPRi for Increasing 5-Aminolevulinic Acid Production in Escherichia coli.} & 2019 & Joint & \sourcetext{g of hemB Using CRISPRi for Increasing 5-Aminolevulinic Acid} & -- & \href{https://doi.org/10.3389/fmicb.2019.01731}{10.\allowbreak 3389/\allowbreak fmicb.\allowbreak 2019.\allowbreak 01731} \\
\addlinespace[2pt]
13 & \org{E. coli} & \sourcetext{Adaptive laboratory evolution and shuffling of Escherichia coli to enhance its tolerance and production of astaxanthin.} & 2022 & Joint & \sourcetext{astaxanthin} & -- & \href{https://doi.org/10.1186/s13068-022-02118-w}{10.\allowbreak 1186/\allowbreak s13068-\allowbreak 022-\allowbreak 02118-\allowbreak w} \\
\addlinespace[2pt]
14 & \org{E. coli} & \sourcetext{Directed Evolution of Branched-Chain $\alpha$-Keto Acid Decarboxylase for 3-Hydroxypropionic Acid Production in Escherichia coli via Oxaloacetate.} & 2025 & Joint & \sourcetext{Chain $\alpha$-Keto Acid Decarboxylase for 3-Hydroxypropionic Acid} & -- & \href{https://doi.org/10.1021/acssynbio.5c00267}{10.\allowbreak 1021/\allowbreak acssynbio.\allowbreak 5c00267} \\
\addlinespace[2pt]
15 & \org{E. coli} & \sourcetext{Biosensor-driven evolution and metabolic engineering of an Escherichia coli cell factory for L-tryptophan production.} & 2026 & Joint & \sourcetext{neering of an Escherichia coli cell factory for L-tryptophan} & -- & \href{https://doi.org/10.1186/s12934-026-03028-4}{10.\allowbreak 1186/\allowbreak s12934-\allowbreak 026-\allowbreak 03028-\allowbreak 4} \\
\addlinespace[2pt]
16 & \org{E. coli} & \sourcetext{Enhancing the heat tolerance of homolactic Escherichia coli through reverse metabolic engineering.} & 2026 & Joint & \sourcetext{Genome-scale/omics phenotype} & -- & \href{https://doi.org/10.1007/s00253-025-13703-y}{10.\allowbreak 1007/\allowbreak s00253-\allowbreak 025-\allowbreak 13703-\allowbreak y} \\
\addlinespace[2pt]
17 & \org{E. coli} & \sourcetext{Metabolic engineering of Escherichia coli BW25113 for the production of Vitamin K2 based on CRISPR/Cas9 mediated gene knockout and metabolic pathway modification.} & 2026 & Joint & \sourcetext{Vitamin K2 based on CRISPR/Cas9 mediated gene knockout and metabolic pathway modification} & -- & \href{https://doi.org/10.1186/s13036-025-00614-9}{10.\allowbreak 1186/\allowbreak s13036-\allowbreak 025-\allowbreak 00614-\allowbreak 9} \\
\addlinespace[2pt]
18 & \org{E. coli} & \sourcetext{A CRISPRi mediated self-inducible system for dynamic regulation of TCA cycle and improvement of itaconic acid production in Escherichia coli.} & 2022 & Joint & \sourcetext{mic regulation of TCA cycle and improvement of itaconic acid} & -- & \href{https://doi.org/10.1016/j.synbio.2022.05.008}{10.\allowbreak 1016/\allowbreak j.\allowbreak synbio.\allowbreak 2022.\allowbreak 05.\allowbreak 008} \\
\addlinespace[2pt]
19 & \org{E. coli} & \sourcetext{Integrated laboratory evolution and rational engineering of GalP/Glk-dependent Escherichia coli for higher yield and productivity of L-tryptophan biosynthesis.} & 2021 & Joint & \sourcetext{Genome-scale/omics phenotype} & -- & \href{https://doi.org/10.1016/j.mec.2021.e00167}{10.\allowbreak 1016/\allowbreak j.\allowbreak mec.\allowbreak 2021.\allowbreak e00167} \\
\addlinespace[2pt]
20 & \org{E. coli} & \sourcetext{Chromosomal evolution of Escherichia coli for the efficient production of lycopene.} & 2013 & Joint & \sourcetext{lycopene} & adjacent & \href{https://doi.org/10.1186/1472-6750-13-6}{10.\allowbreak 1186/\allowbreak 1472-\allowbreak 6750-\allowbreak 13-\allowbreak 6} \\
\addlinespace[2pt]
21 & \org{E. coli} & \sourcetext{Metabolic engineering of Escherichia coli BW25113 for the production of 5-Aminolevulinic Acid based on CRISPR/Cas9 mediated gene knockout and metabolic pathway modification.} & 2022 & Joint & \sourcetext{5-Aminolevulinic Acid based on CRISPR/Cas9 mediated gene knockout and metabolic pathway modification} & -- & \href{https://doi.org/10.1186/s13036-022-00307-7}{10.\allowbreak 1186/\allowbreak s13036-\allowbreak 022-\allowbreak 00307-\allowbreak 7} \\
\addlinespace[2pt]
23 & \org{E. coli} & \sourcetext{GREACE-assisted adaptive laboratory evolution in endpoint fermentation broth enhances lysine production by Escherichia coli.} & 2019 & Joint & \sourcetext{ory evolution in endpoint fermentation broth enhances lysine} & -- & \href{https://doi.org/10.1186/s12934-019-1153-6}{10.\allowbreak 1186/\allowbreak s12934-\allowbreak 019-\allowbreak 1153-\allowbreak 6} \\
\addlinespace[2pt]
24 & \org{E. coli} & \sourcetext{Deionococcus proteotlycius Genomic Library Exploration Enhances Oxidative Stress Resistance and Poly-3-hydroxybutyrate Production in Recombinant Escherichia coli.} & 2023 & Joint & \sourcetext{ances Oxidative Stress Resistance and Poly-3-hydroxybutyrate} & -- & \href{https://doi.org/10.3390/microorganisms11092135}{10.\allowbreak 3390/\allowbreak microorganisms11092135} \\
\addlinespace[2pt]
25 & \org{E. coli} & \sourcetext{Enhanced Production of Pterostilbene in Escherichia coli Through Directed Evolution and Host Strain Engineering.} & 2021 & Joint & \sourcetext{Pterostilbene} & -- & \href{https://doi.org/10.3389/fmicb.2021.710405}{10.\allowbreak 3389/\allowbreak fmicb.\allowbreak 2021.\allowbreak 710405} \\
\addlinespace[2pt]
26 & \org{E. coli} & \sourcetext{Novel Mode Engineering for $\beta$-Alanine Production in Escherichia coli with the Guide of Adaptive Laboratory Evolution.} & 2021 & Joint & \sourcetext{Novel Mode Engineering for $\beta$-Alanine} & -- & \href{https://doi.org/10.3390/microorganisms9030600}{10.\allowbreak 3390/\allowbreak microorganisms9030600} \\
\addlinespace[2pt]
27 & \org{E. coli} & \sourcetext{Enhancing the Production of Pinene in Escherichia coli by Using a Combination of Shotgun, Product-Tolerance and I-SceI Cleavage Systems.} & 2022 & Joint & \sourcetext{Pinene} & adjacent & \href{https://doi.org/10.3390/biology11101484}{10.\allowbreak 3390/\allowbreak biology11101484} \\
\addlinespace[2pt]
29 & \org{E. coli} & \sourcetext{Directed evolution of alditol oxidase for the production of optically pure D-glycerate from glycerol in the engineered Escherichia coli.} & 2021 & Joint & \sourcetext{optically pure D-glycerate from glycerol in the engineered Escherichia coli} & -- & \href{https://doi.org/10.1093/jimb/kuab041}{10.\allowbreak 1093/\allowbreak jimb/\allowbreak kuab041} \\
\addlinespace[2pt]
30 & \org{E. coli} & \sourcetext{Inhibitor tolerance and bioethanol fermentability of levoglucosan-utilizing Escherichia coli were enhanced by overexpression of stress-responsive gene ycfR: The proteomics-guided metabolic engineering.} & 2021 & Joint & \sourcetext{Genome-scale/omics phenotype} & -- & \href{https://doi.org/10.1016/j.synbio.2021.11.003}{10.\allowbreak 1016/\allowbreak j.\allowbreak synbio.\allowbreak 2021.\allowbreak 11.\allowbreak 003} \\
\addlinespace[2pt]
31 & \org{E. coli} & \sourcetext{Improving phloroglucinol tolerance and production in Escherichia coli by GroESL overexpression.} & 2017 & Joint & \sourcetext{Improving phloroglucinol tolerance and} & -- & \href{https://doi.org/10.1186/s12934-017-0839-x}{10.\allowbreak 1186/\allowbreak s12934-\allowbreak 017-\allowbreak 0839-\allowbreak x} \\
\addlinespace[2pt]
32 & \org{E. coli} & \sourcetext{Restoration of biofuel production levels and increased tolerance under ionic liquid stress is enabled by a mutation in the essential Escherichia coli gene cydC.} & 2018 & Joint & \sourcetext{Restoration of biofuel} & -- & \href{https://doi.org/10.1186/s12934-018-1006-8}{10.\allowbreak 1186/\allowbreak s12934-\allowbreak 018-\allowbreak 1006-\allowbreak 8} \\
\addlinespace[2pt]
33 & \org{E. coli} & \sourcetext{Ferulic acid production by metabolically engineered Escherichia coli.} & 2021 & ME breadth & \sourcetext{Ferulic acid} & -- & \href{https://doi.org/10.1186/s40643-021-00423-0}{10.\allowbreak 1186/\allowbreak s40643-\allowbreak 021-\allowbreak 00423-\allowbreak 0} \\
\addlinespace[2pt]
34 & \org{E. coli} & \sourcetext{In silico identification of gene targets to enhance C12 fatty acid production in Escherichia coli.} & 2025 & ME breadth & \sourcetext{ico identification of gene targets to enhance C12 fatty acid} & -- & \href{https://doi.org/10.1007/s00253-025-13501-6}{10.\allowbreak 1007/\allowbreak s00253-\allowbreak 025-\allowbreak 13501-\allowbreak 6} \\
\addlinespace[2pt]
35 & \org{E. coli} & \sourcetext{Metabolic engineering enables Escherichia coli to grow on 1,3-propanediol.} & 2026 & ME breadth & \sourcetext{Metabolic-engineering phenotype} & -- & \href{https://doi.org/10.1016/j.synbio.2025.11.009}{10.\allowbreak 1016/\allowbreak j.\allowbreak synbio.\allowbreak 2025.\allowbreak 11.\allowbreak 009} \\
\addlinespace[2pt]
36 & \org{E. coli} & \sourcetext{Decoupling sucrose utilization from oxygen-responsive regulation for high-efficiency L-lactic acid production in Escherichia coli.} & 2025 & ME breadth & \sourcetext{ygen-responsive regulation for high-efficiency L-lactic acid} & -- & \href{https://doi.org/10.1186/s13068-025-02700-y}{10.\allowbreak 1186/\allowbreak s13068-\allowbreak 025-\allowbreak 02700-\allowbreak y} \\
\addlinespace[2pt]
37 & \org{E. coli} & \sourcetext{Heterologous Production of 2,2'-Dihydroxy Derivatives of Astaxanthin and Adonirubin in Escherichia coli and Evaluation of Their Antioxidant Activity.} & 2026 & ME breadth & \sourcetext{2} & -- & \href{https://doi.org/10.3390/antiox15030327}{10.\allowbreak 3390/\allowbreak antiox15030327} \\
\addlinespace[2pt]
38 & \org{E. coli} & \sourcetext{Automated high-throughput screening combined with model-assisted analysis enables quantitative identification of process drivers and trade-offs during extracellular Fab production in Escherichia coli.} & 2026 & ME breadth & \sourcetext{n of process drivers and trade-offs during extracellular Fab} & -- & \href{https://doi.org/10.1007/s00449-026-03379-7}{10.\allowbreak 1007/\allowbreak s00449-\allowbreak 026-\allowbreak 03379-\allowbreak 7} \\
\addlinespace[2pt]
39 & \org{E. coli} & \sourcetext{Enhancing genome-scale metabolic models with kinetic data: resolving growth and citramalate production trade-offs in Escherichia coli.} & 2025 & ME breadth & \sourcetext{resolving growth and citramalate} & -- & \href{https://doi.org/10.1093/bioadv/vbaf166}{10.\allowbreak 1093/\allowbreak bioadv/\allowbreak vbaf166} \\
\addlinespace[2pt]
40 & \org{E. coli} & \sourcetext{System metabolic engineering of Escherichia coli W for the production of 2-ketoisovalerate using unconventional feedstock.} & 2023 & ME breadth & \sourcetext{2-ketoisovalerate using unconventional feedstock} & -- & \href{https://doi.org/10.3389/fbioe.2023.1176445}{10.\allowbreak 3389/\allowbreak fbioe.\allowbreak 2023.\allowbreak 1176445} \\
\addlinespace[2pt]
41 & \org{E. coli} & \sourcetext{Evaluation of Various Escherichia coli Strains for Enhanced Lycopene Production.} & 2023 & ME breadth & \sourcetext{on of Various Escherichia coli Strains for Enhanced Lycopene} & adjacent & \href{https://doi.org/10.4014/jmb.2302.02003}{10.\allowbreak 4014/\allowbreak jmb.\allowbreak 2302.\allowbreak 02003} \\
\addlinespace[2pt]
42 & \org{E. coli} & \sourcetext{Rational engineering of Escherichia coli strain for stable and enhanced biosynthesis of pinene.} & 2024 & ME breadth & \sourcetext{pinene} & adjacent & \href{https://doi.org/10.3389/fmicb.2024.1527113}{10.\allowbreak 3389/\allowbreak fmicb.\allowbreak 2024.\allowbreak 1527113} \\
\addlinespace[2pt]
43 & \org{E. coli} & \sourcetext{Revalorisation of brewer's spent grain for biotechnological production of hydrogen with Escherichia coli.} & 2024 & ME breadth & \sourcetext{hydrogen} & -- & \href{https://doi.org/10.3389/fbioe.2024.1473704}{10.\allowbreak 3389/\allowbreak fbioe.\allowbreak 2024.\allowbreak 1473704} \\
\addlinespace[2pt]
44 & \org{E. coli} & \sourcetext{Engineering Escherichia coli for Isobutanol Production from Xylose or Glucose-Xylose Mixture.} & 2023 & ME breadth & \sourcetext{Engineering Escherichia coli for Isobutanol} & -- & \href{https://doi.org/10.3390/microorganisms11102573}{10.\allowbreak 3390/\allowbreak microorganisms11102573} \\
\addlinespace[2pt]
45 & \org{E. coli} & \sourcetext{Metabolic engineering of Escherichia coli for efficient production of L-alanyl-L-glutamine.} & 2020 & ME breadth & \sourcetext{L-alanyl-L-glutamine} & -- & \href{https://doi.org/10.1186/s12934-020-01369-2}{10.\allowbreak 1186/\allowbreak s12934-\allowbreak 020-\allowbreak 01369-\allowbreak 2} \\
\addlinespace[2pt]
46 & \org{E. coli} & \sourcetext{Efficient production of hydroxytyrosol by directed evolution of HpaB in Escherichia coli.} & 2023 & Joint & \sourcetext{hydroxytyrosol by directed evolution of HpaB} & -- & \href{https://doi.org/10.1016/j.bbrc.2023.04.024}{10.\allowbreak 1016/\allowbreak j.\allowbreak bbrc.\allowbreak 2023.\allowbreak 04.\allowbreak 024} \\
\addlinespace[2pt]
47 & \org{E. coli} & \sourcetext{Enhanced production of functional CRISPR-AsCas12a protein in Escherichia coli.} & 2025 & Joint & \sourcetext{functional CRISPR-AsCas12a protein} & -- & \href{https://doi.org/10.1016/j.pep.2025.106722}{10.\allowbreak 1016/\allowbreak j.\allowbreak pep.\allowbreak 2025.\allowbreak 106722} \\
\addlinespace[2pt]
48 & \org{E. coli} & \sourcetext{Increased production of L-serine in Escherichia coli through Adaptive Laboratory Evolution.} & 2017 & Joint & \sourcetext{L-serine} & -- & \href{https://doi.org/10.1016/j.ymben.2016.11.008}{10.\allowbreak 1016/\allowbreak j.\allowbreak ymben.\allowbreak 2016.\allowbreak 11.\allowbreak 008} \\
\addlinespace[2pt]
49 & \org{E. coli} & \sourcetext{Production of Rainbow Colorants by Metabolically Engineered Escherichia coli.} & 2021 & ME breadth & \sourcetext{Rainbow Colorants by Metabolically Engineered Escherichia coli} & -- & \href{https://doi.org/10.1002/advs.202100743}{10.\allowbreak 1002/\allowbreak advs.\allowbreak 202100743} \\
\addlinespace[2pt]
50 & \org{E. coli} & \sourcetext{Enhanced production of methylated resveratrol in engineered Escherichia coli and its anticancer activity.} & 2025 & ME breadth & \sourcetext{methylated resveratrol} & -- & \href{https://doi.org/10.1007/s13205-025-04449-5}{10.\allowbreak 1007/\allowbreak s13205-\allowbreak 025-\allowbreak 04449-\allowbreak 5} \\
\addlinespace[2pt]
51 & \org{E. coli} & \sourcetext{Multi-level metabolic engineering of Escherichia coli for high-titer biosynthesis of 2'-fucosyllactose and 3-fucosyllactose.} & 2022 & ME breadth & \sourcetext{2'-fucosyllactose and 3-fucosyllactose} & -- & \href{https://doi.org/10.1111/1751-7915.14152}{10.\allowbreak 1111/\allowbreak 1751-\allowbreak 7915.\allowbreak 14152} \\
\addlinespace[2pt]
52 & \org{E. coli} & \sourcetext{Metabolic engineering of Escherichia coli for optimized biosynthesis of nicotinamide mononucleotide, a noncanonical redox cofactor.} & 2020 & ME breadth & \sourcetext{nicotinamide mononucleotide} & -- & \href{https://doi.org/10.1186/s12934-020-01415-z}{10.\allowbreak 1186/\allowbreak s12934-\allowbreak 020-\allowbreak 01415-\allowbreak z} \\
\addlinespace[2pt]
53 & \org{E. coli} & \sourcetext{Synergetic engineering of Escherichia coli for efficient production of l-tyrosine.} & 2023 & ME breadth & \sourcetext{l-tyrosine} & -- & \href{https://doi.org/10.1016/j.synbio.2023.10.005}{10.\allowbreak 1016/\allowbreak j.\allowbreak synbio.\allowbreak 2023.\allowbreak 10.\allowbreak 005} \\
\addlinespace[2pt]
54 & \org{E. coli} & \sourcetext{Metabolic engineering and pathway construction for O-acetyl-L-homoserine production in Escherichia coli.} & 2023 & ME breadth & \sourcetext{gineering and pathway construction for O-acetyl-L-homoserine} & -- & \href{https://doi.org/10.1007/s13205-023-03564-5}{10.\allowbreak 1007/\allowbreak s13205-\allowbreak 023-\allowbreak 03564-\allowbreak 5} \\
\addlinespace[2pt]
55 & \org{E. coli} & \sourcetext{Metabolic engineering of Escherichia\,coli for biosynthesis of $\beta$-nicotinamide mononucleotide from nicotinamide.} & 2021 & ME breadth & \sourcetext{$\beta$-nicotinamide mononucleotide from nicotinamide} & -- & \href{https://doi.org/10.1111/1751-7915.13901}{10.\allowbreak 1111/\allowbreak 1751-\allowbreak 7915.\allowbreak 13901} \\
\addlinespace[2pt]
56 & \org{E. coli} & \sourcetext{Systems metabolic engineering of Escherichia coli for hyper-production of 5-aminolevulinic acid.} & 2023 & ME breadth & \sourcetext{5-aminolevulinic acid} & -- & \href{https://doi.org/10.1186/s13068-023-02280-9}{10.\allowbreak 1186/\allowbreak s13068-\allowbreak 023-\allowbreak 02280-\allowbreak 9} \\
\addlinespace[2pt]
57 & \org{E. coli} & \sourcetext{Engineering Escherichia coli membrane phospholipid head distribution improves tolerance and production of biorenewables.} & 2017 & Joint & \sourcetext{biorenewables} & -- & \href{https://doi.org/10.1016/j.ymben.2017.08.006}{10.\allowbreak 1016/\allowbreak j.\allowbreak ymben.\allowbreak 2017.\allowbreak 08.\allowbreak 006} \\
\addlinespace[2pt]
58 & \org{E. coli} & \sourcetext{Activating the d-Tagatose Production Capacity of Escherichia coli with Structural Insights into C4 Epimerase Specificity.} & 2025 & ME breadth & \sourcetext{Activating the d-Tagatose} & -- & \href{https://doi.org/10.1021/acs.jafc.4c12842}{10.\allowbreak 1021/\allowbreak acs.\allowbreak jafc.\allowbreak 4c12842} \\
\addlinespace[2pt]
59 & \org{E. coli} & \sourcetext{Enhanced cyanophycin production by Escherichia coli overexpressing the heterologous cphA gene from a deep sea metagenomic library.} & 2017 & Joint & \sourcetext{Enhanced cyanophycin} & -- & \href{https://doi.org/10.1016/j.jbiosc.2016.08.008}{10.\allowbreak 1016/\allowbreak j.\allowbreak jbiosc.\allowbreak 2016.\allowbreak 08.\allowbreak 008} \\
\addlinespace[2pt]
60 & \org{E. coli} & \sourcetext{A robust platform streamlining aromatic noncanonical amino acid biosynthesis and genetic code expansion in Escherichia coli.} & 2025 & ME breadth & \sourcetext{Metabolic-engineering phenotype} & -- & \href{https://doi.org/10.1038/s41467-025-63679-6}{10.\allowbreak 1038/\allowbreak s41467-\allowbreak 025-\allowbreak 63679-\allowbreak 6} \\
\addlinespace[2pt]
61 & \org{E. coli} & \sourcetext{Efficient production of 2'-fucosyllactose from fructose through metabolically engineered recombinant Escherichia coli.} & 2024 & ME breadth & \sourcetext{2'-fucosyllactose from fructose through metabolically engineered recombinant Escherichia coli} & -- & \href{https://doi.org/10.1186/s12934-024-02312-5}{10.\allowbreak 1186/\allowbreak s12934-\allowbreak 024-\allowbreak 02312-\allowbreak 5} \\
\addlinespace[2pt]
62 & \org{E. coli} & \sourcetext{Efficient Production of 3'-Sialyllactose Using Escherichia coli.} & 2024 & ME breadth & \sourcetext{3'-Sialyllactose} & -- & \href{https://doi.org/10.1021/acs.jafc.4c08703}{10.\allowbreak 1021/\allowbreak acs.\allowbreak jafc.\allowbreak 4c08703} \\
\addlinespace[2pt]
63 & \org{E. coli} & \sourcetext{Metabolic engineering strategies of de novo pathway for enhancing 2'-fucosyllactose synthesis in Escherichia coli.} & 2022 & ME breadth & \sourcetext{Metabolic-engineering phenotype} & -- & \href{https://doi.org/10.1111/1751-7915.13977}{10.\allowbreak 1111/\allowbreak 1751-\allowbreak 7915.\allowbreak 13977} \\
\addlinespace[2pt]
64 & \org{E. coli} & \sourcetext{De Novo Biosynthesis of Antidepressant Psilocybin in Escherichia coli.} & 2025 & ME breadth & \sourcetext{Antidepressant Psilocybin} & -- & \href{https://doi.org/10.1111/1751-7915.70135}{10.\allowbreak 1111/\allowbreak 1751-\allowbreak 7915.\allowbreak 70135} \\
\addlinespace[2pt]
65 & \org{E. coli} & \sourcetext{Awakening the natural capability of psicose production in Escherichia coli.} & 2023 & ME breadth & \sourcetext{Awakening the natural capability of psicose} & -- & \href{https://doi.org/10.1038/s41538-023-00231-0}{10.\allowbreak 1038/\allowbreak s41538-\allowbreak 023-\allowbreak 00231-\allowbreak 0} \\
\addlinespace[2pt]
66 & \org{E. coli} & \sourcetext{Improved production of class I lanthipeptides in Escherichia coli.} & 2023 & ME breadth & \sourcetext{class I lanthipeptides} & -- & \href{https://doi.org/10.1039/d2sc06597e}{10.\allowbreak 1039/\allowbreak d2sc06597e} \\
\addlinespace[2pt]
67 & \org{E. coli} & \sourcetext{Metabolic engineering of Escherichia coli for efficient production of L-5-hydroxytryptophan from glucose.} & 2022 & ME breadth & \sourcetext{L-5-hydroxytryptophan from glucose} & -- & \href{https://doi.org/10.1186/s12934-022-01920-3}{10.\allowbreak 1186/\allowbreak s12934-\allowbreak 022-\allowbreak 01920-\allowbreak 3} \\
\addlinespace[2pt]
68 & \org{E. coli} & \sourcetext{Metabolic engineering of Escherichia coli for high-yield dopamine production via optimized fermentation strategies.} & 2025 & ME breadth & \sourcetext{olic engineering of Escherichia coli for high-yield dopamine} & -- & \href{https://doi.org/10.1128/aem.00159-25}{10.\allowbreak 1128/\allowbreak aem.\allowbreak 00159-\allowbreak 25} \\
\addlinespace[2pt]
69 & \org{E. coli} & \sourcetext{Engineering Electron Transfer Flux between Cytochrome P450 Enzyme and P450 Reductase to Enhance Serotonin Production in Escherichia Coli.} & 2025 & ME breadth & \sourcetext{tochrome P450 Enzyme and P450 Reductase to Enhance Serotonin} & -- & \href{https://doi.org/10.1002/advs.202414859}{10.\allowbreak 1002/\allowbreak advs.\allowbreak 202414859} \\
\addlinespace[2pt]
70 & \org{E. coli} & \sourcetext{A production platform for disulfide-bonded peptides in the periplasm of Escherichia coli.} & 2024 & ME breadth & \sourcetext{Metabolic-engineering phenotype} & -- & \href{https://doi.org/10.1186/s12934-024-02446-6}{10.\allowbreak 1186/\allowbreak s12934-\allowbreak 024-\allowbreak 02446-\allowbreak 6} \\
\addlinespace[2pt]
71 & \org{E. coli} & \sourcetext{Increased distribution of carbon metabolic flux during de novo cytidine biosynthesis via attenuation of the acetic acid metabolism pathway in Escherichia coli.} & 2025 & ME breadth & \sourcetext{Metabolic-engineering phenotype} & -- & \href{https://doi.org/10.1186/s12934-025-02657-5}{10.\allowbreak 1186/\allowbreak s12934-\allowbreak 025-\allowbreak 02657-\allowbreak 5} \\
\addlinespace[2pt]
72 & \org{E. coli} & \sourcetext{Efficient production of guanosine in Escherichia coli by combinatorial metabolic engineering.} & 2024 & ME breadth & \sourcetext{guanosine} & -- & \href{https://doi.org/10.1186/s12934-024-02452-8}{10.\allowbreak 1186/\allowbreak s12934-\allowbreak 024-\allowbreak 02452-\allowbreak 8} \\
\addlinespace[2pt]
73 & \org{E. coli} & \sourcetext{Step-by-step optimization of a heterologous pathway for de novo naringenin production in Escherichia coli.} & 2024 & ME breadth & \sourcetext{ptimization of a heterologous pathway for de novo naringenin} & -- & \href{https://doi.org/10.1007/s00253-024-13271-7}{10.\allowbreak 1007/\allowbreak s00253-\allowbreak 024-\allowbreak 13271-\allowbreak 7} \\
\addlinespace[2pt]
74 & \org{E. coli} & \sourcetext{Increased Mevalonate Production Using Engineered Citrate Synthase and Phosphofructokinase Variants of Escherichia coli.} & 2025 & ME breadth & \sourcetext{Increased Mevalonate} & adjacent & \href{https://doi.org/10.1002/bit.28902}{10.\allowbreak 1002/\allowbreak bit.\allowbreak 28902} \\
\addlinespace[2pt]
75 & \org{E. coli} & \sourcetext{Engineering an Escherichia coli strain for production of long single-stranded DNA.} & 2024 & ME breadth & \sourcetext{long single-stranded DNA} & -- & \href{https://doi.org/10.1093/nar/gkae189}{10.\allowbreak 1093/\allowbreak nar/\allowbreak gkae189} \\
\addlinespace[2pt]
76 & \org{E. coli} & \sourcetext{Creating Polyploid Escherichia Coli and Its Application in Efficient L-Threonine Production.} & 2023 & ME breadth & \sourcetext{scherichia Coli and Its Application in Efficient L-Threonine} & -- & \href{https://doi.org/10.1002/advs.202302417}{10.\allowbreak 1002/\allowbreak advs.\allowbreak 202302417} \\
\addlinespace[2pt]
77 & \org{E. coli} & \sourcetext{Synergetic Fermentation of Glucose and Glycerol for High-Yield N-Acetylglucosamine Production in Escherichia coli.} & 2022 & ME breadth & \sourcetext{n of Glucose and Glycerol for High-Yield N-Acetylglucosamine} & -- & \href{https://doi.org/10.3390/ijms23020773}{10.\allowbreak 3390/\allowbreak ijms23020773} \\
\addlinespace[2pt]
78 & \org{E. coli} & \sourcetext{High-level de novo biosynthesis of glycosylated zeaxanthin and astaxanthin in Escherichia coli.} & 2021 & ME breadth & \sourcetext{glycosylated zeaxanthin and astaxanthin} & -- & \href{https://doi.org/10.1186/s40643-021-00415-0}{10.\allowbreak 1186/\allowbreak s40643-\allowbreak 021-\allowbreak 00415-\allowbreak 0} \\
\addlinespace[2pt]
79 & \org{E. coli} & \sourcetext{Engineering Escherichia coli for malate production by integrating modular pathway characterization with CRISPRi-guided multiplexed metabolic tuning.} & 2018 & Joint & \sourcetext{Engineering Escherichia coli for malate} & -- & \href{https://doi.org/10.1002/bit.26486}{10.\allowbreak 1002/\allowbreak bit.\allowbreak 26486} \\
\addlinespace[2pt]
80 & \org{E. coli} & \sourcetext{Metabolic engineering of Escherichia coli for squalene overproduction.} & 2025 & ME breadth & \sourcetext{Metabolic-engineering phenotype} & -- & \href{https://doi.org/10.1016/j.synbio.2025.06.003}{10.\allowbreak 1016/\allowbreak j.\allowbreak synbio.\allowbreak 2025.\allowbreak 06.\allowbreak 003} \\
\addlinespace[2pt]
81 & \org{E. coli} & \sourcetext{Acetol biosynthesis enables NADPH balance during nitrogen limitation in engineered Escherichia coli.} & 2025 & ME breadth & \sourcetext{Metabolic-engineering phenotype} & -- & \href{https://doi.org/10.1186/s12934-025-02687-z}{10.\allowbreak 1186/\allowbreak s12934-\allowbreak 025-\allowbreak 02687-\allowbreak z} \\
\addlinespace[2pt]
82 & \org{E. coli} & \sourcetext{pSIG plasmids, MoClo-compatible vectors for efficient production of chimeric double-stranded RNAs in Escherichia coli HT115 (DE3) strain.} & 2025 & ME breadth & \sourcetext{chimeric double-stranded RNAs} & -- & \href{https://doi.org/10.1186/s13007-025-01413-5}{10.\allowbreak 1186/\allowbreak s13007-\allowbreak 025-\allowbreak 01413-\allowbreak 5} \\
\addlinespace[2pt]
83 & \org{E. coli} & \sourcetext{A novel strategy for L-arginine production in engineered Escherichia coli.} & 2023 & ME breadth & \sourcetext{A novel strategy for L-arginine} & -- & \href{https://doi.org/10.1186/s12934-023-02145-8}{10.\allowbreak 1186/\allowbreak s12934-\allowbreak 023-\allowbreak 02145-\allowbreak 8} \\
\addlinespace[2pt]
84 & \org{E. coli} & \sourcetext{Efficient production of myo-inositol in Escherichia coli through metabolic engineering.} & 2020 & ME breadth & \sourcetext{myo-inositol} & -- & \href{https://doi.org/10.1186/s12934-020-01366-5}{10.\allowbreak 1186/\allowbreak s12934-\allowbreak 020-\allowbreak 01366-\allowbreak 5} \\
\addlinespace[2pt]
85 & \org{E. coli} & \sourcetext{Exploiting Mixed Waste Office Paper Containing Lignocellulosic Fibers for Alternatively Producing High-Value Succinic Acid by Metabolically Engineered Escherichia coli KJ122.} & 2025 & ME breadth & \sourcetext{Metabolic-engineering phenotype} & -- & \href{https://doi.org/10.3390/ijms26030982}{10.\allowbreak 3390/\allowbreak ijms26030982} \\
\addlinespace[2pt]
86 & \org{E. coli} & \sourcetext{Reconstitution of Methionine Cycle With ATP Regeneration for Whole-Cell Catalysis of Creatine Production in Engineered Escherichia coli.} & 2025 & ME breadth & \sourcetext{e With ATP Regeneration for Whole-Cell Catalysis of Creatine} & -- & \href{https://doi.org/10.1111/1751-7915.70145}{10.\allowbreak 1111/\allowbreak 1751-\allowbreak 7915.\allowbreak 70145} \\
\addlinespace[2pt]
87 & \org{E. coli} & \sourcetext{Efficient production of hydroxysalidroside in Escherichia coli via enhanced glycosylation and semi-rational design of UGT85A1.} & 2025 & ME breadth & \sourcetext{hydroxysalidroside} & -- & \href{https://doi.org/10.1016/j.synbio.2025.03.002}{10.\allowbreak 1016/\allowbreak j.\allowbreak synbio.\allowbreak 2025.\allowbreak 03.\allowbreak 002} \\
\addlinespace[2pt]
88 & \org{E. coli} & \sourcetext{Enhanced Outer Membrane Vesicle Production in Escherichia coli: From Metabolic Network Model to Designed Strain Lipidomic Profile.} & 2025 & ME breadth & \sourcetext{Enhanced Outer Membrane Vesicle} & -- & \href{https://doi.org/10.3390/ijms26146714}{10.\allowbreak 3390/\allowbreak ijms26146714} \\
\addlinespace[2pt]
89 & \org{E. coli} & \sourcetext{The knowledge driven DBTL cycle provides mechanistic insights while optimising dopamine production in Escherichia coli.} & 2025 & ME breadth & \sourcetext{ycle provides mechanistic insights while optimising dopamine} & -- & \href{https://doi.org/10.1186/s12934-025-02729-6}{10.\allowbreak 1186/\allowbreak s12934-\allowbreak 025-\allowbreak 02729-\allowbreak 6} \\
\addlinespace[2pt]
90 & \org{E. coli} & \sourcetext{De novo biosynthesis of 4,6-dihydroxycoumarin in Escherichia coli.} & 2025 & ME breadth & \sourcetext{4} & -- & \href{https://doi.org/10.1039/d4gc05694a}{10.\allowbreak 1039/\allowbreak d4gc05694a} \\
\addlinespace[2pt]
91 & \org{E. coli} & \sourcetext{Systems engineering of Escherichia coli for high-level glutarate production from glucose.} & 2024 & ME breadth & \sourcetext{ems engineering of Escherichia coli for high-level glutarate} & -- & \href{https://doi.org/10.1038/s41467-024-45448-z}{10.\allowbreak 1038/\allowbreak s41467-\allowbreak 024-\allowbreak 45448-\allowbreak z} \\
\addlinespace[2pt]
92 & \org{E. coli} & \sourcetext{Metabolic Engineering of Escherichia coli for Ectoine Production With a Fermentation Strategy of Supplementing the Amino Donor.} & 2022 & ME breadth & \sourcetext{Metabolic Engineering of Escherichia coli for Ectoine} & -- & \href{https://doi.org/10.3389/fbioe.2022.824859}{10.\allowbreak 3389/\allowbreak fbioe.\allowbreak 2022.\allowbreak 824859} \\
\addlinespace[2pt]
93 & \org{E. coli} & \sourcetext{Systematic Metabolic Engineering and Model-Guided Optimization for High-Level Production of L-Theanine from Xylose in Escherichia coli.} & 2026 & ME breadth & \sourcetext{L-Theanine from Xylose} & -- & \href{https://doi.org/10.1002/advs.202521440}{10.\allowbreak 1002/\allowbreak advs.\allowbreak 202521440} \\
\addlinespace[2pt]
94 & \org{E. coli} & \sourcetext{Genome-scale Metabolic Modeling Guided Escherichia coli Engineering for De Novo Biosynthesis of Chrysanthemic Acid.} & 2026 & ME breadth & \sourcetext{Chrysanthemic Acid} & -- & \href{https://doi.org/10.1002/advs.202512736}{10.\allowbreak 1002/\allowbreak advs.\allowbreak 202512736} \\
\addlinespace[2pt]
95 & \org{E. coli} & \sourcetext{Engineering Escherichia coli for robust Co-utilization of glucose and xylose enables high-titer succinate production from lignocellulosic hydrolysates.} & 2026 & ME breadth & \sourcetext{ilization of glucose and xylose enables high-titer succinate} & -- & \href{https://doi.org/10.1016/j.synbio.2026.01.006}{10.\allowbreak 1016/\allowbreak j.\allowbreak synbio.\allowbreak 2026.\allowbreak 01.\allowbreak 006} \\
\addlinespace[2pt]
96 & \org{E. coli} & \sourcetext{Heterologous Biosynthesis of Prenylflavonoids in Escherichia coli Based on Fungus Screening of Prenyltransferases.} & 2025 & ME breadth & \sourcetext{Prenylflavonoids} & -- & \href{https://doi.org/10.1021/acsomega.4c05007}{10.\allowbreak 1021/\allowbreak acsomega.\allowbreak 4c05007} \\
\addlinespace[2pt]
97 & \org{E. coli} & \sourcetext{Monitoring of the Single-Cell Behavior of an Escherichia coli Reporter Strain Producing L-phenylalanine in a Scale-Down Bioreactor by Automated Real-Time Flow Cytometry.} & 2025 & ME breadth & \sourcetext{Metabolic-engineering phenotype} & -- & \href{https://doi.org/10.3390/biotech14030054}{10.\allowbreak 3390/\allowbreak biotech14030054} \\
\addlinespace[2pt]
98 & \org{E. coli} & \sourcetext{Multivariate modular metabolic engineering for enhanced L-methionine biosynthesis in Escherichia coli.} & 2023 & ME breadth & \sourcetext{Metabolic-engineering phenotype} & -- & \href{https://doi.org/10.1186/s13068-023-02347-7}{10.\allowbreak 1186/\allowbreak s13068-\allowbreak 023-\allowbreak 02347-\allowbreak 7} \\
\addlinespace[2pt]
99 & \org{E. coli} & \sourcetext{Metabolic engineering of Escherichia coli for high-level production of violaxanthin.} & 2023 & ME breadth & \sourcetext{violaxanthin} & -- & \href{https://doi.org/10.1186/s12934-023-02098-y}{10.\allowbreak 1186/\allowbreak s12934-\allowbreak 023-\allowbreak 02098-\allowbreak y} \\
\addlinespace[2pt]
100 & \org{E. coli} & \sourcetext{A Combinational Optimization Method for Efficient Production of Indigo by the Recombinant Escherichia coli with Expression of Monooxygenase and Malate Dehydrogenase.} & 2023 & ME breadth & \sourcetext{Indigo by the Recombinant Escherichia coli with Expression of Monooxygenase and Malate Dehydrogenase} & -- & \href{https://doi.org/10.3390/foods12030502}{10.\allowbreak 3390/\allowbreak foods12030502} \\
\addlinespace[2pt]
101 & \org{E. coli} & \sourcetext{Enhancing 5-aminolevulinic acid tolerance and production by engineering the antioxidant defense system of Escherichia coli.} & 2019 & Joint & \sourcetext{Enhancing 5-aminolevulinic acid tolerance and} & -- & \href{https://doi.org/10.1002/bit.26981}{10.\allowbreak 1002/\allowbreak bit.\allowbreak 26981} \\
\addlinespace[2pt]
102 & \org{E. coli} & \sourcetext{Metabolic engineering of endogenous MEP pathway for enhanced lycopene production in Escherichia coli.} & 2026 & ME breadth & \sourcetext{engineering of endogenous MEP pathway for enhanced lycopene} & adjacent & \href{https://doi.org/10.1016/j.synbio.2026.01.014}{10.\allowbreak 1016/\allowbreak j.\allowbreak synbio.\allowbreak 2026.\allowbreak 01.\allowbreak 014} \\
\addlinespace[2pt]
103 & \org{E. coli} & \sourcetext{Biosynthesis and Export of Membrane-Enveloped Selenium Nanoparticles by Escherichia coli.} & 2026 & ME breadth & \sourcetext{Metabolic-engineering phenotype} & -- & \href{https://doi.org/10.1021/acs.est.5c10008}{10.\allowbreak 1021/\allowbreak acs.\allowbreak est.\allowbreak 5c10008} \\
\addlinespace[2pt]
104 & \org{E. coli} & \sourcetext{Flavin Biosynthesis Enhances Extracellular Electron Transfer in Bioengineered Escherichia coli.} & 2026 & ME breadth & \sourcetext{Metabolic-engineering phenotype} & -- & \href{https://doi.org/10.1002/advs.202412230}{10.\allowbreak 1002/\allowbreak advs.\allowbreak 202412230} \\
\addlinespace[2pt]
105 & \org{E. coli} & \sourcetext{High-level production of vitamin K2 in Escherichia coli via modular molecular engineering.} & 2026 & ME breadth & \sourcetext{vitamin K2} & -- & \href{https://doi.org/10.1016/j.synbio.2025.11.001}{10.\allowbreak 1016/\allowbreak j.\allowbreak synbio.\allowbreak 2025.\allowbreak 11.\allowbreak 001} \\
\addlinespace[2pt]
106 & \org{E. coli} & \sourcetext{Combination of enzyme engineering and quorum sensing system for efficient de novo biosynthesis of $\beta$-arbutin in Escherichia coli.} & 2026 & ME breadth & \sourcetext{$\beta$-arbutin} & -- & \href{https://doi.org/10.1016/j.synbio.2025.08.010}{10.\allowbreak 1016/\allowbreak j.\allowbreak synbio.\allowbreak 2025.\allowbreak 08.\allowbreak 010} \\
\addlinespace[2pt]
107 & \org{E. coli} & \sourcetext{Enhanced vitamin B6 production in engineered Escherichia coli via restricted mix-carbon feeding and pressure-controlled fermentation.} & 2026 & ME breadth & \sourcetext{Enhanced vitamin B6} & -- & \href{https://doi.org/10.1016/j.synbio.2025.09.004}{10.\allowbreak 1016/\allowbreak j.\allowbreak synbio.\allowbreak 2025.\allowbreak 09.\allowbreak 004} \\
\addlinespace[2pt]
108 & \org{E. coli} & \sourcetext{De Novo Production of Xanthohumol by a Metabolically Engineered Escherichia coli.} & 2025 & ME breadth & \sourcetext{Xanthohumol by a Metabolically Engineered Escherichia coli} & -- & \href{https://doi.org/10.1021/acssynbio.5c00221}{10.\allowbreak 1021/\allowbreak acssynbio.\allowbreak 5c00221} \\
\addlinespace[2pt]
109 & \org{E. coli} & \sourcetext{Exploring single-cell biosynthetic noise and dynamics for enhanced betaxanthin production in Escherichia coli.} & 2025 & ME breadth & \sourcetext{ell biosynthetic noise and dynamics for enhanced betaxanthin} & -- & \href{https://doi.org/10.1038/s41467-025-67733-1}{10.\allowbreak 1038/\allowbreak s41467-\allowbreak 025-\allowbreak 67733-\allowbreak 1} \\
\addlinespace[2pt]
110 & \org{E. coli} & \sourcetext{Engineering Escherichia coli Biofilms for Curcumin Production.} & 2025 & ME breadth & \sourcetext{Engineering Escherichia coli Biofilms for Curcumin} & -- & \href{https://doi.org/10.3390/molecules30092031}{10.\allowbreak 3390/\allowbreak molecules30092031} \\
\addlinespace[2pt]
111 & \org{E. coli} & \sourcetext{Enhanced isoprenoid production in Escherichia coli cells harboring the archaeal mevalonate pathway.} & 2025 & ME breadth & \sourcetext{Enhanced isoprenoid} & adjacent & \href{https://doi.org/10.1016/j.bbrep.2025.102307}{10.\allowbreak 1016/\allowbreak j.\allowbreak bbrep.\allowbreak 2025.\allowbreak 102307} \\
\addlinespace[2pt]
112 & \org{E. coli} & \sourcetext{Synergistic Production of Lycopene and $\beta$-Alanine Through Engineered Redox Balancing in Escherichia coli.} & 2025 & ME breadth & \sourcetext{Lycopene and $\beta$-Alanine Through Engineered Redox Balancing} & adjacent & \href{https://doi.org/10.3390/ijms26146727}{10.\allowbreak 3390/\allowbreak ijms26146727} \\
\addlinespace[2pt]
113 & \org{E. coli} & \sourcetext{Scale up of fermentation of recombinant Escherichia coli for efficient production of spider drag silk protein MaSp1s and its dimers.} & 2025 & ME breadth & \sourcetext{spider drag silk protein MaSp1s and its dimers} & -- & \href{https://doi.org/10.1186/s12934-025-02734-9}{10.\allowbreak 1186/\allowbreak s12934-\allowbreak 025-\allowbreak 02734-\allowbreak 9} \\
\addlinespace[2pt]
114 & \org{E. coli} & \sourcetext{Efficient production of salicylic acid through CmeR-PcmeO biosensor-assisted multiplexing pathway optimization in Escherichia coli.} & 2025 & ME breadth & \sourcetext{salicylic acid through CmeR-PcmeO biosensor-assisted multiplexing pathway optimization} & -- & \href{https://doi.org/10.1186/s13068-025-02637-2}{10.\allowbreak 1186/\allowbreak s13068-\allowbreak 025-\allowbreak 02637-\allowbreak 2} \\
\addlinespace[2pt]
115 & \org{E. coli} & \sourcetext{Robust production of N-acetyl-glucosamine in engineered Escherichia coli from glycerol-glucose mixture.} & 2025 & ME breadth & \sourcetext{N-acetyl-glucosamine} & -- & \href{https://doi.org/10.1016/j.synbio.2025.05.003}{10.\allowbreak 1016/\allowbreak j.\allowbreak synbio.\allowbreak 2025.\allowbreak 05.\allowbreak 003} \\
\addlinespace[2pt]
116 & \org{E. coli} & \sourcetext{Production of minicell-like structures by Escherichia coli biosynthesizing cadmium fluorescent nanoparticles: a novel response to heavy metal exposure.} & 2025 & ME breadth & \sourcetext{minicell-like structures} & -- & \href{https://doi.org/10.1186/s12951-025-03188-2}{10.\allowbreak 1186/\allowbreak s12951-\allowbreak 025-\allowbreak 03188-\allowbreak 2} \\
\addlinespace[2pt]
117 & \org{E. coli} & \sourcetext{Efficient production of phenyllactic acid in Escherichia coli via metabolic engineering and fermentation optimization strategies.} & 2024 & ME breadth & \sourcetext{phenyllactic acid} & -- & \href{https://doi.org/10.3389/fmicb.2024.1457628}{10.\allowbreak 3389/\allowbreak fmicb.\allowbreak 2024.\allowbreak 1457628} \\
\addlinespace[2pt]
118 & \org{E. coli} & \sourcetext{Optimizing Archaeal Lipid Biosynthesis in Escherichia coli.} & 2024 & ME breadth & \sourcetext{Metabolic-engineering phenotype} & -- & \href{https://doi.org/10.1021/acssynbio.4c00235}{10.\allowbreak 1021/\allowbreak acssynbio.\allowbreak 4c00235} \\
\addlinespace[2pt]
119 & \org{E. coli} & \sourcetext{Systematic Engineering of Escherichia coli for Efficient Production of Cytidine 5'-Monophosphate.} & 2024 & ME breadth & \sourcetext{Cytidine 5'-Monophosphate} & -- & \href{https://doi.org/10.1021/acsomega.3c07713}{10.\allowbreak 1021/\allowbreak acsomega.\allowbreak 3c07713} \\
\addlinespace[2pt]
120 & \org{E. coli} & \sourcetext{Budding and explosive membrane vesicle production by hypervesiculating Escherichia coli strain $\Delta$rodZ.} & 2024 & ME breadth & \sourcetext{Budding and explosive membrane vesicle} & -- & \href{https://doi.org/10.3389/fmicb.2024.1400434}{10.\allowbreak 3389/\allowbreak fmicb.\allowbreak 2024.\allowbreak 1400434} \\
\addlinespace[2pt]
121 & \org{E. coli} & \sourcetext{Reprogramming Metabolic Flux in Escherichia Coli to Enhance Chondroitin Production.} & 2024 & ME breadth & \sourcetext{ng Metabolic Flux} & -- & \href{https://doi.org/10.1002/advs.202307351}{10.\allowbreak 1002/\allowbreak advs.\allowbreak 202307351} \\
\addlinespace[2pt]
122 & \org{E. coli} & \sourcetext{Systems metabolic engineering for hydroxytyrosol production in Escherichia coli.} & 2026 & ME breadth & \sourcetext{Systems metabolic engineering for hydroxytyrosol} & -- & \href{https://doi.org/10.1128/aem.02455-25}{10.\allowbreak 1128/\allowbreak aem.\allowbreak 02455-\allowbreak 25} \\
\addlinespace[2pt]
123 & \org{E. coli} & \sourcetext{Rare-codon-tuned fluorescent biosensor for high-throughput selection of Escherichia coli strain capable of L-arginine overproduction.} & 2026 & ME breadth & \sourcetext{Metabolic-engineering phenotype} & -- & \href{https://doi.org/10.1186/s12866-026-05154-w}{10.\allowbreak 1186/\allowbreak s12866-\allowbreak 026-\allowbreak 05154-\allowbreak w} \\
\addlinespace[2pt]
124 & \org{E. coli} & \sourcetext{Systems metabolic engineering of dual pathways with stage-specific activation for high-level 5-aminolevulinic acid production in Escherichia coli.} & 2026 & ME breadth & \sourcetext{age-specific activation for high-level 5-aminolevulinic acid} & -- & \href{https://doi.org/10.1016/j.synbio.2025.09.007}{10.\allowbreak 1016/\allowbreak j.\allowbreak synbio.\allowbreak 2025.\allowbreak 09.\allowbreak 007} \\
\addlinespace[2pt]
125 & \org{E. coli} & \sourcetext{Bioconversion of lignin and food waste into vanillin and biohydrogen by Escherichia coli MC3: a novel dual-functional bacteria.} & 2026 & ME breadth & \sourcetext{lignin and food waste into vanillin and biohydrogen} & -- & \href{https://doi.org/10.1186/s13068-025-02728-0}{10.\allowbreak 1186/\allowbreak s13068-\allowbreak 025-\allowbreak 02728-\allowbreak 0} \\
\addlinespace[2pt]
126 & \org{E. coli} & \sourcetext{Systematic engineering of Escherichia coli for biosynthesis of 3-hydroxypropionic acid from glucose and malonate.} & 2026 & ME breadth & \sourcetext{3-hydroxypropionic acid from glucose and malonate} & -- & \href{https://doi.org/10.1186/s12896-026-01140-2}{10.\allowbreak 1186/\allowbreak s12896-\allowbreak 026-\allowbreak 01140-\allowbreak 2} \\
\addlinespace[2pt]
127 & \org{E. coli} & \sourcetext{De novo synthesis of L-2-aminobutyric acid in Escherichia coli based on multi-layered metabolic engineering strategies.} & 2026 & ME breadth & \sourcetext{Metabolic-engineering phenotype} & -- & \href{https://doi.org/10.1016/j.synbio.2026.01.004}{10.\allowbreak 1016/\allowbreak j.\allowbreak synbio.\allowbreak 2026.\allowbreak 01.\allowbreak 004} \\
\addlinespace[2pt]
128 & \org{E. coli} & \sourcetext{Integrated amplification of NADPH-regenerating modules enhances cytidine biosynthesis in Escherichia coli.} & 2026 & ME breadth & \sourcetext{Metabolic-engineering phenotype} & -- & \href{https://doi.org/10.1016/j.synbio.2025.12.009}{10.\allowbreak 1016/\allowbreak j.\allowbreak synbio.\allowbreak 2025.\allowbreak 12.\allowbreak 009} \\
\addlinespace[2pt]
129 & \org{E. coli} & \sourcetext{De novo biosynthesis of putrescine from glucose using engineered Escherichia coli in a two-step process.} & 2026 & ME breadth & \sourcetext{putrescine from glucose} & -- & \href{https://doi.org/10.1016/j.synbio.2026.02.008}{10.\allowbreak 1016/\allowbreak j.\allowbreak synbio.\allowbreak 2026.\allowbreak 02.\allowbreak 008} \\
\addlinespace[2pt]
130 & \org{E. coli} & \sourcetext{Combining computer-aided enzyme design and chromosomal integration for plasmid-free biosynthesis of 1,5-pentanediol in Escherichia coli.} & 2026 & ME breadth & \sourcetext{1} & -- & \href{https://doi.org/10.1186/s12934-026-02929-8}{10.\allowbreak 1186/\allowbreak s12934-\allowbreak 026-\allowbreak 02929-\allowbreak 8} \\
\addlinespace[2pt]
131 & \org{E. coli} & \sourcetext{Heterologous Production of Torularhodin, the Monocyclic Carotenoid with a Terminal Carboxyl Group, in Escherichia coli.} & 2026 & ME breadth & \sourcetext{Torularhodin} & adjacent & \href{https://doi.org/10.3390/biotech15010003}{10.\allowbreak 3390/\allowbreak biotech15010003} \\
\addlinespace[2pt]
132 & \org{E. coli} & \sourcetext{Insights into the Regulation of Indigo Production in an Engineered Escherichia coli Strain via Overexpression of Specific Transporter Genes and Proteomic Analyzes.} & 2026 & ME breadth & \sourcetext{Insights into the Regulation of Indigo} & -- & \href{https://doi.org/10.3390/foods15081385}{10.\allowbreak 3390/\allowbreak foods15081385} \\
\addlinespace[2pt]
133 & \org{E. coli} & \sourcetext{High-titer 2-phenylethanol production in Escherichia coli via decarboxylation-step optimization and a vegetable-oil overlay for biocompatible in situ recovery.} & 2026 & ME breadth & \sourcetext{High-titer 2-phenylethanol} & -- & \href{https://doi.org/10.1186/s13036-026-00667-4}{10.\allowbreak 1186/\allowbreak s13036-\allowbreak 026-\allowbreak 00667-\allowbreak 4} \\
\addlinespace[2pt]
134 & \org{E. coli} & \sourcetext{Development of a transcription factor-based biosensor strain for reporting $\alpha$-terpineol production via the alcohol-dependent hemiterpene pathway in Escherichia coli.} & 2026 & ME breadth & \sourcetext{tion factor-based biosensor strain for reporting $\alpha$-terpineol} & adjacent & \href{https://doi.org/10.1039/d5cb00310e}{10.\allowbreak 1039/\allowbreak d5cb00310e} \\
\addlinespace[2pt]
135 & \org{E. coli} & \sourcetext{Metabolic engineering of Escherichia coli for de novo production of lauryl glucoside.} & 2025 & ME breadth & \sourcetext{lauryl glucoside} & -- & \href{https://doi.org/10.1016/j.bidere.2025.100045}{10.\allowbreak 1016/\allowbreak j.\allowbreak bidere.\allowbreak 2025.\allowbreak 100045} \\
\addlinespace[2pt]
136 & \org{E. coli} & \sourcetext{Phase-driven rewiring in~Escherichia coli enhances coenzyme Q10 biosynthesis via temporal and energetic coordination.} & 2025 & ME breadth & \sourcetext{Metabolic-engineering phenotype} & -- & \href{https://doi.org/10.1007/s00253-025-13619-7}{10.\allowbreak 1007/\allowbreak s00253-\allowbreak 025-\allowbreak 13619-\allowbreak 7} \\
\addlinespace[2pt]
137 & \org{E. coli} & \sourcetext{Engineering a Robust Escherichia coli W Platform for Scalable Production of Flavonoid-O-Glucosides.} & 2025 & ME breadth & \sourcetext{Flavonoid-O-Glucosides} & -- & \href{https://doi.org/10.1111/1751-7915.70226}{10.\allowbreak 1111/\allowbreak 1751-\allowbreak 7915.\allowbreak 70226} \\
\addlinespace[2pt]
138 & \org{E. coli} & \sourcetext{Enhancing Escherichia coli production of material proteins using circular mRNAs.} & 2025 & ME breadth & \sourcetext{material proteins using circular mRNAs} & -- & \href{https://doi.org/10.1128/aem.01579-25}{10.\allowbreak 1128/\allowbreak aem.\allowbreak 01579-\allowbreak 25} \\
\addlinespace[2pt]
139 & \org{E. coli} & \sourcetext{Development of Sucrose-Utilizing Escherichia coli Nissle 1917 for Efficient Heparosan Biosynthesis.} & 2025 & ME breadth & \sourcetext{Metabolic-engineering phenotype} & -- & \href{https://doi.org/10.3390/metabo15060410}{10.\allowbreak 3390/\allowbreak metabo15060410} \\
\addlinespace[2pt]
140 & \org{E. coli} & \sourcetext{Biosynthesis of Two Types of Exogenous Antigenic Polysaccharides in a Single Escherichia coli Chassis Cell.} & 2025 & ME breadth & \sourcetext{Two Types of Exogenous Antigenic Polysaccharides in a Single Escherichia coli Chassis Cell} & -- & \href{https://doi.org/10.3390/life15060858}{10.\allowbreak 3390/\allowbreak life15060858} \\
\addlinespace[2pt]
141 & \org{E. coli} & \sourcetext{Manipulating Intracellular Oxidative Conditions to Enhance Porphyrin Production in Escherichia coli.} & 2025 & ME breadth & \sourcetext{ting Intracellular Oxidative Conditions to Enhance Porphyrin} & -- & \href{https://doi.org/10.3390/bioengineering12010083}{10.\allowbreak 3390/\allowbreak bioengineering12010083} \\
\addlinespace[2pt]
142 & \org{E. coli} & \sourcetext{Improving sustainable isopropanol production in engineered Escherichia coli W via oxygen limitation.} & 2025 & ME breadth & \sourcetext{Improving sustainable isopropanol} & -- & \href{https://doi.org/10.1186/s12934-025-02720-1}{10.\allowbreak 1186/\allowbreak s12934-\allowbreak 025-\allowbreak 02720-\allowbreak 1} \\
\addlinespace[2pt]
143 & \org{E. coli} & \sourcetext{Extracellular peptide production in Escherichia coli by inducible downregulation of lipoprotein Lpp via MicL sRNA.} & 2025 & ME breadth & \sourcetext{Extracellular peptide} & -- & \href{https://doi.org/10.1007/s00253-025-13524-z}{10.\allowbreak 1007/\allowbreak s00253-\allowbreak 025-\allowbreak 13524-\allowbreak z} \\
\addlinespace[2pt]
144 & \org{E. coli} & \sourcetext{Exchange of the L-cysteine exporter after in-vivo metabolic control analysis improved the L-cysteine production process with engineered Escherichia coli.} & 2025 & ME breadth & \sourcetext{r in-vivo metabolic control analysis improved the L-cysteine} & -- & \href{https://doi.org/10.1186/s12934-025-02715-y}{10.\allowbreak 1186/\allowbreak s12934-\allowbreak 025-\allowbreak 02715-\allowbreak y} \\
\addlinespace[2pt]
145 & \org{E. coli} & \sourcetext{Modular pathway engineering for enhanced production of para-aminobenzoic acid and 4-amino-phenylalanine in Escherichia coli via glucose/xylose co-utilization.} & 2025 & ME breadth & \sourcetext{para-aminobenzoic acid and 4-amino-phenylalanine} & -- & \href{https://doi.org/10.1128/aem.02468-24}{10.\allowbreak 1128/\allowbreak aem.\allowbreak 02468-\allowbreak 24} \\
\addlinespace[2pt]
146 & \org{E. coli} & \sourcetext{Enhancing production and assessing IgE reactivity of dog allergen Can f 6 in Pichia pastoris and Escherichia coli.} & 2025 & ME breadth & \sourcetext{Enhancing} & -- & \href{https://doi.org/10.1007/s00253-025-13465-7}{10.\allowbreak 1007/\allowbreak s00253-\allowbreak 025-\allowbreak 13465-\allowbreak 7} \\
\addlinespace[2pt]
147 & \org{E. coli} & \sourcetext{Mutations in the riboflavin biosynthesis pathway confer resistance to furazolidone and abolish the synergistic interaction between furazolidone and vancomycin in Escherichia coli.} & 2025 & ME breadth & \sourcetext{Metabolic-engineering phenotype} & -- & \href{https://doi.org/10.1099/mgen.0.001356}{10.\allowbreak 1099/\allowbreak mgen.\allowbreak 0.\allowbreak 001356} \\
\addlinespace[2pt]
148 & \org{E. coli} & \sourcetext{Glucoselipid Biosurfactant Biosynthesis Operon of Rouxiella badensis DSM 100043T: Screening, Identification, and Heterologous Expression in Escherichia coli.} & 2025 & ME breadth & \sourcetext{Metabolic-engineering phenotype} & -- & \href{https://doi.org/10.3390/microorganisms13071664}{10.\allowbreak 3390/\allowbreak microorganisms13071664} \\
\addlinespace[2pt]
149 & \org{E. coli} & \sourcetext{Deciphering Transcription-Translation-Folding (TX-TL-FD) for Enhancing Cutinase Production in T7 System and Genetic Chaperone-Equipped Escherichia coli Strains.} & 2025 & ME breadth & \sourcetext{iption-Translation-Folding (TX-TL-FD) for Enhancing Cutinase} & -- & \href{https://doi.org/10.1021/acssynbio.5c00245}{10.\allowbreak 1021/\allowbreak acssynbio.\allowbreak 5c00245} \\
\addlinespace[2pt]
150 & \org{E. coli} & \sourcetext{Multivariate modular metabolic engineering and medium optimization for vitamin B12 production by Escherichia coli.} & 2024 & ME breadth & \sourcetext{etabolic engineering and medium optimization for vitamin B12} & -- & \href{https://doi.org/10.1016/j.synbio.2024.03.017}{10.\allowbreak 1016/\allowbreak j.\allowbreak synbio.\allowbreak 2024.\allowbreak 03.\allowbreak 017} \\
\addlinespace[2pt]
151 & \org{E. coli} & \sourcetext{Engineering a non-oxidative glycolysis pathway in escherichia coli for high-level citramalate production.} & 2024 & ME breadth & \sourcetext{lysis pathway} & -- & \href{https://doi.org/10.1186/s12934-024-02505-y}{10.\allowbreak 1186/\allowbreak s12934-\allowbreak 024-\allowbreak 02505-\allowbreak y} \\
\addlinespace[2pt]
153\,\textsuperscript{\textbf{R}} & \org{E. coli} & \sourcetext{Genomic library screens for genes involved in n-butanol tolerance in Escherichia coli.} & 2011 & Scale anchor & \sourcetext{Genome-scale/omics phenotype} & -- & \href{https://doi.org/10.1371/journal.pone.0017678}{10.\allowbreak 1371/\allowbreak journal.\allowbreak pone.\allowbreak 0017678} \\
\addlinespace[2pt]
154\,\textsuperscript{\textbf{R}} & \org{E. coli} & \sourcetext{Morphological and Transcriptional Responses to CRISPRi Knockdown of Essential Genes in Escherichia coli.} & 2021 & Scale anchor & \sourcetext{Genome-scale/omics phenotype} & -- & \href{https://doi.org/10.1128/mbio.02561-21}{10.\allowbreak 1128/\allowbreak mbio.\allowbreak 02561-\allowbreak 21} \\
\addlinespace[2pt]
156 & \org{E. coli} & \sourcetext{CRISPRi-mediated tunable control of gene expression level with engineered single-guide RNA in Escherichia coli.} & 2023 & Scale anchor & \sourcetext{Genome-scale/omics phenotype} & -- & \href{https://doi.org/10.1093/nar/gkad234}{10.\allowbreak 1093/\allowbreak nar/\allowbreak gkad234} \\
\addlinespace[2pt]
157 & \org{E. coli} & \sourcetext{Improving the Methanol Tolerance of an Escherichia coli Methylotroph via Adaptive Laboratory Evolution Enhances Synthetic Methanol Utilization.} & 2021 & Scale anchor & \sourcetext{Genome-scale/omics phenotype} & -- & \href{https://doi.org/10.3389/fmicb.2021.638426}{10.\allowbreak 3389/\allowbreak fmicb.\allowbreak 2021.\allowbreak 638426} \\
\addlinespace[2pt]
158 & \org{E. coli} & \sourcetext{Adaptive evolution of genomically recoded Escherichia coli.} & 2018 & Scale anchor & \sourcetext{Genome-scale/omics phenotype} & -- & \href{https://doi.org/10.1073/pnas.1715530115}{10.\allowbreak 1073/\allowbreak pnas.\allowbreak 1715530115} \\
\addlinespace[2pt]
159 & \org{E. coli} & \sourcetext{A plasmid toolset for CRISPR-mediated genome editing and CRISPRi gene regulation in Escherichia coli.} & 2021 & Scale anchor & \sourcetext{Genome-scale/omics phenotype} & -- & \href{https://doi.org/10.1111/1751-7915.13780}{10.\allowbreak 1111/\allowbreak 1751-\allowbreak 7915.\allowbreak 13780} \\
\addlinespace[2pt]
160 & \org{E. coli} & \sourcetext{Laboratory Evolution Reveals Transcriptional Mechanisms Underlying Thermal Adaptation of Escherichia coli.} & 2025 & Scale anchor & \sourcetext{Genome-scale/omics phenotype} & -- & \href{https://doi.org/10.1093/gbe/evaf171}{10.\allowbreak 1093/\allowbreak gbe/\allowbreak evaf171} \\
\addlinespace[2pt]
161 & \org{E. coli} & \sourcetext{Modular and signal-responsive transcriptional regulation using CRISPRi-aided genetic switches in Escherichia coli.} & 2025 & Scale anchor & \sourcetext{Genome-scale/omics phenotype} & -- & \href{https://doi.org/10.1186/s13036-025-00526-8}{10.\allowbreak 1186/\allowbreak s13036-\allowbreak 025-\allowbreak 00526-\allowbreak 8} \\
\addlinespace[2pt]
162 & \org{E. coli} & \sourcetext{Non-native Pathway Engineering with CRISPRi for Carbon Dioxide Assimilation and Valued 5-Aminolevulinic Acid Synthesis in Escherichia coli Nissle.} & 2024 & Scale anchor & \sourcetext{Genome-scale/omics phenotype} & -- & \href{https://doi.org/10.1021/acssynbio.4c00318}{10.\allowbreak 1021/\allowbreak acssynbio.\allowbreak 4c00318} \\
\addlinespace[2pt]
163 & \org{E. coli} & \sourcetext{Genome-wide screen reveals cellular functions that counteract rifampicin lethality in Escherichia coli.} & 2024 & Scale anchor & \sourcetext{Genome-scale/omics phenotype} & -- & \href{https://doi.org/10.1128/spectrum.02895-23}{10.\allowbreak 1128/\allowbreak spectrum.\allowbreak 02895-\allowbreak 23} \\
\addlinespace[2pt]
165 & \org{E. coli} & \sourcetext{Genome-wide screen reveals a universal role of ATP in ciprofloxacin tolerance among genetically distinct Escherichia coli persisters.} & 2026 & Scale anchor & \sourcetext{Genome-scale/omics phenotype} & -- & \href{https://doi.org/10.1002/mlf2.70072}{10.\allowbreak 1002/\allowbreak mlf2.\allowbreak 70072} \\
\addlinespace[2pt]
166 & \org{E. coli} & \sourcetext{Using CRISPRi in Escherichia coli to emphasize experimental controls in a molecular microbiology laboratory.} & 2025 & Scale anchor & \sourcetext{Genome-scale/omics phenotype} & -- & \href{https://doi.org/10.1128/jmbe.00135-25}{10.\allowbreak 1128/\allowbreak jmbe.\allowbreak 00135-\allowbreak 25} \\
\addlinespace[2pt]
167 & \org{E. coli} & \sourcetext{Promoters Constrain Evolution of Expression Levels of Essential Genes in Escherichia coli.} & 2024 & Scale anchor & \sourcetext{Genome-scale/omics phenotype} & -- & \href{https://doi.org/10.1093/molbev/msae185}{10.\allowbreak 1093/\allowbreak molbev/\allowbreak msae185} \\
\addlinespace[2pt]
168 & \org{E. coli} & \sourcetext{Synthetic acid stress-tolerance modules improve growth robustness and lysine productivity of industrial Escherichia coli in fermentation at low pH.} & 2022 & Scale anchor & \sourcetext{Genome-scale/omics phenotype} & -- & \href{https://doi.org/10.1186/s12934-022-01795-4}{10.\allowbreak 1186/\allowbreak s12934-\allowbreak 022-\allowbreak 01795-\allowbreak 4} \\
\addlinespace[2pt]
169 & \org{E. coli} & \sourcetext{CRISPRi enables fast growth followed by stable aerobic pyruvate formation in Escherichia coli without auxotrophy.} & 2022 & Scale anchor & \sourcetext{Genome-scale/omics phenotype} & -- & \href{https://doi.org/10.1002/elsc.202100021}{10.\allowbreak 1002/\allowbreak elsc.\allowbreak 202100021} \\
\addlinespace[2pt]
170 & \org{E. coli} & \sourcetext{High-Throughput Screening of a Promoter Library Reveals New Persister Mechanisms in Escherichia Coli.} & 2022 & Scale anchor & \sourcetext{Genome-scale/omics phenotype} & -- & \href{https://doi.org/10.1128/spectrum.02253-21}{10.\allowbreak 1128/\allowbreak spectrum.\allowbreak 02253-\allowbreak 21} \\
\addlinespace[2pt]
171 & \org{E. coli} & \sourcetext{Inhibition of Plasmid Conjugation in Escherichia coli by Targeting rbsB Gene Using CRISPRi System.} & 2023 & Scale anchor & \sourcetext{Genome-scale/omics phenotype} & -- & \href{https://doi.org/10.3390/ijms241310585}{10.\allowbreak 3390/\allowbreak ijms241310585} \\
\addlinespace[2pt]
172 & \org{E. coli} & \sourcetext{Improvement of sabinene tolerance of Escherichia coli using adaptive laboratory evolution and omics technologies.} & 2020 & Scale anchor & \sourcetext{Genome-scale/omics phenotype} & -- & \href{https://doi.org/10.1186/s13068-020-01715-x}{10.\allowbreak 1186/\allowbreak s13068-\allowbreak 020-\allowbreak 01715-\allowbreak x} \\
\addlinespace[2pt]
173 & \org{E. coli} & \sourcetext{Genome-wide screen of genetic determinants that govern Escherichia coli growth and persistence in lake water.} & 2024 & Scale anchor & \sourcetext{Genome-scale/omics phenotype} & -- & \href{https://doi.org/10.1093/ismejo/wrae096}{10.\allowbreak 1093/\allowbreak ismejo/\allowbreak wrae096} \\
\addlinespace[2pt]
174 & \org{E. coli} & \sourcetext{Involvement of catalase and superoxide dismutase in hydrophobic organic solvent tolerance of Escherichia coli.} & 2021 & Scale anchor & \sourcetext{Genome-scale/omics phenotype} & -- & \href{https://doi.org/10.1186/s13568-021-01258-w}{10.\allowbreak 1186/\allowbreak s13568-\allowbreak 021-\allowbreak 01258-\allowbreak w} \\
\addlinespace[2pt]
175 & \org{E. coli} & \sourcetext{Filamentation and restoration of normal growth in Escherichia coli using a combined CRISPRi sgRNA/antisense RNA approach.} & 2018 & Scale anchor & \sourcetext{Genome-scale/omics phenotype} & -- & \href{https://doi.org/10.1371/journal.pone.0198058}{10.\allowbreak 1371/\allowbreak journal.\allowbreak pone.\allowbreak 0198058} \\
\addlinespace[2pt]
176 & \org{E. coli} & \sourcetext{Engineering and adaptive laboratory evolution of Escherichia coli for improving methanol utilization based on a hybrid methanol assimilation pathway.} & 2022 & Scale anchor & \sourcetext{Genome-scale/omics phenotype} & -- & \href{https://doi.org/10.3389/fbioe.2022.1089639}{10.\allowbreak 3389/\allowbreak fbioe.\allowbreak 2022.\allowbreak 1089639} \\
\addlinespace[2pt]
177 & \org{E. coli} & \sourcetext{Adaptive Laboratory Evolution of Flavin Functionality Identifies Dihydrolipoyl Dehydrogenase as One of the Critical Points for the Activity of 7,8-Didemethyl-Riboflavin as a Surrogate for Riboflavin in Escherichia coli.} & 2024 & Scale anchor & \sourcetext{Genome-scale/omics phenotype} & -- & \href{https://doi.org/10.3390/molecules29245891}{10.\allowbreak 3390/\allowbreak molecules29245891} \\
\addlinespace[2pt]
178 & \org{E. coli} & \sourcetext{Conversion of Escherichia coli into Mixotrophic CO2 Assimilation with Malate and Hydrogen Based on Recombinant Expression of 2-Oxoglutarate:Ferredoxin Oxidoreductase Using Adaptive Laboratory Evolution.} & 2023 & Scale anchor & \sourcetext{Genome-scale/omics phenotype} & -- & \href{https://doi.org/10.3390/microorganisms11020253}{10.\allowbreak 3390/\allowbreak microorganisms11020253} \\
\addlinespace[2pt]
179 & \org{E. coli} & \sourcetext{A genome-wide screen reveals the involvement of enterobactin-mediated iron acquisition in Escherichia coli survival during copper stress.} & 2021 & Scale anchor & \sourcetext{Genome-scale/omics phenotype} & -- & \href{https://doi.org/10.1093/mtomcs/mfab052}{10.\allowbreak 1093/\allowbreak mtomcs/\allowbreak mfab052} \\
\addlinespace[2pt]
180 & \org{E. coli} & \sourcetext{Inactivation of a Mismatch-Repair System Diversifies Genotypic Landscape of Escherichia coli During Adaptive Laboratory Evolution.} & 2019 & Scale anchor & \sourcetext{Genome-scale/omics phenotype} & -- & \href{https://doi.org/10.3389/fmicb.2019.01845}{10.\allowbreak 3389/\allowbreak fmicb.\allowbreak 2019.\allowbreak 01845} \\
\addlinespace[2pt]
181 & \org{E. coli} & \sourcetext{Improved furfural tolerance in Escherichia coli mediated by heterologous NADH-dependent benzyl alcohol dehydrogenases.} & 2022 & Scale anchor & \sourcetext{Genome-scale/omics phenotype} & -- & \href{https://doi.org/10.1042/bcj20210811}{10.\allowbreak 1042/\allowbreak bcj20210811} \\
\addlinespace[2pt]
182 & \org{E. coli} & \sourcetext{Escherichia coli Data-Driven Strain Design Using Aggregated Adaptive Laboratory Evolution Mutational Data.} & 2021 & Scale anchor & \sourcetext{Genome-scale/omics phenotype} & -- & \href{https://doi.org/10.1021/acssynbio.1c00337}{10.\allowbreak 1021/\allowbreak acssynbio.\allowbreak 1c00337} \\
\addlinespace[2pt]
183 & \org{E. coli} & \sourcetext{Dysregulated NAD(H) homeostasis associated with ciprofloxacin tolerance in Escherichia coli investigated on a single-cell level with the Peredox [NADH:NAD+] biosensor.} & 2023 & Scale anchor & \sourcetext{Genome-scale/omics phenotype} & -- & \href{https://doi.org/10.3389/fmicb.2023.1191968}{10.\allowbreak 3389/\allowbreak fmicb.\allowbreak 2023.\allowbreak 1191968} \\
\addlinespace[2pt]
184 & \org{E. coli} & \sourcetext{Functional Heterologous Expression of Mature Lipase LipA from Pseudomonas aeruginosa PSA01 in Escherichia coli SHuffle and BL21 (DE3): Effect of the Expression Host on Thermal Stability and Solvent Tolerance of the Enzyme Produced.} & 2020 & Scale anchor & \sourcetext{Genome-scale/omics phenotype} & -- & \href{https://doi.org/10.3390/ijms21113925}{10.\allowbreak 3390/\allowbreak ijms21113925} \\
\addlinespace[2pt]
185 & \org{E. coli} & \sourcetext{Significantly improved solvent tolerance of Escherichia coli by global transcription machinery engineering.} & 2015 & Scale anchor & \sourcetext{Genome-scale/omics phenotype} & -- & \href{https://doi.org/10.1186/s12934-015-0368-4}{10.\allowbreak 1186/\allowbreak s12934-\allowbreak 015-\allowbreak 0368-\allowbreak 4} \\
\addlinespace[2pt]
186 & \org{E. coli} & \sourcetext{Enhancing butanol tolerance of Escherichia coli reveals hydrophobic interaction of multi-tasking chaperone SecB.} & 2019 & Scale anchor & \sourcetext{Genome-scale/omics phenotype} & -- & \href{https://doi.org/10.1186/s13068-019-1507-7}{10.\allowbreak 1186/\allowbreak s13068-\allowbreak 019-\allowbreak 1507-\allowbreak 7} \\
\addlinespace[2pt]
187 & \org{E. coli} & \sourcetext{Interplay between tolerance mechanisms to copper and acid stress in Escherichia coli.} & 2017 & Scale anchor & \sourcetext{Genome-scale/omics phenotype} & -- & \href{https://doi.org/10.1073/pnas.1620232114}{10.\allowbreak 1073/\allowbreak pnas.\allowbreak 1620232114} \\
\addlinespace[2pt]
188 & \org{E. coli} & \sourcetext{A Novel Butanol Tolerance-Promoting Function of the Transcription Factor Rob in Escherichia coli.} & 2020 & Scale anchor & \sourcetext{Genome-scale/omics phenotype} & -- & \href{https://doi.org/10.3389/fbioe.2020.524198}{10.\allowbreak 3389/\allowbreak fbioe.\allowbreak 2020.\allowbreak 524198} \\
\addlinespace[2pt]
189 & \org{E. coli} & \sourcetext{Genome-wide mapping of furfural tolerance genes in Escherichia coli.} & 2014 & Scale anchor & \sourcetext{Genome-scale/omics phenotype} & -- & \href{https://doi.org/10.1371/journal.pone.0087540}{10.\allowbreak 1371/\allowbreak journal.\allowbreak pone.\allowbreak 0087540} \\
\addlinespace[2pt]
190 & \org{E. coli} & \sourcetext{An ultra-dense library resource for rapid deconvolution of mutations that cause phenotypes in Escherichia coli.} & 2016 & Scale anchor & \sourcetext{Genome-scale/omics phenotype} & -- & \href{https://doi.org/10.1093/nar/gkv1131}{10.\allowbreak 1093/\allowbreak nar/\allowbreak gkv1131} \\
\addlinespace[2pt]
191 & \org{E. coli} & \sourcetext{Expression of TolC and Organic Solvent Tolerance of Escherichia Coli Ciprofloxacin Resistant Mutants.} & 2017 & Scale anchor & \sourcetext{Genome-scale/omics phenotype} & -- & \href{https://pmc.ncbi.nlm.nih.gov/articles/PMC5610773/}{PMC5610773} \\
\addlinespace[2pt]
192 & \org{E. coli} & \sourcetext{Creation of a Dense Transposon Insertion Library Using Bacterial Conjugation in Enterobacterial Strains Such As Escherichia Coli or Shigella flexneri.} & 2017 & Scale anchor & \sourcetext{Genome-scale/omics phenotype} & -- & \href{https://doi.org/10.3791/56216}{10.\allowbreak 3791/\allowbreak 56216} \\
\addlinespace[2pt]
193 & \org{E. coli} & \sourcetext{Adaptive laboratory evolution of Escherichia coli under acid stress.} & 2020 & Scale anchor & \sourcetext{Genome-scale/omics phenotype} & -- & \href{https://doi.org/10.1099/mic.0.000867}{10.\allowbreak 1099/\allowbreak mic.\allowbreak 0.\allowbreak 000867} \\
\addlinespace[2pt]
195 & \org{E. coli} & \sourcetext{A genome-wide screen in Escherichia coli reveals that ubiquinone is a key antioxidant for metabolism of long-chain fatty acids.} & 2017 & Scale anchor & \sourcetext{Genome-scale/omics phenotype} & -- & \href{https://doi.org/10.1074/jbc.m117.806240}{10.\allowbreak 1074/\allowbreak jbc.\allowbreak m117.\allowbreak 806240} \\
\addlinespace[2pt]
196 & \org{E. coli} & \sourcetext{Bioprospecting of Native Efflux Pumps To Enhance Furfural Tolerance in Ethanologenic Escherichia coli.} & 2019 & Scale anchor & \sourcetext{Genome-scale/omics phenotype} & -- & \href{https://doi.org/10.1128/aem.02985-18}{10.\allowbreak 1128/\allowbreak aem.\allowbreak 02985-\allowbreak 18} \\
\addlinespace[2pt]
197 & \org{E. coli} & \sourcetext{Benefits of a Recombination-Proficient Escherichia coli System for Adaptive Laboratory Evolution.} & 2016 & Scale anchor & \sourcetext{Genome-scale/omics phenotype} & -- & \href{https://doi.org/10.1128/aem.01850-16}{10.\allowbreak 1128/\allowbreak aem.\allowbreak 01850-\allowbreak 16} \\
\addlinespace[2pt]
198 & \org{E. coli} & \sourcetext{Visualizing evolution in real time to determine the molecular mechanisms of n-butanol tolerance in Escherichia coli.} & 2012 & Scale anchor & \sourcetext{Genome-scale/omics phenotype} & -- & \href{https://doi.org/10.1016/j.ymben.2012.05.002}{10.\allowbreak 1016/\allowbreak j.\allowbreak ymben.\allowbreak 2012.\allowbreak 05.\allowbreak 002} \\
\addlinespace[2pt]
199 & \org{E. coli} & \sourcetext{Adaptive laboratory evolution resolves energy depletion to maintain high aromatic metabolite phenotypes in Escherichia coli strains lacking the Phosphotransferase System.} & 2018 & Scale anchor & \sourcetext{Genome-scale/omics phenotype} & -- & \href{https://doi.org/10.1016/j.ymben.2018.06.005}{10.\allowbreak 1016/\allowbreak j.\allowbreak ymben.\allowbreak 2018.\allowbreak 06.\allowbreak 005} \\
\addlinespace[2pt]
200 & \org{E. coli} & \sourcetext{Adaptive laboratory evolution for improved tolerance of isobutyl acetate in Escherichia coli.} & 2022 & Scale anchor & \sourcetext{Genome-scale/omics phenotype} & -- & \href{https://doi.org/10.1016/j.ymben.2021.11.002}{10.\allowbreak 1016/\allowbreak j.\allowbreak ymben.\allowbreak 2021.\allowbreak 11.\allowbreak 002} \\
\addlinespace[2pt]
\midrule
1\,\textsuperscript{\textbf{R}} & \org{P. putida} & \sourcetext{Muconic acid production from glucose and xylose in Pseudomonas putida via evolution and metabolic engineering.} & 2022 & Joint & \sourcetext{Muconic acid} & -- & \href{https://doi.org/10.1038/s41467-022-32296-y}{10.\allowbreak 1038/\allowbreak s41467-\allowbreak 022-\allowbreak 32296-\allowbreak y} \\
\addlinespace[2pt]
2\,\textsuperscript{\textbf{R}} & \org{P. putida} & \sourcetext{Tuning a high performing multiplexed-CRISPRi Pseudomonas putida strain to further enhance indigoidine production.} & 2022 & Joint & \sourcetext{PRi Pseudomonas putida strain to further enhance indigoidine} & -- & \href{https://doi.org/10.1016/j.mec.2022.e00206}{10.\allowbreak 1016/\allowbreak j.\allowbreak mec.\allowbreak 2022.\allowbreak e00206} \\
\addlinespace[2pt]
9 & \org{P. putida} & \sourcetext{Development of a CRISPR/Cas9n-based tool for metabolic engineering of Pseudomonas putida for ferulic acid-to-polyhydroxyalkanoate bioconversion.} & 2020 & Joint & \sourcetext{Genome-scale/omics phenotype} & -- & \href{https://doi.org/10.1038/s42003-020-0824-5}{10.\allowbreak 1038/\allowbreak s42003-\allowbreak 020-\allowbreak 0824-\allowbreak 5} \\
\addlinespace[2pt]
10 & \org{P. putida} & \sourcetext{Systems biology of electrogenic Pseudomonas putida - multi-omics insights and metabolic engineering for enhanced 2-ketogluconate production.} & 2024 & Joint & \sourcetext{ights and metabolic engineering for enhanced 2-ketogluconate} & -- & \href{https://doi.org/10.1186/s12934-024-02509-8}{10.\allowbreak 1186/\allowbreak s12934-\allowbreak 024-\allowbreak 02509-\allowbreak 8} \\
\addlinespace[2pt]
11 & \org{P. putida} & \sourcetext{Enhanced chaotrope tolerance and (S)-2-hydroxypropiophenone production by recombinant Pseudomonas putida engineered with Pprl from Deinococcus radiodurans.} & 2024 & Joint & \sourcetext{Enhanced chaotrope tolerance and (S)-2-hydroxypropiophenone} & -- & \href{https://doi.org/10.1111/1751-7915.14448}{10.\allowbreak 1111/\allowbreak 1751-\allowbreak 7915.\allowbreak 14448} \\
\addlinespace[2pt]
12 & \org{P. putida} & \sourcetext{Enhancement of polyhydroxyalkanoate production by co-feeding lignin derivatives with glycerol in Pseudomonas putida KT2440.} & 2021 & ME breadth & \sourcetext{Enhancement of polyhydroxyalkanoate} & -- & \href{https://doi.org/10.1186/s13068-020-01861-2}{10.\allowbreak 1186/\allowbreak s13068-\allowbreak 020-\allowbreak 01861-\allowbreak 2} \\
\addlinespace[2pt]
13 & \org{P. putida} & \sourcetext{Metabolic engineering of Pseudomonas putida for increased polyhydroxyalkanoate production from lignin.} & 2020 & ME breadth & \sourcetext{ing of Pseudomonas putida for increased polyhydroxyalkanoate} & -- & \href{https://doi.org/10.1111/1751-7915.13481}{10.\allowbreak 1111/\allowbreak 1751-\allowbreak 7915.\allowbreak 13481} \\
\addlinespace[2pt]
14\,\textsuperscript{\textbf{R}} & \org{P. putida} & \sourcetext{Engineering Pseudomonas putida for production of 3-hydroxyacids using hybrid type I polyketide synthases.} & 2025 & ME breadth & \sourcetext{3-hydroxyacids using hybrid type I polyketide synthases} & -- & \href{https://doi.org/10.1016/j.mec.2025.e00261}{10.\allowbreak 1016/\allowbreak j.\allowbreak mec.\allowbreak 2025.\allowbreak e00261} \\
\addlinespace[2pt]
16\,\textsuperscript{\textbf{R}} & \org{P. putida} & \sourcetext{Lignin conversion to $\beta$-ketoadipic acid by Pseudomonas putida via metabolic engineering and bioprocess development.} & 2023 & ME breadth & \sourcetext{Metabolic-engineering phenotype} & -- & \href{https://doi.org/10.1126/sciadv.adj0053}{10.\allowbreak 1126/\allowbreak sciadv.\allowbreak adj0053} \\
\addlinespace[2pt]
17 & \org{P. putida} & \sourcetext{A reduction in growth rate of Pseudomonas putida KT2442 counteracts productivity advances in medium-chain-length polyhydroxyalkanoate production from gluconate.} & 2011 & ME breadth & \sourcetext{ctivity advances in medium-chain-length polyhydroxyalkanoate} & -- & \href{https://doi.org/10.1186/1475-2859-10-25}{10.\allowbreak 1186/\allowbreak 1475-\allowbreak 2859-\allowbreak 10-\allowbreak 25} \\
\addlinespace[2pt]
18 & \org{P. putida} & \sourcetext{Quantifying NAD(P)H production in the upper Entner-Doudoroff pathway from Pseudomonas putida KT2440.} & 2015 & ME breadth & \sourcetext{Quantifying NAD(P)H} & -- & \href{https://doi.org/10.1016/j.fob.2015.11.002}{10.\allowbreak 1016/\allowbreak j.\allowbreak fob.\allowbreak 2015.\allowbreak 11.\allowbreak 002} \\
\addlinespace[2pt]
19 & \org{P. putida} & \sourcetext{Integration of Genetic and Process Engineering for Optimized Rhamnolipid Production Using Pseudomonas putida.} & 2020 & ME breadth & \sourcetext{of Genetic and Process Engineering for Optimized Rhamnolipid} & -- & \href{https://doi.org/10.3389/fbioe.2020.00976}{10.\allowbreak 3389/\allowbreak fbioe.\allowbreak 2020.\allowbreak 00976} \\
\addlinespace[2pt]
20 & \org{P. putida} & \sourcetext{Leveraging host metabolism for bisdemethoxycurcumin production in Pseudomonas putida.} & 2020 & ME breadth & \sourcetext{Leveraging host metabolism for bisdemethoxycurcumin} & -- & \href{https://doi.org/10.1016/j.mec.2019.e00119}{10.\allowbreak 1016/\allowbreak j.\allowbreak mec.\allowbreak 2019.\allowbreak e00119} \\
\addlinespace[2pt]
21 & \org{P. putida} & \sourcetext{Production of medium chain length polyhydroxyalkanoate in metabolic flux optimized Pseudomonas putida.} & 2014 & ME breadth & \sourcetext{medium chain length polyhydroxyalkanoate in metabolic flux optimized Pseudomonas putida} & -- & \href{https://doi.org/10.1186/1475-2859-13-88}{10.\allowbreak 1186/\allowbreak 1475-\allowbreak 2859-\allowbreak 13-\allowbreak 88} \\
\addlinespace[2pt]
22 & \org{P. putida} & \sourcetext{Valorization of Polyethylene Terephthalate to Muconic Acid by Engineering Pseudomonas Putida.} & 2022 & ME breadth & \sourcetext{Polyethylene Terephthalate to Muconic Acid by Engineering Pseudomonas Putida} & -- & \href{https://doi.org/10.3390/ijms231910997}{10.\allowbreak 3390/\allowbreak ijms231910997} \\
\addlinespace[2pt]
23 & \org{P. putida} & \sourcetext{Large-scale kinetic metabolic models of Pseudomonas putida KT2440 for consistent design of metabolic engineering strategies.} & 2020 & ME breadth & \sourcetext{Metabolic-engineering phenotype} & -- & \href{https://doi.org/10.1186/s13068-020-1665-7}{10.\allowbreak 1186/\allowbreak s13068-\allowbreak 020-\allowbreak 1665-\allowbreak 7} \\
\addlinespace[2pt]
24 & \org{P. putida} & \sourcetext{Machine learning-led semi-automated medium optimization reveals salt as key for flaviolin production in Pseudomonas putida.} & 2025 & ME breadth & \sourcetext{omated medium optimization reveals salt as key for flaviolin} & -- & \href{https://doi.org/10.1038/s42003-025-08039-2}{10.\allowbreak 1038/\allowbreak s42003-\allowbreak 025-\allowbreak 08039-\allowbreak 2} \\
\addlinespace[2pt]
25 & \org{P. putida} & \sourcetext{Anoxic metabolism and biochemical production in Pseudomonas putida F1 driven by a bioelectrochemical system.} & 2016 & ME breadth & \sourcetext{Anoxic metabolism and biochemical} & -- & \href{https://doi.org/10.1186/s13068-016-0452-y}{10.\allowbreak 1186/\allowbreak s13068-\allowbreak 016-\allowbreak 0452-\allowbreak y} \\
\addlinespace[2pt]
27 & \org{P. putida} & \sourcetext{Biocatalytic production of perillyl alcohol from limonene by using a novel Mycobacterium sp. cytochrome P450 alkane hydroxylase expressed in Pseudomonas putida.} & 2005 & ME breadth & \sourcetext{perillyl alcohol from limonene by using a novel Mycobacterium sp} & -- & \href{https://doi.org/10.1128/aem.71.4.1737-1744.2005}{10.\allowbreak 1128/\allowbreak aem.\allowbreak 71.\allowbreak 4.\allowbreak 1737-\allowbreak 1744.\allowbreak 2005} \\
\addlinespace[2pt]
28 & \org{P. putida} & \sourcetext{Fed-Batch mcl- Polyhydroxyalkanoates Production in Pseudomonas putida KT2440 and $\Delta$phaZ Mutant on Biodiesel-Derived Crude Glycerol.} & 2021 & ME breadth & \sourcetext{Fed-Batch mcl- Polyhydroxyalkanoates} & -- & \href{https://doi.org/10.3389/fbioe.2021.642023}{10.\allowbreak 3389/\allowbreak fbioe.\allowbreak 2021.\allowbreak 642023} \\
\addlinespace[2pt]
29 & \org{P. putida} & \sourcetext{Leveraging Engineered Pseudomonas putida Minicells for Bioconversion of Organic Acids into Short-Chain Methyl Ketones.} & 2025 & ME breadth & \sourcetext{Organic Acids into Short-Chain Methyl Ketones} & -- & \href{https://doi.org/10.1021/acssynbio.4c00700}{10.\allowbreak 1021/\allowbreak acssynbio.\allowbreak 4c00700} \\
\addlinespace[2pt]
30 & \org{P. putida} & \sourcetext{Production of selenium nanoparticles occurs through an interconnected pathway of sulphur metabolism and oxidative stress response in Pseudomonas putida KT2440.} & 2023 & ME breadth & \sourcetext{selenium nanoparticles occurs through an interconnected pathway of sulphur metabolism and oxidative} & -- & \href{https://doi.org/10.1111/1751-7915.14215}{10.\allowbreak 1111/\allowbreak 1751-\allowbreak 7915.\allowbreak 14215} \\
\addlinespace[2pt]
31 & \org{P. putida} & \sourcetext{Efficient recombinant production of prodigiosin in Pseudomonas putida.} & 2015 & ME breadth & \sourcetext{prodigiosin} & -- & \href{https://doi.org/10.3389/fmicb.2015.00972}{10.\allowbreak 3389/\allowbreak fmicb.\allowbreak 2015.\allowbreak 00972} \\
\addlinespace[2pt]
32 & \org{P. putida} & \sourcetext{De novo production of the monoterpenoid geranic acid by metabolically engineered Pseudomonas putida.} & 2014 & ME breadth & \sourcetext{the monoterpenoid geranic acid by metabolically engineered Pseudomonas putida} & adjacent & \href{https://doi.org/10.1186/s12934-014-0170-8}{10.\allowbreak 1186/\allowbreak s12934-\allowbreak 014-\allowbreak 0170-\allowbreak 8} \\
\addlinespace[2pt]
33 & \org{P. putida} & \sourcetext{Production of Vanillin From Ferulic Acid by Pseudomonas putida KT2440 Using Metabolic Engineering and In~Situ Product Recovery.} & 2025 & ME breadth & \sourcetext{Vanillin From Ferulic Acid} & -- & \href{https://doi.org/10.1111/1751-7915.70152}{10.\allowbreak 1111/\allowbreak 1751-\allowbreak 7915.\allowbreak 70152} \\
\addlinespace[2pt]
34 & \org{P. putida} & \sourcetext{Enhanced biosynthesis of poly(3-hydroxybutyrate) in engineered strains of Pseudomonas putida via increased malonyl-CoA availability.} & 2024 & ME breadth & \sourcetext{poly(3-hydroxybutyrate) in engineered strains of Pseudomonas putida via increased malonyl-CoA availa} & -- & \href{https://doi.org/10.1111/1751-7915.70044}{10.\allowbreak 1111/\allowbreak 1751-\allowbreak 7915.\allowbreak 70044} \\
\addlinespace[2pt]
35 & \org{P. putida} & \sourcetext{Microbial Valorization of Lignin to Bioplastic by Genome-Reduced Pseudomonas putida.} & 2022 & ME breadth & \sourcetext{Lignin to Bioplastic by Genome-Reduced Pseudomonas putida} & -- & \href{https://doi.org/10.3389/fmicb.2022.923664}{10.\allowbreak 3389/\allowbreak fmicb.\allowbreak 2022.\allowbreak 923664} \\
\addlinespace[2pt]
36 & \org{P. putida} & \sourcetext{Exploring engineered vesiculation by Pseudomonas putida KT2440 for natural product biosynthesis.} & 2024 & ME breadth & \sourcetext{Metabolic-engineering phenotype} & -- & \href{https://doi.org/10.1111/1751-7915.14312}{10.\allowbreak 1111/\allowbreak 1751-\allowbreak 7915.\allowbreak 14312} \\
\addlinespace[2pt]
37\,\textsuperscript{\textbf{R}} & \org{P. putida} & \sourcetext{Metabolic Engineering of Pseudomonas putida KT2440 for the Production of para-Hydroxy Benzoic Acid.} & 2016 & ME breadth & \sourcetext{para-Hydroxy Benzoic Acid} & -- & \href{https://doi.org/10.3389/fbioe.2016.00090}{10.\allowbreak 3389/\allowbreak fbioe.\allowbreak 2016.\allowbreak 00090} \\
\addlinespace[2pt]
38 & \org{P. putida} & \sourcetext{Bioprocess exploitation of microaerobic auto-induction using the example of rhamnolipid biosynthesis in Pseudomonas putida KT2440.} & 2025 & ME breadth & \sourcetext{Metabolic-engineering phenotype} & -- & \href{https://doi.org/10.1186/s13036-025-00478-z}{10.\allowbreak 1186/\allowbreak s13036-\allowbreak 025-\allowbreak 00478-\allowbreak z} \\
\addlinespace[2pt]
39 & \org{P. putida} & \sourcetext{Synthetic pathways for microbial biosynthesis of valuable pyrazine derivatives using genetically modified Pseudomonas putida KT2440.} & 2025 & ME breadth & \sourcetext{valuable pyrazine derivatives using genetically modified Pseudomonas putida KT2440} & -- & \href{https://doi.org/10.1016/j.mec.2025.e00258}{10.\allowbreak 1016/\allowbreak j.\allowbreak mec.\allowbreak 2025.\allowbreak e00258} \\
\addlinespace[2pt]
40 & \org{P. putida} & \sourcetext{Pseudomonas putida as a platform for medium-chain length $\alpha$,$\omega$-diol production: Opportunities and challenges.} & 2024 & ME breadth & \sourcetext{diol} & -- & \href{https://doi.org/10.1111/1751-7915.14423}{10.\allowbreak 1111/\allowbreak 1751-\allowbreak 7915.\allowbreak 14423} \\
\addlinespace[2pt]
41 & \org{P. putida} & \sourcetext{Portable bacterial CRISPR transcriptional activation enables metabolic engineering in Pseudomonas putida.} & 2021 & Joint & \sourcetext{Genome-scale/omics phenotype} & -- & \href{https://doi.org/10.1016/j.ymben.2021.04.002}{10.\allowbreak 1016/\allowbreak j.\allowbreak ymben.\allowbreak 2021.\allowbreak 04.\allowbreak 002} \\
\addlinespace[2pt]
42 & \org{P. putida} & \sourcetext{Biosynthesis of Novel Aromatic Copolyesters from Insoluble 11-Phenoxyundecanoic Acid by Pseudomonas putida BM01.} & 1996 & ME breadth & \sourcetext{Novel Aromatic Copolyesters from Insoluble 11-Phenoxyundecanoic Acid} & -- & \href{https://doi.org/10.1128/aem.62.2.536-544.1996}{10.\allowbreak 1128/\allowbreak aem.\allowbreak 62.\allowbreak 2.\allowbreak 536-\allowbreak 544.\allowbreak 1996} \\
\addlinespace[2pt]
43 & \org{P. putida} & \sourcetext{Engineered Membrane Vesicle Production via oprF or oprI Deletion Has Distinct Phenotypic Effects in Pseudomonas putida.} & 2025 & ME breadth & \sourcetext{Engineered Membrane Vesicle} & -- & \href{https://doi.org/10.1021/acssynbio.5c00171}{10.\allowbreak 1021/\allowbreak acssynbio.\allowbreak 5c00171} \\
\addlinespace[2pt]
44 & \org{P. putida} & \sourcetext{Metabolic Engineering of Pseudomonas putida KT2440 for $\beta$-Nicotinamide Mononucleotide Biosynthesis from Glucose and Aspartate.} & 2026 & ME breadth & \sourcetext{Metabolic-engineering phenotype} & -- & \href{https://doi.org/10.3390/microorganisms14091877}{10.\allowbreak 3390/\allowbreak microorganisms14091877} \\
\addlinespace[2pt]
45 & \org{P. putida} & \sourcetext{High-level prodigiosin production in Pseudomonas putida enabled by combinatorial metabolic engineering.} & 2026 & ME breadth & \sourcetext{High-level prodigiosin} & -- & \href{https://doi.org/10.1016/j.synbio.2026.01.015}{10.\allowbreak 1016/\allowbreak j.\allowbreak synbio.\allowbreak 2026.\allowbreak 01.\allowbreak 015} \\
\addlinespace[2pt]
46 & \org{P. putida} & \sourcetext{Characterisation of metabolic burden in Pseudomonas putida reveals precursor limitation in heterologous lycopene production.} & 2026 & ME breadth & \sourcetext{putida reveals precursor limitation in heterologous lycopene} & adjacent & \href{https://doi.org/10.1186/s12934-026-03108-5}{10.\allowbreak 1186/\allowbreak s12934-\allowbreak 026-\allowbreak 03108-\allowbreak 5} \\
\addlinespace[2pt]
47 & \org{P. putida} & \sourcetext{Engineering Pseudomonas putida for the Production of 2,3-Butanediol Isomers and Acetoin.} & 2026 & ME breadth & \sourcetext{2} & -- & \href{https://doi.org/10.1021/acssynbio.6c00243}{10.\allowbreak 1021/\allowbreak acssynbio.\allowbreak 6c00243} \\
\addlinespace[2pt]
48 & \org{P. putida} & \sourcetext{Acetate as alternative carbon source for production of mono- and di-rhamnolipids in Pseudomonas putida KT2440.} & 2026 & ME breadth & \sourcetext{mono- and di-rhamnolipids} & -- & \href{https://doi.org/10.1186/s12934-026-03050-6}{10.\allowbreak 1186/\allowbreak s12934-\allowbreak 026-\allowbreak 03050-\allowbreak 6} \\
\addlinespace[2pt]
49 & \org{P. putida} & \sourcetext{Development of an ethanol-based microbial platform for 3-hydroxypropionic acid production using engineered Pseudomonas putida KT2440.} & 2025 & ME breadth & \sourcetext{ethanol-based microbial platform for 3-hydroxypropionic acid} & -- & \href{https://doi.org/10.1186/s13036-025-00595-9}{10.\allowbreak 1186/\allowbreak s13036-\allowbreak 025-\allowbreak 00595-\allowbreak 9} \\
\addlinespace[2pt]
50 & \org{P. putida} & \sourcetext{Combinatorial engineering pinpoints shikimate pathway bottlenecks in para-aminobenzoic acid production in Pseudomonas putida.} & 2025 & ME breadth & \sourcetext{ints shikimate pathway bottlenecks in para-aminobenzoic acid} & -- & \href{https://doi.org/10.1186/s13036-025-00553-5}{10.\allowbreak 1186/\allowbreak s13036-\allowbreak 025-\allowbreak 00553-\allowbreak 5} \\
\addlinespace[2pt]
51 & \org{P. putida} & \sourcetext{Citrate Synthase Overexpression of Pseudomonas putida Increases Succinate Production from Acetate in Microaerobic Cultivation.} & 2023 & ME breadth & \sourcetext{ase Overexpression of Pseudomonas putida Increases Succinate} & -- & \href{https://doi.org/10.1021/acsomega.3c02520}{10.\allowbreak 1021/\allowbreak acsomega.\allowbreak 3c02520} \\
\addlinespace[2pt]
52 & \org{P. putida} & \sourcetext{Biosynthesis of fragrance 2-phenylethanol from sugars by Pseudomonas putida.} & 2024 & ME breadth & \sourcetext{fragrance 2-phenylethanol from sugars} & -- & \href{https://doi.org/10.1186/s13068-024-02498-1}{10.\allowbreak 1186/\allowbreak s13068-\allowbreak 024-\allowbreak 02498-\allowbreak 1} \\
\addlinespace[2pt]
53\,\textsuperscript{\textbf{R}} & \org{P. putida} & \sourcetext{Engineering Pseudomonas putida KT2440 for chain length tailored free fatty acid and oleochemical production.} & 2022 & ME breadth & \sourcetext{0 for chain length tailored free fatty acid and oleochemical} & -- & \href{https://doi.org/10.1038/s42003-022-04336-2}{10.\allowbreak 1038/\allowbreak s42003-\allowbreak 022-\allowbreak 04336-\allowbreak 2} \\
\addlinespace[2pt]
54 & \org{P. putida} & \sourcetext{Metabolic engineering of Pseudomonas putida for production of vanillylamine from lignin-derived substrates.} & 2021 & ME breadth & \sourcetext{vanillylamine from lignin-derived substrates} & -- & \href{https://doi.org/10.1111/1751-7915.13764}{10.\allowbreak 1111/\allowbreak 1751-\allowbreak 7915.\allowbreak 13764} \\
\addlinespace[2pt]
55 & \org{P. putida} & \sourcetext{Engineering Pseudomonas putida for isoprenoid production by manipulating endogenous and shunt pathways supplying precursors.} & 2019 & ME breadth & \sourcetext{Engineering Pseudomonas putida for isoprenoid} & adjacent & \href{https://doi.org/10.1186/s12934-019-1204-z}{10.\allowbreak 1186/\allowbreak s12934-\allowbreak 019-\allowbreak 1204-\allowbreak z} \\
\addlinespace[2pt]
56 & \org{P. putida} & \sourcetext{Strategy for Optimizing Vitamin B12 Production in Pseudomonas putida KT2440 Using Metabolic Modeling.} & 2024 & ME breadth & \sourcetext{Strategy for Optimizing Vitamin B12} & -- & \href{https://doi.org/10.3390/metabo14110636}{10.\allowbreak 3390/\allowbreak metabo14110636} \\
\addlinespace[2pt]
57 & \org{P. putida} & \sourcetext{Bioconversion of L-Tyrosine into p-Coumaric Acid by Tyrosine Ammonia-Lyase Heterologue of Rhodobacter sphaeroides Produced in Pseudomonas putida KT2440.} & 2024 & ME breadth & \sourcetext{L-Tyrosine into p-Coumaric Acid by Tyrosine Ammonia-Lyase Heterologue of Rhodobacter sphaeroides Pro} & -- & \href{https://doi.org/10.3390/cimb46090603}{10.\allowbreak 3390/\allowbreak cimb46090603} \\
\addlinespace[2pt]
58 & \org{P. putida} & \sourcetext{Quantum modeling simulates nutrient effect of bioplastic polyhydroxyalkanoate (PHA) production in Pseudomonas putida.} & 2024 & ME breadth & \sourcetext{tes nutrient effect of bioplastic polyhydroxyalkanoate (PHA)} & -- & \href{https://doi.org/10.1038/s41598-024-68727-7}{10.\allowbreak 1038/\allowbreak s41598-\allowbreak 024-\allowbreak 68727-\allowbreak 7} \\
\addlinespace[2pt]
59 & \org{P. putida} & \sourcetext{Heterologous constitutive production of short-chain-length polyhydroxyalkanoates in Pseudomonas putida KT2440: the involvement of IbpA inclusion body protein.} & 2023 & ME breadth & \sourcetext{short-chain-length polyhydroxyalkanoates} & -- & \href{https://doi.org/10.3389/fbioe.2023.1275036}{10.\allowbreak 3389/\allowbreak fbioe.\allowbreak 2023.\allowbreak 1275036} \\
\addlinespace[2pt]
60 & \org{P. putida} & \sourcetext{Enhanced Protocatechuic Acid Production From Glucose Using Pseudomonas putida 3-Dehydroshikimate Dehydratase Expressed in a Phenylalanine-Overproducing Mutant of Escherichia coli.} & 2021 & ME breadth & \sourcetext{Enhanced Protocatechuic Acid} & -- & \href{https://doi.org/10.3389/fbioe.2021.695704}{10.\allowbreak 3389/\allowbreak fbioe.\allowbreak 2021.\allowbreak 695704} \\
\addlinespace[2pt]
61 & \org{P. putida} & \sourcetext{Channelling carbon flux through the meta-cleavage route for improved poly(3-hydroxyalkanoate) production from benzoate and lignin-based aromatics in Pseudomonas putida H.} & 2021 & ME breadth & \sourcetext{he meta-cleavage route for improved poly(3-hydroxyalkanoate)} & -- & \href{https://doi.org/10.1111/1751-7915.13705}{10.\allowbreak 1111/\allowbreak 1751-\allowbreak 7915.\allowbreak 13705} \\
\addlinespace[2pt]
62 & \org{P. putida} & \sourcetext{Vanillin Production in Pseudomonas: Whole-Genome Sequencing of Pseudomonas sp. Strain 9.1 and Reannotation of Pseudomonas putida CalA as a Vanillin Reductase.} & 2020 & ME breadth & \sourcetext{Vanillin} & -- & \href{https://doi.org/10.1128/aem.02442-19}{10.\allowbreak 1128/\allowbreak aem.\allowbreak 02442-\allowbreak 19} \\
\addlinespace[2pt]
63 & \org{P. putida} & \sourcetext{Engineering Pseudomonas putida KT2440 for the production of isobutanol.} & 2020 & ME breadth & \sourcetext{isobutanol} & -- & \href{https://doi.org/10.1002/elsc.201900151}{10.\allowbreak 1002/\allowbreak elsc.\allowbreak 201900151} \\
\addlinespace[2pt]
64 & \org{P. putida} & \sourcetext{Effect of pseudobactin 358 production by Pseudomonas putida WCS358 on suppression of fusarium wilt of carnations by nonpathogenic Fusarium oxysporum Fo47.} & 1992 & ME breadth & \sourcetext{Effect of pseudobactin 358} & -- & \href{https://doi.org/10.1128/aem.58.9.2978-2982.1992}{10.\allowbreak 1128/\allowbreak aem.\allowbreak 58.\allowbreak 9.\allowbreak 2978-\allowbreak 2982.\allowbreak 1992} \\
\addlinespace[2pt]
65 & \org{P. putida} & \sourcetext{Optimizing Systems for Robust Heterologous Production of Biosurfactants Rhamnolipid and Lyso-Ornithine Lipid in Pseudomonas putida KT2440.} & 2024 & ME breadth & \sourcetext{Biosurfactants Rhamnolipid and Lyso-Ornithine Lipid} & -- & \href{https://doi.org/10.3390/molecules29143288}{10.\allowbreak 3390/\allowbreak molecules29143288} \\
\addlinespace[2pt]
66 & \org{P. putida} & \sourcetext{Volatiles of Shiraia fruiting body-associated Pseudomonas putida No.24 stimulate fungal hypocrellin production.} & 2023 & ME breadth & \sourcetext{24 stimulate fungal hypocrellin} & -- & \href{https://doi.org/10.1016/j.synbio.2023.06.004}{10.\allowbreak 1016/\allowbreak j.\allowbreak synbio.\allowbreak 2023.\allowbreak 06.\allowbreak 004} \\
\addlinespace[2pt]
67 & \org{P. putida} & \sourcetext{Genomic Epidemiology of MBL-Producing Pseudomonas putida Group Isolates in Poland.} & 2022 & ME breadth & \sourcetext{Metabolic-engineering phenotype} & -- & \href{https://doi.org/10.1007/s40121-022-00659-z}{10.\allowbreak 1007/\allowbreak s40121-\allowbreak 022-\allowbreak 00659-\allowbreak z} \\
\addlinespace[2pt]
68 & \org{P. putida} & \sourcetext{Micro-aerobic production of isobutanol with engineered Pseudomonas putida.} & 2021 & ME breadth & \sourcetext{isobutanol} & -- & \href{https://doi.org/10.1002/elsc.202000116}{10.\allowbreak 1002/\allowbreak elsc.\allowbreak 202000116} \\
\addlinespace[2pt]
69 & \org{P. putida} & \sourcetext{Metabolic Engineering of Pseudomonas putida KT2440 to Produce Anthranilate from Glucose.} & 2015 & ME breadth & \sourcetext{Metabolic-engineering phenotype} & -- & \href{https://doi.org/10.3389/fmicb.2015.01310}{10.\allowbreak 3389/\allowbreak fmicb.\allowbreak 2015.\allowbreak 01310} \\
\addlinespace[2pt]
70 & \org{P. putida} & \sourcetext{Metabolic engineering of Pseudomonas putida for increased polyhydroxyalkanoate production from lignin.} & 2020 & ME breadth & \sourcetext{ing of Pseudomonas putida for increased polyhydroxyalkanoate} & -- & \href{https://doi.org/10.1111/1751-7915.13547}{10.\allowbreak 1111/\allowbreak 1751-\allowbreak 7915.\allowbreak 13547} \\
\addlinespace[2pt]
71 & \org{P. putida} & \sourcetext{A novel programmable lysozyme-based lysis system in Pseudomonas putida for biopolymer production.} & 2017 & ME breadth & \sourcetext{zyme-based lysis system} & -- & \href{https://doi.org/10.1038/s41598-017-04741-2}{10.\allowbreak 1038/\allowbreak s41598-\allowbreak 017-\allowbreak 04741-\allowbreak 2} \\
\addlinespace[2pt]
72 & \org{P. putida} & \sourcetext{Metabolic engineering of Pseudomonas putida for production of the natural sweetener 5-ketofructose from fructose or sucrose by periplasmic oxidation with a heterologous fructose dehydrogenase.} & 2021 & ME breadth & \sourcetext{the natural sweetener 5-ketofructose from fructose or sucrose by periplasmic oxidation with a hetero} & -- & \href{https://doi.org/10.1111/1751-7915.13913}{10.\allowbreak 1111/\allowbreak 1751-\allowbreak 7915.\allowbreak 13913} \\
\addlinespace[2pt]
73 & \org{P. putida} & \sourcetext{Comparative one-factor-at-a-time, response surface (statistical) and bench-scale bioreactor level optimization of thermoalkaline protease production from a psychrotrophic Pseudomonas putida SKG-1 isolate.} & 2011 & ME breadth & \sourcetext{ale bioreactor level optimization of thermoalkaline protease} & -- & \href{https://doi.org/10.1186/1475-2859-10-114}{10.\allowbreak 1186/\allowbreak 1475-\allowbreak 2859-\allowbreak 10-\allowbreak 114} \\
\addlinespace[2pt]
74 & \org{P. putida} & \sourcetext{Fermentative Production of N-Methylglutamate From Glycerol by Recombinant Pseudomonas putida.} & 2018 & ME breadth & \sourcetext{N-Methylglutamate From Glycerol by Recombinant Pseudomonas putida} & -- & \href{https://doi.org/10.3389/fbioe.2018.00159}{10.\allowbreak 3389/\allowbreak fbioe.\allowbreak 2018.\allowbreak 00159} \\
\addlinespace[2pt]
75 & \org{P. putida} & \sourcetext{Comparison of mcl-Poly(3-hydroxyalkanoates) synthesis by different Pseudomonas putida strains from crude glycerol: citrate accumulates at high titer under PHA-producing conditions.} & 2014 & ME breadth & \sourcetext{Metabolic-engineering phenotype} & -- & \href{https://doi.org/10.1186/s12896-014-0110-z}{10.\allowbreak 1186/\allowbreak s12896-\allowbreak 014-\allowbreak 0110-\allowbreak z} \\
\addlinespace[2pt]
76 & \org{P. putida} & \sourcetext{Selective Biosynthesis of Furoic Acid From Furfural by Pseudomonas Putida and Identification of Molybdate Transporter Involvement in Furfural Oxidation.} & 2020 & ME breadth & \sourcetext{Furoic Acid From Furfural} & -- & \href{https://doi.org/10.3389/fchem.2020.587456}{10.\allowbreak 3389/\allowbreak fchem.\allowbreak 2020.\allowbreak 587456} \\
\addlinespace[2pt]
77\,\textsuperscript{\textbf{R}} & \org{P. putida} & \sourcetext{Adaptive laboratory evolution of Pseudomonas putida KT2440 improves p-coumaric and ferulic acid catabolism and tolerance.} & 2020 & Scale anchor & \sourcetext{Genome-scale/omics phenotype} & -- & \href{https://doi.org/10.1016/j.mec.2020.e00143}{10.\allowbreak 1016/\allowbreak j.\allowbreak mec.\allowbreak 2020.\allowbreak e00143} \\
\addlinespace[2pt]
78 & \org{P. putida} & \sourcetext{Development of a high efficiency integration system and promoter library for rapid modification of Pseudomonas putida KT2440.} & 2017 & Scale anchor & \sourcetext{Genome-scale/omics phenotype} & -- & \href{https://doi.org/10.1016/j.meteno.2017.04.001}{10.\allowbreak 1016/\allowbreak j.\allowbreak meteno.\allowbreak 2017.\allowbreak 04.\allowbreak 001} \\
\addlinespace[2pt]
80 & \org{P. putida} & \sourcetext{Integration of ARTP Mutation and Adaptive Laboratory Evolution to Reveal 1,4-Butanediol Degradation in Pseudomonas putida KT2440.} & 2023 & Scale anchor & \sourcetext{Genome-scale/omics phenotype} & -- & \href{https://doi.org/10.1128/spectrum.04988-22}{10.\allowbreak 1128/\allowbreak spectrum.\allowbreak 04988-\allowbreak 22} \\
\addlinespace[2pt]
81\,\textsuperscript{\textbf{R}} & \org{P. putida} & \sourcetext{Unraveling the mechanism of furfural tolerance in engineered Pseudomonas putida by genomics.} & 2022 & Scale anchor & \sourcetext{Genome-scale/omics phenotype} & -- & \href{https://doi.org/10.3389/fmicb.2022.1035263}{10.\allowbreak 3389/\allowbreak fmicb.\allowbreak 2022.\allowbreak 1035263} \\
\addlinespace[2pt]
82 & \org{P. putida} & \sourcetext{Adaptive Laboratory Evolution Unlocks Membrane Permeability as a Key Limitation in Long-Chain Alcohol Metabolism by Pseudomonas putida KT2440.} & 2026 & Scale anchor & \sourcetext{Genome-scale/omics phenotype} & -- & \href{https://doi.org/10.1111/1751-7915.70399}{10.\allowbreak 1111/\allowbreak 1751-\allowbreak 7915.\allowbreak 70399} \\
\addlinespace[2pt]
83 & \org{P. putida} & \sourcetext{Cyclopropane fatty acids are involved in organic solvent tolerance but not in acid stress resistance in Pseudomonas putida DOT-T1E.} & 2009 & Scale anchor & \sourcetext{Genome-scale/omics phenotype} & -- & \href{https://doi.org/10.1111/j.1751-7915.2009.00084.x}{10.\allowbreak 1111/\allowbreak j.\allowbreak 1751-\allowbreak 7915.\allowbreak 2009.\allowbreak 00084.\allowbreak x} \\
\addlinespace[2pt]
84 & \org{P. putida} & \sourcetext{Adaptive Laboratory Evolution Restores Solvent Tolerance in Plasmid-Cured Pseudomonas putida S12: a Molecular Analysis.} & 2021 & Scale anchor & \sourcetext{Genome-scale/omics phenotype} & -- & \href{https://doi.org/10.1128/aem.00041-21}{10.\allowbreak 1128/\allowbreak aem.\allowbreak 00041-\allowbreak 21} \\
\addlinespace[2pt]
85 & \org{P. putida} & \sourcetext{Understanding butanol tolerance and assimilation in Pseudomonas putida BIRD-1: an integrated omics approach.} & 2016 & Scale anchor & \sourcetext{Genome-scale/omics phenotype} & -- & \href{https://doi.org/10.1111/1751-7915.12328}{10.\allowbreak 1111/\allowbreak 1751-\allowbreak 7915.\allowbreak 12328} \\
\addlinespace[2pt]
86 & \org{P. putida} & \sourcetext{Novel Toxin-Antitoxin Module SlvT-SlvA Regulates Megaplasmid Stability and Incites Solvent Tolerance in Pseudomonas putida S12.} & 2020 & Scale anchor & \sourcetext{Genome-scale/omics phenotype} & -- & \href{https://doi.org/10.1128/aem.00686-20}{10.\allowbreak 1128/\allowbreak aem.\allowbreak 00686-\allowbreak 20} \\
\addlinespace[2pt]
87 & \org{P. putida} & \sourcetext{Phospholipid biosynthesis and solvent tolerance in Pseudomonas putida strains.} & 1997 & Scale anchor & \sourcetext{Genome-scale/omics phenotype} & -- & \href{https://doi.org/10.1128/jb.179.13.4219-4226.1997}{10.\allowbreak 1128/\allowbreak jb.\allowbreak 179.\allowbreak 13.\allowbreak 4219-\allowbreak 4226.\allowbreak 1997} \\
\addlinespace[2pt]
88 & \org{P. putida} & \sourcetext{Engineering thermal stability and solvent tolerance of the soluble quinoprotein PedE from Pseudomonas putida KT2440 with a heterologous whole-cell screening approach.} & 2018 & Scale anchor & \sourcetext{Genome-scale/omics phenotype} & -- & \href{https://doi.org/10.1111/1751-7915.13036}{10.\allowbreak 1111/\allowbreak 1751-\allowbreak 7915.\allowbreak 13036} \\
\addlinespace[2pt]
89 & \org{P. putida} & \sourcetext{Regulation of solvent tolerance in Pseudomonas putida S12 mediated by mobile elements.} & 2017 & Scale anchor & \sourcetext{Genome-scale/omics phenotype} & -- & \href{https://doi.org/10.1111/1751-7915.12495}{10.\allowbreak 1111/\allowbreak 1751-\allowbreak 7915.\allowbreak 12495} \\
\addlinespace[2pt]
90 & \org{P. putida} & \sourcetext{Dynamic Response of Pseudomonas putida S12 to Sudden Addition of Toluene and the Potential Role of the Solvent Tolerance Gene trgI.} & 2015 & Scale anchor & \sourcetext{Genome-scale/omics phenotype} & -- & \href{https://doi.org/10.1371/journal.pone.0132416}{10.\allowbreak 1371/\allowbreak journal.\allowbreak pone.\allowbreak 0132416} \\
\addlinespace[2pt]
91 & \org{P. putida} & \sourcetext{Identification of conditionally essential genes for growth of Pseudomonas putida KT2440 on minimal medium through the screening of a genome-wide mutant library.} & 2010 & Scale anchor & \sourcetext{Genome-scale/omics phenotype} & -- & \href{https://doi.org/10.1111/j.1462-2920.2010.02166.x}{10.\allowbreak 1111/\allowbreak j.\allowbreak 1462-\allowbreak 2920.\allowbreak 2010.\allowbreak 02166.\allowbreak x} \\
\addlinespace[2pt]
92 & \org{P. putida} & \sourcetext{Proteomic analysis of Pseudomonas putida reveals an organic solvent tolerance-related gene mmsB.} & 2013 & Scale anchor & \sourcetext{Genome-scale/omics phenotype} & -- & \href{https://doi.org/10.1371/journal.pone.0055858}{10.\allowbreak 1371/\allowbreak journal.\allowbreak pone.\allowbreak 0055858} \\
\addlinespace[2pt]
93 & \org{P. putida} & \sourcetext{Transcriptional phase variation at the flhB gene of Pseudomonas putida DOT-T1E is involved in response to environmental changes and suggests the participation of the flagellar export system in solvent tolerance.} & 2004 & Scale anchor & \sourcetext{Genome-scale/omics phenotype} & -- & \href{https://doi.org/10.1128/jb.186.6.1905-1909.2004}{10.\allowbreak 1128/\allowbreak jb.\allowbreak 186.\allowbreak 6.\allowbreak 1905-\allowbreak 1909.\allowbreak 2004} \\
\addlinespace[2pt]
94 & \org{P. putida} & \sourcetext{An antirepressor, SrpR, is involved in transcriptional regulation of the SrpABC solvent tolerance efflux pump of Pseudomonas putida S12.} & 2011 & Scale anchor & \sourcetext{Genome-scale/omics phenotype} & -- & \href{https://doi.org/10.1128/jb.00149-11}{10.\allowbreak 1128/\allowbreak jb.\allowbreak 00149-\allowbreak 11} \\
\addlinespace[2pt]
95 & \org{P. putida} & \sourcetext{Adaptive laboratory evolution and genetic engineering improved terephthalate utilization in Pseudomonas putida KT2440.} & 2025 & Scale anchor & \sourcetext{Genome-scale/omics phenotype} & -- & \href{https://doi.org/10.1016/j.ymben.2024.12.006}{10.\allowbreak 1016/\allowbreak j.\allowbreak ymben.\allowbreak 2024.\allowbreak 12.\allowbreak 006} \\
\addlinespace[2pt]
97 & \org{P. putida} & \sourcetext{Adaptive laboratory evolution rewires  Pseudomonas putida  for resource-efficient acetate assimilation} & 2026 & Scale anchor & \sourcetext{Genome-scale/omics phenotype} & -- & \href{https://doi.org/10.64898/2026.08.20.746122}{10.\allowbreak 64898/\allowbreak 2026.\allowbreak 08.\allowbreak 20.\allowbreak 746122} \\
\addlinespace[2pt]
98 & \org{P. putida} & \sourcetext{Automated mini-bioreactors reveal the temporal dynamics and multi-omics responses of CRISPRi knockdowns in  Pseudomonas putida} & 2026 & Scale anchor & \sourcetext{Genome-scale/omics phenotype} & -- & \href{https://doi.org/10.64898/2026.03.06.709552}{10.\allowbreak 64898/\allowbreak 2026.\allowbreak 03.\allowbreak 06.\allowbreak 709552} \\
\addlinespace[2pt]
99 & \org{P. putida} & \sourcetext{Adaptive laboratory evolution unlocks membrane permeability as a key limitation in long-chain alcohol metabolism by  Pseudomonas putida  KT2440} & 2026 & Scale anchor & \sourcetext{Genome-scale/omics phenotype} & -- & \href{https://doi.org/10.64898/2026.01.19.700371}{10.\allowbreak 64898/\allowbreak 2026.\allowbreak 01.\allowbreak 19.\allowbreak 700371} \\
\addlinespace[2pt]
100 & \org{P. putida} & \sourcetext{Functional analysis of the fatty acid and alcohol metabolism of  Pseudomonas putida  using RB-TnSeq} & 2020 & Scale anchor & \sourcetext{Genome-scale/omics phenotype} & -- & \href{https://doi.org/10.1101/2020.07.04.188060}{10.\allowbreak 1101/\allowbreak 2020.\allowbreak 07.\allowbreak 04.\allowbreak 188060} \\
\addlinespace[2pt]
\end{longtable}
\endgroup
\end{landscape}


%% notes-tex/database/database-expansion-bacteria/sections/analogs.tex
%% [[experiments.database.expansion-bacteria]]
%% https://github.com/Mjvolk3/torchcell/tree/main/notes-tex/database/database-expansion-bacteria/sections/analogs.tex
%% SOURCE: analogs table from
%% experiments/database/scripts/build_bacteria_candidate_datasets_table.py;
%% surrounding prose is hand-authored

\section{What each row maps onto in the supported set\secstatus{todo}}\label{sec:analogs}

A bacterial row is cheap or expensive depending on whether the substrate already holds its
record type. Table~\ref{tab:banalogs} names, for every row, the built \org{S. cerevisiae}
dataset it mirrors and why. A named analog means the loader writes a record the schema
already validates, against a different reference genome, so the work is retrieval and
provenance. A row with no analog is naming a phenotype the substrate has never held.

%% GENERATED FILE -- do not hand-edit.
%% SOURCE: experiments/database/scripts/build_bacteria_candidate_datasets_table.py
\begin{landscape}
\begingroup
\footnotesize
\setlength{\tabcolsep}{4pt}
\renewcommand{\arraystretch}{1.15}
\begin{longtable}{@{}r@{\hspace{4pt}} L{40mm} L{40mm} L{74mm} L{74mm}@{}}
\caption[]{What each of the fifty recommended builds, rows 1--50 of
Table~\ref{tab:bfinal}, maps onto in the supported set; the ranked reserve below the cut is
not listed. A named analog means
the loader writes a record type the schema already holds, against a different reference
genome, so the work is the retrieval and the provenance rather than a new phenotype
class. \emph{What the schema still needs} is per-row and excludes the two blockers every
row shares (Table~\ref{tab:bschema}); a dash means those two are the whole cost.}
\label{tab:banalogs}\\
\toprule
\textbf{\#} & \textbf{Bacterial row} & \textbf{Built yeast analog} & \textbf{Why it is the analog} & \textbf{What the schema still needs} \\
\midrule
\endfirsthead
\multicolumn{5}{@{}l}{\footnotesize\emph{Table~\ref{tab:banalogs}, continued}}\\
\toprule
\textbf{\#} & \textbf{Bacterial row} & \textbf{Built yeast analog} & \textbf{Why it is the analog} & \textbf{What the schema still needs} \\
\midrule
\endhead
\bottomrule
\endfoot
1 & Fuhrer 2017 \org{E. coli} & \textbf{Mulleder 2016 amino-acid metabolome} & whole deletion collection metabolite profiling & -- \\
\addlinespace[5pt]
2 & Wetmore 2015 \org{E. coli} & \textbf{Hillenmeyer 2008 HIP/HOP} & pooled barcoded competitive fitness assay & -- \\
\addlinespace[5pt]
3 & Mutalik 2020 \org{E. coli} & \textbf{Hillenmeyer 2008 HIP/HOP} & pooled barcoded mutants across selective conditions & -- \\
\addlinespace[5pt]
4 & Tong 2020 \org{E. coli} & \textbf{Smith 2006 chemogenomic} & colony growth per deletion per condition & -- \\
\addlinespace[5pt]
5 & PRECISE-1K \org{E. coli} & \textbf{Kemmeren 2014 deletion transcriptome} & deletion-and-condition expression compendium & -- \\
\addlinespace[5pt]
6 & Carruthers 2025 \org{P. putida} & \textbf{Lian 2019 CRISPR-AID} & multiplexed guide combinations scored by product & -- \\
\addlinespace[5pt]
7 & Lim 2022 putidaPRECISE321 \org{P. putida} & \textbf{Caudal 2024 pan-transcriptome} & aggregated RNA-seq compendium across many backgrounds & -- \\
\addlinespace[5pt]
8 & Borchert 2024 fModules \org{P. putida} & \textbf{Hillenmeyer 2008 HIP/HOP} & pooled mutant fitness across hundreds of conditions & -- \\
\addlinespace[5pt]
9 & Caglar 2017 \org{E. coli} & \textbf{Zelezniak 2018 proteome and metabolome} & two layers on identical samples & -- \\
\addlinespace[5pt]
10 & Nichols 2011 \org{E. coli} & \textbf{Hillenmeyer 2008 HIP/HOP} & deletion collection across many chemicals & -- \\
\addlinespace[5pt]
11 & Goodall 2018 \org{E. coli} & \textbf{SGD essentiality} & genome-wide essential gene call set & -- \\
\addlinespace[5pt]
12 & Foo 2014 isopentenol tolerance \org{E. coli} & \textbf{Lopez 2024 isobutanol} & alcohol tolerance genes scored on a production chassis & a transcriptome measured against a dosed inhibitor on an unperturbed genotype, which is an environment-response record rather than a gene-perturbation record \\
\addlinespace[5pt]
13 & Yunus 2026 \org{P. putida} & \textbf{Lian 2019 CRISPR-AID} & predicted multiplexed knockdowns scored by titer & -- \\
\addlinespace[5pt]
14 & de Siqueira 2025 \org{P. putida} & \textbf{Mormino 2022 CRISPRi acetic acid} & acetate tolerance with a product readout & -- \\
\addlinespace[5pt]
15 & Wang 2015 isoprenol tolerance \org{E. coli} & \textbf{Lopez 2024 isobutanol} & deletion collection scored for an alcohol phenotype & -- \\
\addlinespace[5pt]
16 & Lim 2025 isoprenol TALE \org{P. putida} & \textbf{Mormino 2022 CRISPRi acetic acid} & inhibitor tolerance scored by growth & -- \\
\addlinespace[5pt]
17 & Tian 2019 isopentenol CRISPRi \org{E. coli} & \textbf{Lian 2019 CRISPR-AID} & multiplexed guide arrays over one gene set & a guide-array perturbation carrying more than one target at a graded induction level, and a titer measured on a chassis whose pathway is plasmid-borne rather than integrated \\
\addlinespace[5pt]
18 & Kang 2026 isoprenyl acetate \org{P. putida} & \textbf{Lopez 2024 isobutanol} & engineered strain series producing an alcohol ester & -- \\
\addlinespace[5pt]
19 & Wang 2022 P. putida isoprenoids \org{P. putida} & \textbf{Ozaydin 2013 beta-carotene} & heterologous isoprenoid pathway plus host deletions & a titer for several products measured on one genotype, so the phenotype is a vector over product identity rather than one label \\
\addlinespace[5pt]
20 & Menasalvas 2025 \org{P. putida} & \textbf{Lopez 2024 isobutanol} & iterative engineered panel, alcohol titer & -- \\
\addlinespace[5pt]
21 & Price 2018 \org{E. coli} & \textbf{Hillenmeyer 2008 HIP/HOP} & pooled barcoded fitness across many conditions & -- \\
\addlinespace[5pt]
22 & Choe 2025 \org{E. coli} & \textbf{Smith 2016 CRISPRi} & pooled CRISPRi knockdown across drugs & -- \\
\addlinespace[5pt]
23 & Rapp 2026 CRISPRi metabolome \org{E. coli} & \textbf{Zelezniak 2018 metabolome} & both put metabolite concentrations and a second molecular layer on one strain axis & -- \\
\addlinespace[5pt]
24 & Schmidt 2022 nitrogen \org{P. putida} & \textbf{Hillenmeyer 2008 HIP/HOP} & pooled mutant fitness across nitrogen sources & -- \\
\addlinespace[5pt]
25 & Wang 2018 \org{E. coli} & \textbf{Smith 2016 CRISPRi} & pooled CRISPRi guide fitness screen & -- \\
\addlinespace[5pt]
26 & Rousset 2018 \org{E. coli} & \textbf{Smith 2016 CRISPRi} & pooled CRISPRi guide fitness screen & -- \\
\addlinespace[5pt]
27 & Shiver 2016 \org{E. coli} & \textbf{Wildenhain 2015 drug tolerance} & neglected-compound panel against deletion collection & -- \\
\addlinespace[5pt]
28 & Thompson 2020 fatty acid and alcohol \org{P. putida} & \textbf{Hillenmeyer 2008 HIP/HOP} & pooled mutant fitness across a carbon-source panel & -- \\
\addlinespace[5pt]
29 & Borchert 2023 lignin tolerance \org{P. putida} & \textbf{Hillenmeyer 2008 HIP/HOP} & pooled mutant fitness across chemical stresses & -- \\
\addlinespace[5pt]
30 & Schastnaya 2021 \org{E. coli} & \textbf{Mulleder 2016 amino-acid metabolome} & point-mutant panel, metabolite readout & -- \\
\addlinespace[5pt]
31 & Cui 2018 \org{E. coli} & \textbf{Smith 2016 CRISPRi} & pooled CRISPRi guide fitness screen & -- \\
\addlinespace[5pt]
32 & Schmidt 2016 \org{E. coli} & \textbf{Messner 2023 proteome} & both give panel-wide absolute proteome depth & -- \\
\addlinespace[5pt]
33 & Butland 2008 \org{E. coli} & \textbf{Costanzo 2016 SGA} & conjugation analog of yeast SGA & -- \\
\addlinespace[5pt]
34 & Mori 2021 \org{E. coli} & \textbf{Messner 2023 proteome} & absolute proteome across a large condition panel & -- \\
\addlinespace[5pt]
35 & Ishii 2007 \org{E. coli} & \textbf{Zelezniak 2018 proteome and metabolome} & several layers, one strain panel & -- \\
\addlinespace[5pt]
36 & Campos 2018 \org{E. coli} & \textbf{Kemmeren 2014 deletion transcriptome} & high-dimensional profile per deletion strain & -- \\
\addlinespace[5pt]
37 & Typas 2008 \org{E. coli} & \textbf{Costanzo 2016 SGA} & conjugation analog of yeast SGA & -- \\
\addlinespace[5pt]
38 & Girgis 2009 \org{E. coli} & \textbf{Hillenmeyer 2008 HIP/HOP} & pooled mutant library across drug panel & -- \\
\addlinespace[5pt]
39 & Banerjee 2025 \org{P. putida} & \textbf{Zelezniak 2018 proteome} & engineered panel paired with quantitative proteomes & -- \\
\addlinespace[5pt]
40 & Rachwalski 2024 mobile CRISPRi \org{E. coli} & \textbf{Costanzo 2016 SGA} & pairwise gene-gene fitness across a cataloged array & a genotype combining a cataloged deletion with a guide-borne knockdown, so one perturbation is an edit to the genome and the other is not \\
\addlinespace[5pt]
41 & Fang 2025 FFA CRISPRi-FACS \org{E. coli} & \textbf{Xue 2025 free fatty acids} & both read free fatty acid output as the label on a perturbed strain & -- \\
\addlinespace[5pt]
42 & Babu 2014 \org{E. coli} & \textbf{Costanzo 2016 SGA} & conjugation digenic epistasis colony scores & -- \\
\addlinespace[5pt]
43 & Gupta 2024 turnover \org{E. coli} & \textbf{Messner 2023 proteome} & proteome-wide per-protein quantitative depth & -- \\
\addlinespace[5pt]
44 & MCF2Chem 2023 \org{E. coli} & \textbf{Ozaydin 2013 beta-carotene screen} & both score a small-molecule product per engineered strain over a genome-scale perturbation panel & -- \\
\addlinespace[5pt]
45 & Choe 2019 genome-reduced ALE \org{E. coli} & \textbf{Caudal 2024 pan-transcriptome} & expression across genetically distinct backgrounds & a genotype that is a population allele-frequency vector rather than one strain, and a reference that is itself a reduced derivative with no deposited assembly \\
\addlinespace[5pt]
46 & Hawkins 2020 mismatch-CRISPRi \org{E. coli} & \textbf{Smith 2016 CRISPRi} & essential-gene knockdown reachable only by interference & a graded knockdown whose level is a model prediction rather than a measured quantity, so the dose carries its own uncertainty \\
\addlinespace[5pt]
47 & Royet 2025 KT2440 metal Tn-seq \org{P. putida} & \textbf{Hillenmeyer 2008 HIP/HOP} & genome-wide fitness across dosed chemical stress & a fitness record whose genotype is a gene rather than a strain, because a non-barcoded insertion pool is never resolvable to a clone \\
\addlinespace[5pt]
48 & Niu 2019 pinene evolved \org{E. coli} & \textbf{Lian 2019 CRISPR-AID} & activation and interference over one gene set & an evolved genotype released as a variant list against a reference that is not its own parent, so the call set needs a provenance gap rather than a clean edit list \\
\addlinespace[5pt]
49 & Yang 2019 mevalonate from ethanol \org{P. putida} & \textbf{Ozaydin 2013 beta-carotene} & stacked host deletions feeding a heterologous isoprenoid route & a genotype stacking five markerless deletions with four heterologous genes, where the deletion boundary is given by primer rather than by coordinate \\
\addlinespace[5pt]
50 & Li 2021 mevalonate flux \org{E. coli} & \textbf{Zelezniak 2018 metabolome} & internal flux state under targeted perturbation & a fitted flux distribution as the phenotype, where each value carries an interval from the fit rather than a replicate standard deviation \\
\addlinespace[5pt]
\end{longtable}
\endgroup
\end{landscape}


%% notes-tex/database/database-expansion-bacteria/sections/caveats.tex
%% [[experiments.database.expansion-bacteria]]
%% https://github.com/Mjvolk3/torchcell/tree/main/notes-tex/database/database-expansion-bacteria/sections/caveats.tex
%% SOURCE: excluded table from
%% experiments/database/scripts/build_bacteria_candidate_datasets_table.py; prose
%% hand-authored

\section{Before a loader is written\secstatus{todo}}\label{sec:caveats}

Every figure in Table~\ref{tab:bfinal} is an axis count rather than a verified record
count, and the qualifications below decide which rows can be trusted at their stated size.

\notesec{Instance counts marked $\dagger$ are products, not counts} A screen reported as
$n$ strains against $k$ conditions enters as $nk$ instances. Real screens drop strains per
condition, so each such figure is an upper bound. The ones marked $\ddagger$ are weaker
still, being order-of-magnitude estimates for libraries whose exact size was not confirmed.
Neither should be quoted as a dataset size without opening the source.

\notesec{Dimensionality is the other half of every rank, so an estimated depth moves the
order} Because rows are ranked on instances times dimensionality, a wrong panel depth moves
a row directly rather than at the margin. The generated table records whether each depth was
read from a source or estimated, and the summary counts the split. A row whose position
looks surprising should have its depth checked before its instance count.

\notesec{A corpus is not a dataset} Rows marked \textbf{A} are curated literature
collections. Their record counts are counts of reported observations extracted from other
papers, so they are not net new until the corpus is resolved to primary sources, and
several of those sources are already elsewhere in this table. They are ranked because
their aggregate size is real and useful for discovery, and they are flagged because
counting them as independent datasets would double-count.

\notesec{Rows marked \textbf{B} have no released per-record values} A screen whose released
form is a figure or a table of selected hits is not ingestible as a genotype-to-phenotype
dataset, whatever its reputation. These stay in the table with the flag rather than being
dropped, because an author request or a supplementary re-check can change the state.

\subsection{The transposon rows overlap each other, and the largest is a superset\secstatus{todo}}\label{sec:overlap}

Several of the highest-ranking rows draw on the same barcoded libraries, so ingesting them
in rank order would load the same measurements more than once. This is the one place where
the ranking and the build order should come apart.

For \org{P. putida}, one released fitness matrix is dense and machine-readable at 4,732
genes by 332 samples with no missing values, and it is a superset of the other
transposon-fitness rows in this table. Those rows should be de-duplicated against it
rather than ingested in parallel, which turns four loaders into one plus a provenance
record of which source experiments the matrix covers. The library behind them is one
library, not several: the reported insertion counts belong to a single barcoded collection
of 185,401 insertions, 152,810 of them in genes covering 5,213 of 5,661 genes, and a
smaller figure quoted elsewhere is fitness coverage rather than library size.

For \org{E. coli}, the same relation holds between the method paper that introduced the
barcoded library and the later compendium built on it: the compendium is a superset of that
paper's \org{E. coli} experiments. The deletion-collection rows have the cleaner
relationship, since they share a strain axis without sharing measurements, and that shared
axis is what makes them joinable rather than redundant.

\subsection{Access findings that change what a loader can do\secstatus{todo}}\label{sec:access}

Four findings are specific enough to name, because each one changes the work rather than
the confidence.

\notesec{One \org{P. putida} campaign is deposited as seven separate proteomics
accessions, one per design cycle} Ingesting the one accession a reader finds first would
drop six sevenths of the proteome. The campaign is one dataset with linked assets, and its
loader has to enumerate all seven.

\notesec{No \org{P. putida} knockdown library in this table releases a guide-by-sample
matrix} The library itself is real, roughly three guides per gene over 5,591 coding
sequences, but the deposit holds baseline per-guide counts, a lost-guide list, per-round
enriched hit lists and representative rather than per-replicate sequencing. A per-guide
fitness matrix cannot be reconstructed from it, so the row enters on its per-strain
production measurements and not as a pooled screen.

\notesec{Two accessions did not resolve when checked} One resequencing project identifier
returns nothing at the sequence archives, so the mutation calls exist only in a
supplementary table, and one deletion-collection portal has moved host. Neither is a
reason to drop a row, and both are a reason the provenance record has to pin the file that
was actually read.

\notesec{One fitness portal is behind a challenge that blocks scripted retrieval} Four
rows whose only data location is that portal are therefore unconfirmed at the accession
level, and the coverage figures circulating for it come from archived snapshots rather than
a live read. They are recorded as unconfirmed rather than quoted as current.

\section{What was dropped\secstatus{todo}}\label{sec:excluded}

%% GENERATED FILE -- do not hand-edit.
%% SOURCE: experiments/database/scripts/build_bacteria_candidate_datasets_table.py
\begingroup
\footnotesize
\begin{longtable}{@{}L{58mm} L{28mm} L{83mm}@{}}
\caption[]{Considered and dropped. \emph{no-sequence} is the hard gate: without a route
to the strain's genomic content there is no genotype to map a phenotype from.
\emph{no-per-record-data} is a real experiment whose released form is a figure or a
summary rather than per-strain values.}
\label{tab:bexcluded}\\
\toprule
Dataset or group & Rule & Reason \\
\midrule
\endfirsthead
\toprule
Dataset or group & Rule & Reason \\
\midrule
\endhead
\bottomrule
\endfoot
Mi 2014, geranic acid in P. putida DSM 12264 & no-sequence & The host is DSM 12264, a solvent-tolerant strain that is not KT2440, and no genome accession for it appears anywhere in the paper. PP\_ locus tags therefore do not apply and the chassis genome is unidentified, so there is no route to the total genomic content of any strain. The enzyme responsible for the key endogenous geraniol oxidation is also unidentified, stated by the authors as a hypothesis. Two comparisons made in KT2440, the host actually of interest, are reported as data not shown \\
\addlinespace[5pt]
Chen 2013, lycopene by chromosomal evolution in E. coli & no-sequence & The perturbation is triclosan-driven tandem gene amplification, and its extent is quantified only as about 30 copies by quantitative PCR of a single gene against one reference gene. The amplicon boundaries were never sequenced or mapped and no strain was resequenced, so an array of unstated extent cannot be written as a genome sequence, and mutations hitchhiking through six rounds of escalating selection are unknown. The released phenotype tables are good, with mean and SD for 20 strains, but the perturbation axis itself, triclosan dose against copy number, is figure-only \\
\addlinespace[5pt]
Glebes 2014, SCALEs furfural dosage map in E. coli & no-per-record-data & The headline is the largest in the whole queue and the release is the smallest part of it. More than ten to the fifth genomic fragments were assayed by microarray at about 125 nucleotide resolution across the genome, under furfural at 0.75 grams per liter against a control, and the only table deposited is Table S2, 268 enriched genes with one fitness score each. The microarray data are not in GEO or anywhere else named in the paper, so the per-clone frequencies the method is built on are gone, and the surviving 268 rows are the enriched tail with every neutral and depleted fragment discarded. Ranking this by its abstract would have placed it near the top of the list on a measurement count that no longer exists \\
\addlinespace[5pt]
Casanova-Hampton 2021, copper stress across the Keio collection & no-per-record-data & The whole Keio collection was stamped onto LB agar at 0, 3 and 6 millimolar copper sulfate, which would be about twelve thousand mutant by dose observations, but the paper states that phenotypic response was determined qualitatively relative to on-plate controls at 24, 48 and 72 hours. There is no colony size, no growth index and no score of any kind, and what is released is the call: 43 sensitive and 25 resistant mutants in Tables S4 and S5, plus a 36-mutant targeted rescreen of iron genes. A qualitative call carries no uncertainty and cannot be compared against the quantitative colony scores the store already holds for the same collection \\
\addlinespace[5pt]
Mei 2026, ciprofloxacin tolerance across the Keio collection & no-per-record-data & A two-step genome-wide screen of the Keio library, stationary-phase starvation then dilution into antibiotic, and the data availability statement says all data are in the Supporting Information, which is a single PDF. Nothing machine-readable is deposited, and what the PDF carries is the outcome, 37 ciprofloxacin-sensitive and 11 tolerant mutants out of the roughly four thousand screened. The 48 named mutants are a hit list, and the survival values behind the other mutants are what would have made it a dataset \\
\addlinespace[5pt]
Huang 2022, pinene by shotgun library in E. coli MG1655 & no-per-record-data & The data availability statement reads Not applicable, there is no results table of any kind, and only two absolute titers plus three percentages are recoverable from prose. The library screen also did not work as a screen: a library of more than ten to the fifth transformants yielded four surviving colonies, which the authors attribute to loss during construction or pinene toxicity, so there is no genotype-phenotype matrix to ingest. Insert boundaries for the four clones and insertion coordinates for the three promoter replacements are both unstated \\
\addlinespace[5pt]
Mireles 2026, long-chain alcohol metabolism in P. putida & no-per-record-data & Sequencing data are available from the corresponding author upon reasonable request, with no deposit and no accession, and the evolved-clone phenotype is figure-only, so genotype and phenotype cannot both be taken as released values. The released variant table also disagrees with the text, 84 per-clone calls against 91 claimed, and labels one mutation PP\_1802 where the text argues from PP\_1801. Separately it is off this table's phenotype axis: the alcohols are C16 and C20 used as sole carbon sources and the paper states that toxicity is absent, so it measures catabolism rather than tolerance \\
\addlinespace[5pt]
Hernandez-Arranz 2019, isoprenoid precursors in P. putida & no-per-record-data & Lycopene is reported in two incompatible ways that cannot be joined. The one table of measurements gives nanograms per milliliter by chromatography for four strains, all of them the lowest producers, while every other result is an absorbance normalized to cell density and expressed as a fold change against a control. The headline fifty-fold strain has no absolute value, so ten of fourteen genotypes carry only a ratio. There is no strain table, no supplementary file, and no accession of any kind, and the perturbed native genes are never given a locus tag \\
\addlinespace[5pt]
Multi-host cell-factory capacity resource (modeled yields) & not-a-dataset & Genome-scale model predictions of production capacity across host and product combinations. Every value is computed rather than measured, so it belongs in a model-derived partition and must not enter an experimental instance or dataset count. Useful for negative sampling and as a feature source, which is a different role from training data \\
\addlinespace[5pt]
Simulated FSEOF rows inside the curated engineering corpora & not-a-dataset & The larger metabolic-engineering corpora carry both experimentally reported modifications and rows generated by flux-scanning enforced objective flux. The simulated rows are predictions of what an intervention would do and are excluded or separately typed; mixing them into an experimental count is the specific error these corpora invite \\
\addlinespace[5pt]
Adaptive laboratory evolution without resequencing & no-sequence & An evolved population differs from its parent at unmapped positions, so there is no genotype to map a phenotype from. Only resequenced clones and explicitly reconstructed mutations from such a campaign are admitted, which is why the tolerance-evolution rows enter on their sequenced isolates rather than on their lineages \\
\addlinespace[5pt]
Genome shuffling, chemical and UV mutagenesis strain panels & no-sequence & Same gate as unsequenced evolution, and it bites harder: the selected isolate carries an unknown number of edits at unknown positions. No amount of phenotypic value substitutes for a missing genotype \\
\addlinespace[5pt]
Other bacterial hosts (B. subtilis, C. glutamicum, Z. mobilis, cyanobacteria) & off-host & Real and in several cases large, but outside the two hosts this list covers. The generalization axis is single-cell bioproduction hosts, so these become candidates once either bacterial host is actually served and the per-host schema and genome-tier cost has been paid once \\
\addlinespace[5pt]
EcoCyc and the curated annotation databases & not-a-dataset & EcoCyc releases per-class flat files covering 4,688 genes, 479 pathways and 6,124 regulatory interactions for the reference strain (BioCyc v30.0, checked live). Those are annotations of one genome rather than measured genotype-by-environment records, so under a ranking on instances times dimensionality the curation would outrank real screens. It belongs in the reference layer the graph already has, alongside the gene ontology, not in a candidate list of experiments \\
\addlinespace[5pt]
Review and perspective tables of reported titers & not-a-dataset & A titer quoted in a review is a second-hand number without the strain construction, the medium or the replicate structure that makes it a record. The primary paper is the candidate; the review is a discovery index, and the curated corpora already serve that role better \\
\addlinespace[5pt]
\end{longtable}
\endgroup


%% notes-tex/database/database-expansion-bacteria/sections/summary.tex
%% [[experiments.database.expansion-bacteria]]
%% https://github.com/Mjvolk3/torchcell/tree/main/notes-tex/database/database-expansion-bacteria/sections/summary.tex
%% SOURCE: experiments/database/scripts/build_bacteria_candidate_datasets_table.py

\clearpage
\section{Summary statistics\secstatus{todo}}\label{sec:summary}

%% GENERATED FILE -- do not hand-edit.
%% SOURCE: experiments/database/scripts/build_bacteria_candidate_datasets_table.py
\begin{table}[H]\centering
\small
\caption[]{Summary of the fifty, split at the tranche-1 line. \emph{Tr. 1} and
\emph{Tr. 2} count rows; \emph{Genotypes}, \emph{Instances} and \emph{Meas.} sum the row
axes over all fifty and are lower bounds, since a row with no count contributes nothing.
The fifty name 61 joins, 48 of them to a dataset already built.}
\label{tab:bsummary}
\begin{tabular}{@{}l r r r r r@{}}
\toprule
 & Tr. 1 & Tr. 2 & Genotypes & Instances & Meas. \\
\midrule
\multicolumn{6}{@{}l}{\emph{By host}}\\
E. coli & 11 & 25 & 3,137,690 & $3.0\times 10^{7}$ & $3.0\times 10^{8}$ \\
P. putida & 9 & 5 & 19,813 & $2.3\times 10^{6}$ & $6.3\times 10^{6}$ \\
\midrule
\multicolumn{6}{@{}l}{\emph{By tier}}\\
tier 1 & 7 & 12 & 2,839,383 & $3.1\times 10^{7}$ & $3.0\times 10^{8}$ \\
tier 2 & 2 & 6 & 306,616 & $1.2\times 10^{6}$ & $1.3\times 10^{6}$ \\
tier 3 & 3 & 4 & 11,120 & $1.4\times 10^{4}$ & $7.1\times 10^{6}$ \\
tier 4 & 8 & 8 & 384 & $2.4\times 10^{3}$ & $3.9\times 10^{6}$ \\
\midrule
\multicolumn{6}{@{}l}{\emph{By class}}\\
Aggregation / support & 0 & 1 & 8,888 & $8.9\times 10^{3}$ & $3.6\times 10^{4}$ \\
CRISPR library screen & 0 & 6 & 303,848 & $1.2\times 10^{6}$ & $1.7\times 10^{6}$ \\
Combinatorial design & 1 & 1 & 104 & 161 & 322 \\
Fitness / chemical genomics & 2 & 2 & 15,977 & $1.6\times 10^{6}$ & $9.8\times 10^{6}$ \\
Genetic interaction & 0 & 4 & 1,089,977 & $5.1\times 10^{5}$ & $5.1\times 10^{5}$ \\
Metabolome / flux & 1 & 2 & 3,899 & $3.4\times 10^{4}$ & $2.6\times 10^{8}$ \\
Multi-omics campaign & 1 & 2 & 29 & 361 & $1.5\times 10^{6}$ \\
Production campaign & 6 & 2 & 732 & $3.5\times 10^{3}$ & $2.2\times 10^{6}$ \\
Proteome & 0 & 2 & 6 & 132 & $3.1\times 10^{5}$ \\
Tolerance / robustness & 3 & 0 & 100 & 271 & 271 \\
Transcriptome & 2 & 1 & 2,023 & $9.1\times 10^{3}$ & $6.3\times 10^{6}$ \\
Translation / turnover & 0 & 1 & 1 & 13 & $4.2\times 10^{4}$ \\
Transposon fitness & 4 & 6 & 1,731,919 & $2.9\times 10^{7}$ & $2.9\times 10^{7}$ \\
\midrule
\multicolumn{6}{@{}l}{\emph{By sequence basis}}\\
K-12+guide & 0 & 6 & 303,848 & $1.2\times 10^{6}$ & $1.7\times 10^{6}$ \\
K-12+promoter-reporter & 0 & 1 & 1,930 & $7.7\times 10^{3}$ & $6.9\times 10^{4}$ \\
K-12+transposon & 3 & 3 & 1,712,899 & $2.7\times 10^{7}$ & $2.7\times 10^{7}$ \\
K-12-KO & 5 & 4 & 19,935 & $1.7\times 10^{6}$ & $2.7\times 10^{8}$ \\
K-12-KO x KO & 0 & 3 & 1,077,573 & $4.5\times 10^{5}$ & $4.5\times 10^{5}$ \\
K-12-KO+guide & 0 & 1 & 12,404 & $6.4\times 10^{4}$ & $6.4\times 10^{4}$ \\
KT2440+guide & 2 & 0 & 492 & $1.5\times 10^{3}$ & $2.1\times 10^{6}$ \\
KT2440+promoter & 0 & 1 & 12 & 36 & $7.2\times 10^{4}$ \\
KT2440+transposon & 1 & 3 & 19,020 & $2.3\times 10^{6}$ & $2.3\times 10^{6}$ \\
engineered-chassis & 4 & 2 & 317 & $2.4\times 10^{3}$ & $2.1\times 10^{5}$ \\
engineered-chassis+guide & 1 & 0 & 24 & 84 & 168 \\
evolved-WGS & 2 & 2 & 136 & 264 & $3.5\times 10^{4}$ \\
reference-only & 2 & 4 & 8,913 & $9.7\times 10^{3}$ & $3.5\times 10^{6}$ \\
\midrule
\multicolumn{6}{@{}l}{\emph{Other attributes}}\\
carries a time axis & 9 & 5 & 6,048 & $1.2\times 10^{5}$ & $9.8\times 10^{6}$ \\
mirrors a built yeast dataset & 20 & 30 & 3,157,503 & $3.3\times 10^{7}$ & $3.1\times 10^{8}$ \\
has a join to a built dataset & 20 & 28 & 3,151,882 & $3.3\times 10^{7}$ & $3.1\times 10^{8}$ \\
figures sourced this pass & 20 & 30 & 3,157,503 & $3.3\times 10^{7}$ & $3.1\times 10^{8}$ \\
data location confirmed live & 14 & 21 & 1,138,156 & $5.3\times 10^{6}$ & $2.8\times 10^{8}$ \\
per-record values not released & 2 & 3 & 147,379 & $1.5\times 10^{5}$ & $1.5\times 10^{5}$ \\
a corpus, not net new until split & 2 & 1 & 13,637 & $1.6\times 10^{6}$ & $3.4\times 10^{6}$ \\
instances a reported count & 7 & 12 & 1,080,639 & $2.3\times 10^{6}$ & $2.7\times 10^{8}$ \\
instances a product of the axes & 11 & 15 & 1,576,806 & $3.0\times 10^{7}$ & $3.3\times 10^{7}$ \\
instances an estimate & 2 & 3 & 500,058 & $7.3\times 10^{4}$ & $1.9\times 10^{5}$ \\
dimensionality reported, not estimated & 16 & 27 & 3,096,444 & $3.2\times 10^{7}$ & $3.0\times 10^{8}$ \\
\midrule
Total & 20 & 30 & 3,157,503 & $3.3\times 10^{7}$ & $3.1\times 10^{8}$ \\
\bottomrule
\end{tabular}
\end{table}


\clearpage
%% GENERATED FILE -- do not hand-edit.
%% SOURCE: experiments/database/scripts/build_bacteria_candidate_datasets_table.py
\begin{table}[H]\centering
\small
\caption[]{Candidates by class and by host, split at the two tranche lines.
\emph{Genotypes}, \emph{Instances} and \emph{Meas.} sum the per-row axes over all
tranches; a row with no count contributes nothing, so every total is a lower bound.}
\label{tab:bcounts}
\begin{tabular}{@{}l r r r r r r@{}}
\toprule
 & Tr. 1 & Tr. 2 & Reserve & Genotypes & Instances & Meas. \\
\midrule
Aggregation / support & 0 & 1 & 3 & 18,222 & $1.9\times 10^{4}$ & $7.0\times 10^{4}$ \\
CRISPR library screen & 0 & 6 & 4 & 331,032 & $1.3\times 10^{6}$ & $1.7\times 10^{6}$ \\
Combinatorial design & 1 & 1 & 5 & 757 & $1.2\times 10^{3}$ & $1.4\times 10^{3}$ \\
Fitness / chemical genomics & 2 & 2 & 0 & 15,977 & $1.6\times 10^{6}$ & $9.8\times 10^{6}$ \\
Genetic interaction & 0 & 4 & 0 & 930,545 & $3.5\times 10^{5}$ & $3.5\times 10^{5}$ \\
Metabolome / flux & 1 & 2 & 1 & 3,990 & $3.5\times 10^{4}$ & $2.6\times 10^{8}$ \\
Modality / backbone & 0 & 0 & 2 & 189,386 & $2.3\times 10^{4}$ & $2.3\times 10^{4}$ \\
Morphology / imaging & 0 & 0 & 1 & 548 & 585 & $2.3\times 10^{3}$ \\
Multi-omics campaign & 1 & 2 & 5 & 1,120 & $1.7\times 10^{3}$ & $1.5\times 10^{6}$ \\
Production campaign & 6 & 2 & 10 & 868 & $4.2\times 10^{3}$ & $2.2\times 10^{6}$ \\
Proteome & 0 & 2 & 1 & 7 & 134 & $3.1\times 10^{5}$ \\
Tolerance / robustness & 3 & 0 & 4 & 249 & 813 & 881 \\
Transcriptome & 2 & 1 & 0 & 2,023 & $9.1\times 10^{3}$ & $6.3\times 10^{6}$ \\
Translation / turnover & 0 & 1 & 1 & 2 & 15 & $4.8\times 10^{4}$ \\
Transposon fitness & 4 & 6 & 5 & 1,755,369 & $2.9\times 10^{7}$ & $2.9\times 10^{7}$ \\
\midrule
\multicolumn{7}{@{}l}{\emph{The same rows, by host}}\\
E. coli & 11 & 25 & 27 & 3,029,348 & $3.0\times 10^{7}$ & $3.0\times 10^{8}$ \\
P. putida & 9 & 5 & 15 & 220,747 & $2.4\times 10^{6}$ & $6.4\times 10^{6}$ \\
\midrule
Total & 20 & 30 & 42 & 3,250,095 & $3.3\times 10^{7}$ & $3.1\times 10^{8}$ \\
\bottomrule
\end{tabular}
\end{table}



\end{document}
